% Radioactivity — Physics Upper Sixth — Cours signé OnBuch+
% Official MINESEC syllabus. Compiler avec : tectonic PHY12-radioactivity.tex
\newcommand{\DOCMATIERE}{Physics}
\newcommand{\DOCNIVEAU}{Upper Sixth}
\newcommand{\DOCMODULE}{Upper Sixth — Physics}
\newcommand{\DOCLECON}{12}
\newcommand{\DOCTITRE}{Radioactivity}
% OnBuch+ — préambule « cours signés » ENRICHI (compilé avec tectonic / XeLaTeX)
% Système visuel « Pop » d'OnBuch+ : fond crème (jamais blanc), bordures encre
% épaisses, ombres « dures » décalées, titres Archivo Black en capitales.
% Version MATHÉMATIQUES : graphes (pgfplots/tikz), boîtes pédagogiques
% (définition, propriété, méthode, attention, le savais-tu, exercice résolu,
% exercice, TP, bilan), page de garde et signature OnBuch+ ancrée en bas.
%
% Macros attendues dans main.tex AVANT \input{preamble} :
%   \DOCMATIERE  \DOCNIVEAU  \DOCMODULE  \DOCLECON (numéro)  \DOCTITRE
\documentclass[11pt]{article}

\usepackage{fontspec}
\usepackage[english]{babel}
\usepackage[a4paper,margin=2cm,top=3.4cm,bottom=2.8cm,headheight=1.4cm,footskip=1.3cm]{geometry}
\usepackage{amsmath,amssymb,amsfonts}
\usepackage{mathtools} % \xmapsto, \coloneqq... souvent utilisés par les modèles
\usepackage{cancel} % simplifications barrées (bilans de forces)
% \abs{z} (module, valeur absolue) et \norm{u} (norme), utilisés par les modèles
\providecommand{\abs}{}\let\abs\relax\DeclarePairedDelimiter{\abs}{\lvert}{\rvert}
\providecommand{\norm}{}\let\norm\relax\DeclarePairedDelimiter{\norm}{\lVert}{\rVert}
\usepackage[table]{xcolor} % [table] : fournit aussi \rowcolors (alternance de lignes)
\usepackage{pagecolor}
\usepackage{tikz}
\usetikzlibrary{arrows.meta,positioning,calc,shapes.geometric,shapes.misc,shapes.arrows,decorations.pathmorphing,decorations.pathreplacing,decorations.markings,patterns,shadows,mindmap,trees,fit,backgrounds,matrix,intersections,angles,quotes}
\usepackage{pgfplots}
\usepackage[version=4]{mhchem} % équations nucléaires \ce{^{235}_{92}U}
\usepackage[european,siunitx]{circuitikz} % schémas électriques
\pgfplotsset{compat=1.18}
\usepgfplotslibrary{fillbetween,groupplots} % aires entre courbes (calcul intégral)
\usepackage{enumitem}
\usepackage{fancyhdr}
% [most] charge toutes les bibliothèques tcolorbox (listings, minted, raster,
% documentation...) et gonfle inutilement la mémoire TeX sur un document long
% (~300 boîtes) : on ne charge que ce qui est réellement utilisé.
\usepackage[skins,breakable]{tcolorbox}
\usepackage{graphicx}
\usepackage{colortbl}
\usepackage{booktabs}
\usepackage{tabularx}
\usepackage{diagbox} % case d'en-tête diagonale des tableaux à double entrée
% Colonnes extensibles centrée / gauche / droite (C, L, R), souvent utilisées par les modèles
\newcolumntype{C}{>{\centering\arraybackslash}X}
\newcolumntype{L}{>{\raggedright\arraybackslash}X}
\newcolumntype{R}{>{\raggedleft\arraybackslash}X}
\usepackage{array}
\usepackage{multicol}
\usepackage{siunitx}
% parse-numbers=false : accepte aussi \SI{2,5 \times 10^{-3}}{...}
\sisetup{locale=UK,output-decimal-marker={.},group-minimum-digits=4,parse-numbers=false}
\DeclareSIUnit\ohm{\ensuremath{\symup{\Omega}}} % U+2126 absent de la police
\DeclareSIUnit\division{div} % réglages d'oscilloscope (ms/div, V/div)
\DeclareSIUnit\tr{tr} % tours (vitesse de rotation, ex. \SI{1500}{\tr\per\minute})
\usepackage{qrcode}
% \degree : commande fréquemment utilisée par les modèles pour un angle ou une
% température en dehors de \SI{}{} (non fournie par les paquets chargés).
\providecommand{\degree}{^{\circ}}
% \qed (fin de démonstration), utilisé par les modèles sans amsthm
\providecommand{\qed}{\hfill$\square$}
% \barre : double barre des valeurs interdites dans un tableau de variations (tkz-tab)
\providecommand{\barre}{\ensuremath{\|}}
% environnement proof (amsthm non chargé) utilisé par les modèles
\ifdefined\proof\else\newenvironment{proof}[1][Proof]{\par\noindent\textit{#1.}\ }{\qed\par}\fi
\usepackage{lastpage}
\usepackage{hyperref}
\hypersetup{hidelinks}

% ============================================================
% Palette « Pop » OnBuch+ — mêmes hex que la classe P (plus_ui.dart)
% ============================================================
\definecolor{poporange}{HTML}{F59321}
\definecolor{popdark}{HTML}{E07A0C}
\definecolor{popcream}{HTML}{FFF6E8}
\definecolor{popink}{HTML}{1C1714}
\definecolor{poppurple}{HTML}{7A5AE0}
\definecolor{popgreen}{HTML}{2E9E6A}
\definecolor{poppink}{HTML}{E0457B}
\definecolor{popblue}{HTML}{2F7DE1}
\definecolor{popgold}{HTML}{C98A12}
\definecolor{popmuted}{HTML}{8A7F76}
\definecolor{popline}{HTML}{EADFD2}
\definecolor{popgray}{HTML}{8A7F76}
\colorlet{popgrey}{popgray}
% Alias tolérants (le modèle nomme parfois les couleurs différemment)
\colorlet{popwhite}{white}
\colorlet{popblack}{popink}
\colorlet{popred}{poppink}
\colorlet{popyellow}{popgold}
% Teintes claires pour fonds de tableaux / zones
\colorlet{popblueL}{popblue!10!white}
\colorlet{poporangeL}{poporange!14!white}
\colorlet{popgreenL}{popgreen!12!white}
\colorlet{poppurpleL}{poppurple!12!white}
\colorlet{poppinkL}{poppink!12!white}
\colorlet{popgoldL}{popgold!14!white}

\pagecolor{popcream}

% ============================================================
% Typographie
% ============================================================
\defaultfontfeatures{Ligatures=TeX}
\setmainfont{PlusJakartaSans-Medium.ttf}[Path=../../../fonts/,BoldFont=PlusJakartaSans-Bold.ttf]
% Maths en sans-serif (Fira Math) avec chiffres et lettres droites de Plus
% Jakarta, pour que les formules se fondent dans le texte.
\usepackage[math-style=ISO,bold-style=ISO]{unicode-math}
\setmathfont{FiraMath-Regular.otf}[Path=../../../fonts/]
\setmathfont{PlusJakartaSans-Medium.ttf}[Path=../../../fonts/,range={up/num,up/{Latin,latin},\mathup}]
% (icomma retiré : décimales avec point en anglais)
% Caractères mathématiques Unicode absents des polices de texte (Plus Jakarta,
% Archivo Black) : rendus par leur équivalent mathématique, en texte comme en
% formule — sinon ils s'affichent en carré vide (ex. « Arithmétique dans ℤ »).
\usepackage{newunicodechar}
\newunicodechar{ℝ}{\ensuremath{\mathbb{R}}}
\newunicodechar{ℤ}{\ensuremath{\mathbb{Z}}}
\newunicodechar{ℂ}{\ensuremath{\mathbb{C}}}
\newunicodechar{ℕ}{\ensuremath{\mathbb{N}}}
\newunicodechar{ℚ}{\ensuremath{\mathbb{Q}}}
\newunicodechar{𝔻}{\ensuremath{\mathbb{D}}}
\newunicodechar{α}{\ensuremath{\alpha}}
\newunicodechar{β}{\ensuremath{\beta}}
\newunicodechar{γ}{\ensuremath{\gamma}}
\newunicodechar{δ}{\ensuremath{\delta}}
\newunicodechar{ε}{\ensuremath{\varepsilon}}
\newunicodechar{θ}{\ensuremath{\theta}}
\newunicodechar{λ}{\ensuremath{\lambda}}
\newunicodechar{μ}{\ensuremath{\mu}}
\newunicodechar{σ}{\ensuremath{\sigma}}
\newunicodechar{τ}{\ensuremath{\tau}}
\newunicodechar{φ}{\ensuremath{\varphi}}
\newunicodechar{ψ}{\ensuremath{\psi}}
\newunicodechar{ω}{\ensuremath{\omega}}
\newunicodechar{Σ}{\ensuremath{\Sigma}}
% majuscules produites par \MakeUppercase dans les titres (α-aminés -> Α)
\newunicodechar{Α}{\ensuremath{\alpha}}
\newunicodechar{Β}{\ensuremath{\beta}}
\newunicodechar{Γ}{\ensuremath{\Gamma}}
\newunicodechar{Δ}{\ensuremath{\Delta}}
\newunicodechar{Ε}{\ensuremath{\varepsilon}}
\newunicodechar{Θ}{\ensuremath{\Theta}}
\newunicodechar{Λ}{\ensuremath{\Lambda}}
\newunicodechar{Μ}{\ensuremath{\mu}}
\newunicodechar{Τ}{\ensuremath{\tau}}
\newunicodechar{Φ}{\ensuremath{\Phi}}
\newunicodechar{Ψ}{\ensuremath{\Psi}}
\newunicodechar{Ω}{\ensuremath{\Omega}}
\newunicodechar{↦}{\ensuremath{\mapsto}}
\newunicodechar{∈}{\ensuremath{\in}}
\newunicodechar{∉}{\ensuremath{\notin}}
\newunicodechar{∧}{\ensuremath{\wedge}}
\newunicodechar{⇔}{\ensuremath{\Leftrightarrow}}
\newunicodechar{⟺}{\ensuremath{\Longleftrightarrow}}
\newunicodechar{⇒}{\ensuremath{\Rightarrow}}
\newunicodechar{⟹}{\ensuremath{\Longrightarrow}}
\newunicodechar{∘}{\ensuremath{\circ}}
\newunicodechar{∪}{\ensuremath{\cup}}
\newunicodechar{∩}{\ensuremath{\cap}}
\newunicodechar{≡}{\ensuremath{\equiv}}
\newunicodechar{⊂}{\ensuremath{\subset}}
\newunicodechar{⊥}{\ensuremath{\perp}}
\newunicodechar{∃}{\ensuremath{\exists}}
\newunicodechar{∀}{\ensuremath{\forall}}
\newunicodechar{∛}{\ensuremath{\sqrt[3]{\,}}}
\newunicodechar{∜}{\ensuremath{\sqrt[4]{\,}}}
\newunicodechar{⃗}{\ensuremath{{}^{\to}}}
\newunicodechar{ⁿ}{\ifmmode^{n}\else\textsuperscript{\ensuremath{n}}\fi}
\newunicodechar{ˣ}{\ifmmode^{x}\else\textsuperscript{\ensuremath{x}}\fi}
\newunicodechar{ᵇ}{\ifmmode^{b}\else\textsuperscript{\ensuremath{b}}\fi}
\newunicodechar{ʳ}{\ifmmode^{r}\else\textsuperscript{\ensuremath{r}}\fi}
\newunicodechar{ᵘ}{\ifmmode^{u}\else\textsuperscript{\ensuremath{u}}\fi}
\newunicodechar{ʸ}{\ifmmode^{y}\else\textsuperscript{\ensuremath{y}}\fi}
\newunicodechar{ᵃ}{\ifmmode^{a}\else\textsuperscript{\ensuremath{a}}\fi}
\newunicodechar{ᵉ}{\ifmmode^{e}\else\textsuperscript{\ensuremath{e}}\fi}
\newunicodechar{ᵏ}{\ifmmode^{k}\else\textsuperscript{\ensuremath{k}}\fi}
\newunicodechar{ᵖ}{\ifmmode^{p}\else\textsuperscript{\ensuremath{p}}\fi}
\newunicodechar{ᴺ}{\ifmmode^{N}\else\textsuperscript{\ensuremath{N}}\fi}
\newunicodechar{ᶜ}{\ifmmode^{c}\else\textsuperscript{\ensuremath{c}}\fi}
\newunicodechar{ʰ}{\ifmmode^{h}\else\textsuperscript{\ensuremath{h}}\fi}
\newunicodechar{ᵠ}{\ifmmode^{q}\else\textsuperscript{\ensuremath{q}}\fi}
\newunicodechar{ₙ}{\ifmmode_{n}\else\textsubscript{\ensuremath{n}}\fi}
\newunicodechar{ₐ}{\ifmmode_{a}\else\textsubscript{\ensuremath{a}}\fi}
\newunicodechar{ᵢ}{\ifmmode_{i}\else\textsubscript{\ensuremath{i}}\fi}
\newunicodechar{ᵣ}{\ifmmode_{r}\else\textsubscript{\ensuremath{r}}\fi}
\newunicodechar{ₖ}{\ifmmode_{k}\else\textsubscript{\ensuremath{k}}\fi}
\newunicodechar{ᵦ}{\ifmmode_{\beta}\else\textsubscript{\ensuremath{\beta}}\fi}
\newunicodechar{ₓ}{\ifmmode_{x}\else\textsubscript{\ensuremath{x}}\fi}
\newfontfamily\popdisplay{ArchivoBlack-Regular.ttf}[Path=../../../fonts/]
\newfontfamily\poplogo{SpaceGrotesk-Bold.ttf}[Path=../../../fonts/]
\newfontfamily\popsemi{PlusJakartaSans-SemiBold.ttf}[Path=../../../fonts/]
\newfontfamily\popxbold{PlusJakartaSans-ExtraBold.ttf}[Path=../../../fonts/]
\newfontfamily\popxboldls{PlusJakartaSans-ExtraBold.ttf}[Path=../../../fonts/,LetterSpace=3.0]

\setlist[enumerate]{leftmargin=1.7em,itemsep=5pt,topsep=4pt}
\setlist[itemize]{leftmargin=1.7em,itemsep=4pt,topsep=4pt,label={\color{poporange}\textbullet}}
\setlength{\parindent}{0pt}
\setlength{\parskip}{5pt}
\renewcommand{\arraystretch}{1.25}

% pgfplots : style Pop par défaut
\pgfplotsset{
  every axis/.append style={axis line style={popink,line width=1pt},tick style={popink},
    grid style={popline,line width=0.5pt},label style={font=\small},tick label style={font=\footnotesize}},
  popcurve/.style={poporange,line width=1.8pt,no markers,smooth},
}
% Même style disponible dans un \draw TikZ simple (hors pgfplots)
\tikzset{popcurve/.style={poporange,line width=1.8pt}}
\tikzset{popgrid/.style={popline,very thin}}
\colorlet{popcurve}{poporange} % le nom du style est parfois utilisé comme couleur

% ============================================================
% Pill / squiggle
% ============================================================
\newcommand{\poppill}[3]{%
  \tikz[baseline=-0.55ex]\node[draw=popink,line width=1.2pt,rounded corners=2.6pt,%
    fill=#2,inner xsep=6.5pt,inner ysep=3pt]{%
    \popxboldls\fontsize{7}{7}\selectfont\color{#3}\MakeUppercase{#1}};%
}
\newcommand{\popsquiggle}[1]{%
  \tikz[baseline=-0.4ex]{
    \draw[#1,line width=2.6pt,line cap=round]
      plot [smooth] coordinates {(0,0.09) (0.34,0.01) (0.68,0.21) (1.02,0.01) (1.36,0.10)};
  }%
}
% Mot-clé surligné (marqueur orange sous le texte)
\newcommand{\cle}[1]{{\popxbold\color{popdark}#1}}

% ============================================================
% En-tête / pied
% ============================================================
\pagestyle{fancy}
\fancyhf{}
\renewcommand{\headrulewidth}{0pt}
\renewcommand{\footrulewidth}{0pt}
\lhead{\raisebox{-0.15ex}{\poplogo\fontsize{13}{13}\selectfont\color{popink}OnBuch\color{poporange}+}}
\rhead{\poppill{\DOCMATIERE\ \textbullet\ \DOCNIVEAU\ \textbullet\ Chapter \DOCLECON}{white}{popink}}
\lfoot{\popxbold\fontsize{7.6}{7.6}\selectfont\color{popmuted}\MakeUppercase{Signed course OnBuch+}\ \textbullet\ {\popsemi\DOCTITRE}}
\rfoot{\tikz[baseline=-0.6ex]\node[rounded corners=8.5pt,fill=popink,minimum height=17pt,minimum width=17pt,inner xsep=5pt,inner sep=0]{\popxbold\fontsize{8}{8}\selectfont\color{popcream}\thepage\,/\,\pageref*{LastPage}};}

% ============================================================
% Titres
% ============================================================
\newcommand{\chaptitre}[2]{%
  \begin{tcolorbox}[enhanced,colback=poporange,colframe=popink,boxrule=2.6pt,arc=11pt,
      drop shadow=popink,left=15pt,right=15pt,top=12pt,bottom=11pt,before skip=3pt,after skip=18pt]
    \poppill{#2}{popink}{popcream}\par\vspace{9pt}
    {\raggedright\hyphenpenalty=10000\exhyphenpenalty=10000\popdisplay\fontsize{22}{24}\selectfont\color{popcream}\MakeUppercase{#1}\par}
  \end{tcolorbox}
}
% Section numérotée : pastille orange avec le numéro + titre Archivo
\newcounter{popsec}
\newcommand{\coursec}[1]{%
  \stepcounter{popsec}\setcounter{popsub}{0}%
  \par\vspace{16pt}\noindent
  \begin{minipage}{\linewidth}
  \tikz[baseline=-0.7ex]\node[draw=popink,line width=1.6pt,fill=poporange,rounded corners=4pt,minimum size=22pt,inner sep=0]{\popdisplay\fontsize{12}{12}\selectfont\color{popcream}\thepopsec};%
  \hspace{8pt}{\popdisplay\fontsize{15}{17}\selectfont\color{popink}\MakeUppercase{#1}}\par
  \vspace{4pt}\noindent{\color{poporange}\rule{36pt}{3.6pt}}
  \end{minipage}\par\nopagebreak\vspace{8pt}
}
\newcounter{popsub}
\newcommand{\courssub}[1]{%
  \stepcounter{popsub}%
  \par\vspace{9pt}\noindent{\popxbold\fontsize{11.5}{13}\selectfont\color{popink}{\color{poporange}\thepopsec.\thepopsub}\hspace{6pt}#1}\par\nopagebreak\vspace{4pt}
}

% ============================================================
% Boîtes pédagogiques — même recette : fond blanc, bordure épaisse colorée,
% ombre dure encre, badge plein. Titre optionnel [#1].
% ============================================================
\tcbset{popbase/.style={breakable,enhanced,colback=white,boxrule=2.3pt,arc=9pt,
  shadow={3pt}{-3pt}{0pt}{popink},left=14pt,right=14pt,top=11pt,bottom=11pt,before skip=11pt,after skip=15pt,
  fontupper=\popsemi\fontsize{10.6}{14.5}\selectfont\color{popink},
  fontlower=\popsemi\fontsize{10.6}{14.5}\selectfont\color{popink}}}
% Coche dessinée (le glyphe ✓ n'existe pas dans les polices du document)
\newcommand{\popcheck}[1]{\tikz[baseline=-0.5ex]\draw[#1,line width=1.6pt,line cap=round,line join=round](-0.13,0.0)--(-0.04,-0.09)--(0.13,0.1);}
\providecommand{\coche}{\popcheck{popgreen}}
\providecommand{\resultat}[1]{\par\smallskip\noindent\fcolorbox{popgreen}{popgreenL}{\parbox{\dimexpr\linewidth-2\fboxsep-2\fboxrule\relax}{\textbf{Result.} #1}}\par\smallskip}
\newcommand{\popboxhead}[3]{\poppill{#1}{#2}{white}\ifx\relax#3\relax\else\hspace{6pt}{\popxbold\fontsize{10.6}{12}\selectfont\color{popink}#3}\fi\par\vspace{7pt}}

\newtcolorbox{aretenir}[1][]{popbase,colframe=popgreen,before upper={\popboxhead{Key points}{popgreen}{#1}}}
\newtcolorbox{exemplebox}[1][]{popbase,colframe=poporange,before upper={\popboxhead{Exemple}{poporange}{#1}}}
\newtcolorbox{definition}[1][]{popbase,colframe=popblue,before upper={\popboxhead{Definition}{popblue}{#1}}}
\newtcolorbox{propriete}[1][]{popbase,colframe=poppurple,before upper={\popboxhead{Property}{poppurple}{#1}}}
\newtcolorbox{methode}[1][]{popbase,colframe=popgold,before upper={\popboxhead{Method}{popgold}{#1}}}
\newtcolorbox{attention}[1][]{popbase,colframe=poppink,before upper={\popboxhead{Warning}{poppink}{#1}}}
\newtcolorbox{experience}[1][]{popbase,colframe=popblue,colback=popblueL,before upper={\popboxhead{Experiment}{popblue}{#1}}}
% « Le savais-tu ? » : carte violette pleine (contexte, culture, Cameroun)
\newtcolorbox{savaistu}[1][]{popbase,colback=poppurple,colframe=popink,
  fontupper=\popsemi\fontsize{10.4}{14.2}\selectfont\color{white},
  before upper={\poppill{Did you know?}{white}{poppurple}\ifx\relax#1\relax\else\hspace{6pt}{\popxbold\color{white}#1}\fi\par\vspace{7pt}}}
% Exercice résolu : énoncé en haut, \tcblower puis la solution sur fond crème
\newcounter{popexo}\newcounter{popexr}
\newtcolorbox{exoresolu}[1][]{popbase,colframe=poporange,colbacklower=popcream,
  segmentation style={popink,line width=1.2pt,dashed},
  before upper={\stepcounter{popexr}\popboxhead{Worked example \thepopexr}{poporange}{#1}},
  before lower={\poppill{Solution}{popink}{popcream}\par\vspace{6pt}}}
% Exercice (non résolu en place) — la correction est dans la partie Corrigés
\newtcolorbox{exercice}[1][]{popbase,colframe=popink,
  before upper={\stepcounter{popexo}\popboxhead{Exercise \thepopexo}{popink}{#1}}}
% Corrigé
\newtcolorbox{corrige}[1][]{popbase,colframe=popgreen,colback=popgreenL,
  before upper={\popboxhead{Solution}{popgreen}{#1}}}
% Objectifs / prérequis / bilan
\newtcolorbox{objectifs}{popbase,colframe=popink,colback=poporangeL,
  before upper={\popboxhead{Chapter objectives}{popink}{}}}
\newtcolorbox{prerequis}{popbase,colframe=popink,colback=popblueL,
  before upper={\popboxhead{Prerequisites}{popblue}{}}}
\newtcolorbox{bilan}{popbase,colframe=popink,colback=white,boxrule=2.8pt,
  before upper={\popboxhead{Fiche bilan}{poporange}{L'essentiel à connaître pour l'examen}}}
% Figure encadrée avec légende
\newtcolorbox{popfigure}[1][]{enhanced,breakable=false,colback=white,colframe=popline,boxrule=1.4pt,arc=8pt,
  left=10pt,right=10pt,top=10pt,bottom=8pt,before skip=10pt,after skip=12pt,halign=center,
  lower separated=false,halign lower=center,
  fontlower=\popsemi\fontsize{9}{11.5}\selectfont\color{popmuted},
  #1}
\newcommand{\legende}[1]{\par\vspace{6pt}{\popsemi\fontsize{9}{11.5}\selectfont\color{popmuted}#1}}

% ============================================================
% Signature OnBuch+
% ============================================================
\newcommand{\poplogoinline}[1]{{\poplogo\fontsize{#1}{#1}\selectfont\color{popink}OnBuch\color{poporange}+}}
\newcommand{\signatureonbuch}{%
  \begin{tcolorbox}[enhanced,colback=popink,colframe=popink,boxrule=2.2pt,arc=10pt,
      drop shadow=poporange,left=14pt,right=14pt,top=10pt,bottom=10pt,before skip=0pt,after skip=0pt]
    \tikz[baseline=-0.7ex]\node[circle,fill=popgreen,draw=popcream,line width=1.3pt,minimum size=22pt,inner sep=0]{\popcheck{white}};%
    \hspace{9pt}\begin{minipage}[c]{0.62\linewidth}
      {\popxbold\fontsize{12}{13}\selectfont\color{popcream}Signed\ \textbullet\ {\poplogo OnBuch\color{poporange}+}}\par\vspace{2pt}
      {\raggedright\popsemi\fontsize{8}{10.5}\selectfont\color{popline}\MakeUppercase{OnBuch+ teaching team}\\\MakeUppercase{Follows the official MINESEC syllabus}\par}
    \end{minipage}\hfill
    {\poplogo\fontsize{22}{22}\selectfont\color{popcream}OnBuch\color{poporange}+}
  \end{tcolorbox}%
}

% ============================================================
% Page de garde
% #1 titre · #2 matière · #3 niveau · #4 module · #5 n° leçon · #6 durée ·
% #7 description / objectifs (texte ou itemize)
% ============================================================
\newcommand{\pagegarde}[7]{%
  \thispagestyle{empty}
  \newgeometry{a4paper,margin=1.7cm,top=1.6cm,bottom=1.4cm}%
  \begin{tikzpicture}[remember picture,overlay]
    \fill[poporange!18] ([xshift=-3.2cm,yshift=-3.6cm]current page.north east) circle (4.6cm);
    \fill[poppurple!14] ([xshift=2.4cm,yshift=7.6cm]current page.south west) circle (3.8cm);
  \end{tikzpicture}%
  {\poplogo\fontsize{30}{30}\selectfont\color{popink}OnBuch\color{poporange}+}\hfill
  \raisebox{0.6ex}{\poppill{Signed course \textbullet\ Premium}{popink}{popcream}}\par
  \vspace{1pt}\popsquiggle{poppurple}\par
  \vspace{8pt}
  \begin{tcolorbox}[enhanced,colback=poporange,colframe=popink,boxrule=2.8pt,arc=13pt,
      drop shadow=popink,left=18pt,right=18pt,top=12pt,bottom=13pt,after skip=10pt]
    \poppill{#2 \textbullet\ #3}{popink}{popcream}\hspace{5pt}\poppill{#4}{white}{popink}\par\vspace{8pt}
    {\popdisplay\fontsize{14}{14}\selectfont\color{popink}CHAPTER #5}\par\vspace{4pt}
    {\raggedright\hyphenpenalty=10000\exhyphenpenalty=10000\popdisplay\fontsize{25}{27}\selectfont\color{popcream}\MakeUppercase{#1}\par}
    \vspace{8pt}
    {\popxbold\fontsize{9.5}{9.5}\selectfont\color{popink}\MakeUppercase{Suggested time: #6}}
  \end{tcolorbox}
  \begin{tcolorbox}[enhanced,colback=white,colframe=popink,boxrule=2.2pt,arc=10pt,
      drop shadow=popink,left=16pt,right=16pt,top=10pt,bottom=10pt,after skip=8pt]
    \poppill{In this chapter}{poppurple}{white}\par\vspace{5pt}
    \popsemi\fontsize{9.3}{12.6}\selectfont\color{popink}#7
  \end{tcolorbox}
  \noindent\begin{minipage}[t]{0.487\linewidth}
    \begin{tcolorbox}[enhanced,colback=popgreen,colframe=popink,boxrule=2.2pt,arc=10pt,
        drop shadow=popink,left=11pt,right=11pt,top=10pt,bottom=10pt,height=3.1cm,valign=center]
      \tcbox[colback=white,colframe=popink,boxrule=1.3pt,arc=5pt,boxsep=0pt,nobeforeafter,
        left=3pt,right=3pt,top=3pt,bottom=3pt,valign=center]{\qrcode[height=1.9cm]{https://play.google.com/store/apps/details?id=com.luvvix.onbuch&hl=en}}%
      \hspace{8pt}\begin{minipage}[c]{0.5\linewidth}
        {\popdisplay\fontsize{11}{12.5}\selectfont\color{white}\MakeUppercase{Download OnBuch}}\par\vspace{4pt}
        {\raggedright\popsemi\fontsize{8.4}{11}\selectfont\color{white}Free \textbullet\ Google Play\\AI tutor, past papers, results\par}
      \end{minipage}
    \end{tcolorbox}
  \end{minipage}\hfill
  \begin{minipage}[t]{0.487\linewidth}
    \begin{tcolorbox}[enhanced,colback=poppurple,colframe=popink,boxrule=2.2pt,arc=10pt,
        drop shadow=popink,left=11pt,right=11pt,top=10pt,bottom=10pt,height=3.1cm,valign=center]
      \tikz[baseline=-0.7ex]\node[circle,fill=white,draw=popink,line width=1.3pt,minimum size=1.6cm,inner sep=0]{\popdisplay\fontsize{24}{24}\selectfont\color{poppurple}+};%
      \hspace{8pt}\begin{minipage}[c]{0.6\linewidth}
        {\popdisplay\fontsize{11}{12.5}\selectfont\color{white}\MakeUppercase{Go OnBuch+}}\par\vspace{4pt}
        {\raggedright\popsemi\fontsize{8.4}{11}\selectfont\color{white}All signed courses, unlimited AI tutor, PDF sheets — OnBuch+ tab in the app\par}
      \end{minipage}
    \end{tcolorbox}
  \end{minipage}
  \vspace{8pt}
  \popcontenu
  \vfill
  \signatureonbuch
  \restoregeometry
  \newpage
}

% Rangée « Ce cours contient » (page de garde)
\newcommand{\poptile}[3]{% #1 couleur #2 gros texte #3 légende
  \begin{tcolorbox}[enhanced,colback=#1,colframe=popink,boxrule=2pt,arc=9pt,shadow={2.5pt}{-2.5pt}{0pt}{popink},
      left=6pt,right=6pt,top=9pt,bottom=9pt,halign=center,valign=center,height=1.75cm,width=\linewidth,nobeforeafter]
    {\popdisplay\fontsize{12.5}{13}\selectfont\color{white}\MakeUppercase{#2}}\par\vspace{3pt}
    {\centering\popsemi\fontsize{7.8}{9.5}\selectfont\color{white}#3\par}
  \end{tcolorbox}}
\newcommand{\popcontenu}{%
  \par\noindent{\popxboldls\fontsize{7.5}{7.5}\selectfont\color{popmuted}THIS SIGNED COURSE CONTAINS}\par\vspace{7pt}
  \noindent\begin{minipage}[t]{0.235\linewidth}\poptile{popblue}{Course}{complete, illustrated, step by step}\end{minipage}\hfill
  \begin{minipage}[t]{0.235\linewidth}\poptile{poporange}{Methods}{and worked examples}\end{minipage}\hfill
  \begin{minipage}[t]{0.235\linewidth}\poptile{poppink}{Exercises}{exam style + solutions}\end{minipage}\hfill
  \begin{minipage}[t]{0.235\linewidth}\poptile{popgreen}{Summary sheet}{the essentials to remember}\end{minipage}\par}

% Sommaire
\newtcolorbox{sommaire}{enhanced,colback=white,colframe=popink,boxrule=2.2pt,arc=10pt,
  drop shadow=popink,left=18pt,right=18pt,top=13pt,bottom=13pt,before skip=6pt,after skip=20pt,
  before upper={\poppill{Contents}{poppurple}{white}\par\vspace{9pt}\popsemi\fontsize{10.6}{14.5}\selectfont\color{popink}}}

% Signature de fin de cours — poussée tout en bas de la dernière page
\newcommand{\findecours}{%
  \par\vfill\nopagebreak
  \begin{center}\popsquiggle{poporange}\end{center}\vspace{4pt}
  \signatureonbuch
}
\AtBeginDocument{\def\llbracket{\mathopen{[\mkern-3.5mu[}}\def\rrbracket{\mathclose{]\mkern-3.5mu]}}} % intervalles d'entiers (glyphe ⟦ absent de la police)
% Glyphes absents de Fira Math : redéfinis à partir de glyphes disponibles
\AtBeginDocument{%
  \def\setminus{\mathbin{\backslash}}\let\smallsetminus\setminus
  \def\vdots{\mathord{\vbox{\baselineskip3.2pt\lineskiplimit0pt\kern4pt\hbox{.}\hbox{.}\hbox{.}}}}%
  \def\ddots{\mathinner{\mkern1mu\raise6.5pt\hbox{.}\mkern2mu\raise3.6pt\hbox{.}\mkern2mu\raise0.7pt\hbox{.}\mkern1mu}}%
  \def\triangle{\mathord{\scalebox{0.9}{$\symup{\Delta}$}}}%
  \def\nabla{\mathord{\raisebox{\depth}{\scalebox{1}[-1]{$\symup{\Delta}$}}}}%
  \def\checkmark{\popcheck{popdark}}%
  \def\star{\mathbin{\ast}}\def\bigstar{\mathord{\ast}}%
  \def\diamond{\mathbin{\scalebox{0.6}{$\Diamond$}}}%
  \def\bigcup{\mathop{\mathchoice{\raisebox{-0.3em}{\scalebox{1.7}{$\cup$}}}{\scalebox{1.25}{$\cup$}}{\cup}{\cup}}\displaylimits}%
  \def\bigcap{\mathop{\mathchoice{\raisebox{-0.3em}{\scalebox{1.7}{$\cap$}}}{\scalebox{1.25}{$\cap$}}{\cap}{\cap}}\displaylimits}%
  \def\lesseqqgtr{\mathrel{\lessgtr}}\def\bigoplus{\mathop{\oplus}\displaylimits}%
}
\providecommand{\ssi}{if and only if} % abréviation fréquente
\AtBeginDocument{\providecommand{\wideparen}{\overparen}} % arc de cercle

%% FIGFIX-BEGIN v3 — correctifs globaux des figures (docs/onbuch-generation/tools/figfix)
\usepackage{adjustbox}\usepackage{etoolbox}
% toute figure TikZ/pgfplots plus large que la ligne est réduite à la largeur disponible
\BeforeBeginEnvironment{tikzpicture}{\begin{adjustbox}{max width=\linewidth}}
\AfterEndEnvironment{tikzpicture}{\end{adjustbox}}
% légendes pgfplots lisibles par-dessus les courbes
\pgfplotsset{every axis/.append style={legend style={fill=white,fill opacity=0.92,text opacity=1,draw=black!15,inner sep=3pt}}}
% étiquettes d'arêtes (syntaxe edge["..."]) avec fond blanc
\tikzset{every edge quotes/.append style={fill=white,inner sep=1.2pt}}
% bulles de cartes mentales : texte à la ligne dans la bulle, bulle un peu plus grande
\tikzset{every concept/.append style={text width=2.0cm,align=center,minimum size=2.3cm,inner sep=2pt}}
%% FIGFIX-END
\begin{document}

\pagegarde{Radioactivity}{Physics}{Upper Sixth}{Upper Sixth — Physics}{12}{2 periods}{This chapter studies the stability of atomic nuclei, the nature and properties of radioactive radiations (alpha, beta, gamma), and the energy released in nuclear processes through the mass-energy relation. It ends with the two great nuclear transformations that power reactors and stars: fission and fusion.
\vspace{4pt}
\begin{multicols}{2}\small
\begin{itemize}
\item By the end of this chapter I can write the symbolic representation of a nuclide and identify isotopes
\item By the end of this chapter I can explain nuclear stability using the N/Z ratio and the valley of stability
\item By the end of this chapter I can write and balance nuclear equations using conservation of nucleon number A and charge number Z
\item By the end of this chapter I can state the nature, charge, mass, penetrating power and ionising power of alpha, beta-minus, beta-plus and gamma radiations
\item By the end of this chapter I can describe practical applications of radioactivity in medicine, agriculture, industry and dating
\item By the end of this chapter I can define and calculate the mass defect and binding energy of a nucleus
\end{itemize}\end{multicols}}

\begin{sommaire}
\begin{multicols}{2}
\begin{enumerate}
\item Atomic Nucleus and Nuclear Stability
\item Types of Radiation, Properties and Applications
\item Mass Defect and Nuclear Energy
\item Nuclear Fission and Nuclear Fusion
\item Integration activity
\item Methods and worked examples
\item Exercises
\item Solutions to the exercises
\item Summary sheet
\end{enumerate}
\end{multicols}
\end{sommaire}

\begin{objectifs}
\begin{itemize}
\item By the end of this chapter I can write the symbolic representation of a nuclide and identify isotopes
\item By the end of this chapter I can explain nuclear stability using the N/Z ratio and the valley of stability
\item By the end of this chapter I can write and balance nuclear equations using conservation of nucleon number A and charge number Z
\item By the end of this chapter I can state the nature, charge, mass, penetrating power and ionising power of alpha, beta-minus, beta-plus and gamma radiations
\item By the end of this chapter I can describe practical applications of radioactivity in medicine, agriculture, industry and dating
\item By the end of this chapter I can define and calculate the mass defect and binding energy of a nucleus
\item By the end of this chapter I can apply Einstein's relation E = mc² to compute energy released in nuclear reactions (in joules and in MeV)
\item By the end of this chapter I can use the binding energy per nucleon curve to predict whether a reaction releases or absorbs energy
\item By the end of this chapter I can distinguish nuclear fission from nuclear fusion, write typical equations, and give their conditions and applications
\end{itemize}
\end{objectifs}

\begin{prerequis}
\begin{itemize}
\item Structure of the atom: protons, neutrons, electrons; atomic number Z and mass number A (Lower Sixth)
\item Conservation of electric charge and of matter in chemical/physical equations
\item Energy units: joule, electron-volt; conversions and powers of ten
\item Mass-energy equivalence idea E = mc² introduced qualitatively in Lower Sixth
\item Exponential behaviour and half-life notions (radioactive decay law from Lower Sixth, useful for dating applications)
\end{itemize}
\end{prerequis}

\newpage

\coursec{Atomic Nucleus and Nuclear Stability}

In 2019, a cancer patient at the Yaound\'e General Hospital received radiotherapy: a tiny radioactive source, no bigger than a grain of rice, destroyed tumour cells without a single surgical cut. Yet the same word ``radiation'' makes market traders in Douala panic when they hear about uranium mining projects in the North Region. How can the same physical phenomenon both heal and harm? And how can a few grams of matter release, through fission or fusion, the energy of tonnes of petrol? The answer lies hidden in a region of space a hundred thousand times smaller than the atom itself: the \cle{nucleus}. In this first section we open the nucleus, identify its inhabitants, and discover why some nuclei sit quietly forever while others spontaneously explode, emitting the radiations that can kill cancer cells --- or healthy ones.

\courssub{Composition of the Nucleus}

\textbf{Observation.} Every atom around you --- the carbon in your exercise book, the oxygen you breathe, the iron in a Douala--Yaound\'e railway line --- is built on the same plan: a tiny, dense, positively charged central \cle{nucleus}, surrounded by a cloud of electrons. The nucleus occupies almost no space (its radius is of order $10^{-15}\ \mathrm{m}$ while the atom's is of order $10^{-10}\ \mathrm{m}$), yet it contains more than $99.97\ \%$ of the atom's mass. Whatever makes atoms heavy, and whatever makes some of them unstable, lives in the nucleus.

The nucleus is made of two kinds of particles, collectively called \cle{nucleons}:

\begin{itemize}
\item the \cle{proton}, of charge $+e = +1.60\times 10^{-19}\ \mathrm{C}$ and mass $m_p = 1.673\times 10^{-27}\ \mathrm{kg}$;
\item the \cle{neutron}, electrically neutral, of mass $m_n = 1.675\times 10^{-27}\ \mathrm{kg}$, very slightly heavier than the proton.
\end{itemize}

\begin{definition}[Nucleon, nuclide, isotope]
A \cle{nucleon} is any particle found in the nucleus: a proton or a neutron.\par
A \cle{nuclide} is a nuclear species characterised by its number $Z$ of protons (\cle{atomic number}) and its total number $A$ of nucleons (\cle{mass number}). It is written symbolically as
\[ {}^{A}_{Z}X \qquad \text{with} \qquad N = A - Z \ \text{neutrons,} \]
where $X$ is the chemical symbol of the element.\par
\cle{Isotopes} are nuclides of the \emph{same element} (same $Z$) but with different numbers of neutrons (different $A$).
\end{definition}

\begin{popfigure}
\begin{center}
\begin{tikzpicture}[scale=1.0]
% Big nuclide symbol
\node[font=\Huge\bfseries, text=popink] at (0,0) {${}^{A}_{Z}X$};
% A annotation
\draw[->, thick, poporange] (-1.15,0.45) -- (-1.15,1.3);
\node[above, text=poporange, align=center, text width=3.4cm] at (-1.15,1.3) {$A$: mass number\\ (protons $+$ neutrons)};
% Z annotation
\draw[->, thick, popblue] (-1.15,-0.45) -- (-1.15,-1.3);
\node[below, text=popblue, align=center, text width=3.4cm] at (-1.15,-1.3) {$Z$: atomic number\\ (number of protons)};
% X annotation
\draw[->, thick, popgreen] (0.15,-0.45) -- (1.6,-1.3);
\node[below, text=popgreen, align=center, text width=3.6cm] at (1.6,-1.3) {$X$: chemical symbol\\ (fixed by $Z$)};
% N annotation
\draw[->, thick, poppurple] (0.15,0.45) -- (1.6,1.3);
\node[above, text=poppurple, align=center, text width=3.6cm] at (1.6,1.3) {$N = A - Z$:\\ number of neutrons};
\end{tikzpicture}
\legende{Anatomy of the nuclide symbol ${}^{A}_{Z}X$: $Z$ fixes the element, $A$ counts all nucleons, and $N = A - Z$ counts the neutrons.}
\end{center}
\end{popfigure}

\begin{exemplebox}[Reading nuclide symbols]
\begin{itemize}
\item $\ensuremath{\mathrm{^{12}_{6}C}}$: $Z = 6$ protons, $A = 12$ nucleons, so $N = 12 - 6 = 6$ neutrons. This is the common carbon of living matter.
\item $\ensuremath{\mathrm{^{14}_{6}C}}$: same element ($Z = 6$), but $N = 8$ neutrons. It is an isotope of carbon --- the famous ``carbon-14'' used for dating archaeological remains.
\item $\ensuremath{\mathrm{^{235}_{92}U}}$ and $\ensuremath{\mathrm{^{238}_{92}U}}$: two isotopes of uranium, with $143$ and $146$ neutrons respectively. Only the first undergoes fission easily in a reactor.
\end{itemize}
\end{exemplebox}

Two neighbouring notions complete the vocabulary:

\begin{definition}[Isobars and isotones]
\cle{Isobars} are nuclides with the \emph{same mass number} $A$ but different $Z$ (e.g.\ $\ensuremath{\mathrm{^{14}_{6}C}}$ and $\ensuremath{\mathrm{^{14}_{7}N}}$).\par
\cle{Isotones} are nuclides with the \emph{same number of neutrons} $N$ but different $Z$ (e.g.\ $\ensuremath{\mathrm{^{13}_{6}C}}$ and $\ensuremath{\mathrm{^{14}_{7}N}}$, both with $N = 7$).
\end{definition}

\begin{attention}[Do not confuse the numbers!]
\begin{itemize}
\item $A$ is a \emph{count of nucleons} (a whole number, no unit). It is \textbf{not} the atomic mass in $\mathrm{u}$: the mass of $\ensuremath{\mathrm{^{12}_{6}C}}$ is $12.000\ \mathrm{u}$, but that of $\ensuremath{\mathrm{^{14}_{6}C}}$ is $14.003\ \mathrm{u}$, not exactly $14$.
\item $Z$ is the number of \emph{protons}, \textbf{not} the number of neutrons. The neutron number is $N = A - Z$.
\item The chemical element is fixed by $Z$ alone: changing $N$ gives an isotope of the same element; changing $Z$ changes the element itself.
\end{itemize}
\end{attention}

\courssub{Nuclear Size, Density and the Strong Nuclear Force}

\textbf{How big is a nucleus?} Scattering experiments (firing alpha particles or electrons at nuclei, in the spirit of Rutherford) show that the nuclear radius grows with the mass number according to
\[ R = r_0\, A^{1/3}, \qquad r_0 \approx 1.2\times 10^{-15}\ \mathrm{m} = 1.2\ \mathrm{fm}. \]
The unit $1\ \mathrm{fm} = 10^{-15}\ \mathrm{m}$ is called the \emph{femtometre} (or ``fermi''). For $\ensuremath{\mathrm{^{238}_{92}U}}$, $R \approx 1.2\times 10^{-15}\times 238^{1/3} \approx 7.4\times 10^{-15}\ \mathrm{m}$.

Since the volume $V = \tfrac{4}{3}\pi R^3 = \tfrac{4}{3}\pi r_0^3 A$ is proportional to $A$, and the mass is also proportional to $A$, the \cle{density of nuclear matter is the same for all nuclei}:
\[ \rho = \frac{A\, m_n}{\tfrac{4}{3}\pi r_0^3 A} = \frac{3 m_n}{4\pi r_0^3} \approx \frac{3\times 1.675\times 10^{-27}}{4\pi (1.2\times 10^{-15})^3} \approx 2.3\times 10^{17}\ \mathrm{kg\,m^{-3}}. \]
One teaspoon of pure nuclear matter would weigh hundreds of millions of tonnes. Nuclei are the densest ordinary objects in the universe.

\textbf{A paradox.} Here is the puzzle that should bother every thoughtful student. The nucleus contains several protons packed into a ball a few femtometres wide. But protons all carry charge $+e$, and Coulomb's law says like charges repel with a force
\[ F_{\text{elec}} = \frac{e^2}{4\pi\varepsilon_0 r^2} \approx \frac{9\times 10^{9}\times (1.6\times 10^{-19})^2}{(2\times 10^{-15})^2} \approx 58\ \mathrm{N}, \]
an enormous force on a particle of mass $10^{-27}\ \mathrm{kg}$! Why does the nucleus not fly apart instantly? There must exist a third force of nature, stronger than electricity at these distances: the \cle{strong nuclear force}.

\begin{propriete}[The strong nuclear force]
The strong nuclear force acts between nucleons (proton--proton, proton--neutron, neutron--neutron alike) and has the following features:
\begin{enumerate}
\item \textbf{Attractive} and about $100$ times stronger than the electrostatic repulsion at distances of order $1\ \mathrm{fm}$;
\item \textbf{Very short range}: it is essentially zero beyond about $2$ to $3\ \mathrm{fm}$;
\item \textbf{Charge-independent}: it does not distinguish protons from neutrons;
\item \textbf{Saturating}: a nucleon binds strongly only to its few nearest neighbours.
\end{enumerate}
(The precise mathematical form of this force is not required at GCE Advanced Level; the four properties above are examinable.)
\end{propriete}

\courssub{The Origin of Nuclear Instability: a Competition of Forces}

The stability of a nucleus is the outcome of a \emph{competition}:
\begin{itemize}
\item the \textbf{strong force} attracts every nucleon towards its immediate neighbours and holds the nucleus together;
\item the \textbf{electrostatic force} makes every proton repel \emph{every other proton} (Coulomb's force has infinite range), and pushes the nucleus apart.
\end{itemize}

Because the strong force saturates (each nucleon feels only its neighbours) while the Coulomb repulsion grows with the \emph{total} number of proton pairs, the repulsion wins more and more as $Z$ increases. This single fact explains the whole architecture of nuclear stability:

\begin{itemize}
\item \textbf{Light nuclei} ($Z \lesssim 20$): a stable nucleus has roughly equal numbers of protons and neutrons, $N \approx Z$. Examples: $\ensuremath{\mathrm{^{4}_{2}He}}$, $\ensuremath{\mathrm{^{12}_{6}C}}$, $\ensuremath{\mathrm{^{16}_{8}O}}$, $\ensuremath{\mathrm{^{40}_{20}Ca}}$.
\item \textbf{Heavy nuclei}: extra neutrons are needed. Neutrons add strong-force attraction \emph{without} adding any Coulomb repulsion --- they act like a nuclear ``cement'' diluting the proton repulsion. Stable heavy nuclei therefore have $N > Z$: for $\ensuremath{\mathrm{^{208}_{82}Pb}}$, $N = 126$ against $Z = 82$, i.e.\ $N/Z \approx 1.5$.
\item \textbf{Beyond $Z = 83$}: no amount of neutron cement can compensate the repulsion of so many protons. \emph{All} nuclei with $Z > 83$ are unstable (radioactive), including uranium.
\end{itemize}

\courssub{The Valley of Stability}

Plotting the neutron number $N$ against the proton number $Z$ for all known nuclides gives one of the most instructive charts in physics. The stable nuclides cluster along a narrow band, the \cle{valley of stability}, which starts along the line $N = Z$ and then curves towards $N > Z$.

\begin{popfigure}
\begin{center}
\begin{tikzpicture}
\begin{axis}[
  width=11cm, height=8.5cm,
  xlabel={$Z$ (protons)}, ylabel={$N$ (neutrons)},
  xmin=0, xmax=100, ymin=0, ymax=160,
  xtick={0,20,40,60,80,100},
  ytick={0,40,80,120,160},
  legend style={at={(0.03,0.97)}, anchor=north west, font=\small},
  axis lines=left, clip=false]
% N = Z line
\addplot[domain=0:100, samples=2, dashed, popmuted] {x};
\addlegendentry{$N = Z$}
% Valley of stability (approximate band)
\addplot[domain=0:92, samples=60, very thick, popblue, name path=valley] {x + 0.006*x^2};
\addlegendentry{valley of stability}
% Beta-minus region (above valley)
\addplot[domain=0:92, samples=60, poppink!60, draw=none, fill=popink, fill opacity=0.10] {x + 0.006*x^2 + 30} \closedcycle;
% Labels
\node[text=popblue, font=\small\bfseries, rotate=32] at (axis cs:38,46) {stable nuclei};
\node[text=popink, font=\small, align=center] at (axis cs:30,75) {$\beta^{-}$ region\\ (too many neutrons)};
\node[text=popgreen, font=\small, align=center] at (axis cs:62,40) {$\beta^{+}$ region\\ (too many protons)};
\node[text=poporange, font=\small, align=center] at (axis cs:88,140) {$\alpha$ region\\ (very heavy, $Z > 83$)};
\end{axis}
\end{tikzpicture}
\legende{The $N$--$Z$ chart. Stable nuclides lie in the valley of stability (blue curve). Nuclides above the valley are neutron-rich and decay by $\beta^{-}$ emission; those below are proton-rich and decay by $\beta^{+}$ emission; very heavy nuclides ($Z > 83$) decay mainly by $\alpha$ emission.}
\end{center}
\end{popfigure}

The chart tells us, at a glance, \emph{how} an unstable nucleus will try to reach stability:
\begin{itemize}
\item A nucleus \textbf{above} the valley has too many neutrons. It converts a neutron into a proton by \cle{beta-minus} ($\beta^{-}$) emission, moving one step down towards the valley.
\item A nucleus \textbf{below} the valley has too many protons. It converts a proton into a neutron by \cle{beta-plus} ($\beta^{+}$) emission (or electron capture).
\item A \textbf{very heavy} nucleus ($Z > 83$) sheds bulk by ejecting a $\ensuremath{\mathrm{^{4}_{2}He}}$ nucleus: \cle{alpha} ($\alpha$) emission.
\end{itemize}

\courssub{Radioactivity and Soddy's Conservation Laws}

\begin{definition}[Radioactivity]
\cle{Radioactivity} is the spontaneous disintegration (transformation) of an unstable nucleus into a different, more stable nucleus, accompanied by the emission of radiation ($\alpha$ particles, $\beta$ particles and/or $\gamma$ photons). It is a \emph{nuclear} phenomenon: it depends only on the nucleus, and is unaffected by temperature, pressure, chemical combination or electric fields.
\end{definition}

In a nuclear transformation, the particles rearrange, but two bookkeeping rules always hold. They were stated by Frederick Soddy and Kazimierz Fajans around 1913.

\begin{propriete}[Soddy's conservation laws]
In every nuclear equation, the following are conserved:
\begin{enumerate}
\item the \textbf{total mass number} $A$: the sum of the superscripts is the same on both sides (conservation of the number of nucleons);
\item the \textbf{total charge number} $Z$: the sum of the subscripts is the same on both sides (conservation of electric charge).
\end{enumerate}
For a reaction ${}^{A_1}_{Z_1}X + {}^{A_2}_{Z_2}Y \to {}^{A_3}_{Z_3}X' + {}^{A_4}_{Z_4}Y'$:
\[ A_1 + A_2 = A_3 + A_4 \qquad \text{and} \qquad Z_1 + Z_2 = Z_3 + Z_4. \]
(The proof is examinable in the trivial sense: you must be able to \emph{apply} and justify both equalities in any nuclear equation.)
\end{propriete}

Note carefully: Soddy's laws conserve the \emph{number of nucleons}, not the mass. As we shall see in Section 3, the total mass is \emph{not} conserved --- the small ``mass defect'' is exactly what releases nuclear energy.

\begin{methode}[Checking (or completing) a nuclear equation]
\begin{enumerate}
\item Write the equation with the symbols ${}^{A}_{Z}X$ for every species, including particles: $\alpha = {}^{4}_{2}He$, $\beta^{-} = {}^{0}_{-1}e$, $\beta^{+} = {}^{0}_{+1}e$, neutron $= {}^{1}_{0}n$, proton $= {}^{1}_{1}p$, $\gamma = {}^{0}_{0}\gamma$.
\item \textbf{Balance $A$}: sum the superscripts on the left; the sum of the superscripts on the right must be equal. Solve for any unknown $A$.
\item \textbf{Balance $Z$}: do the same with the subscripts. Solve for any unknown $Z$.
\item \textbf{Identify the element} from $Z$ using the periodic table ($Z$ alone fixes the chemical symbol).
\item \textbf{Re-read} the final equation and double-check both sums.
\end{enumerate}
\end{methode}

\begin{exoresolu}[Completing a decay equation]
Radium-226, once used in hospital sources, decays by alpha emission into radon (Rn, $Z = 86$). Write the balanced equation.
\tcblower
\textbf{Step 1 --- skeleton equation.} The alpha particle is ${}^{4}_{2}He$:
\[ \ce{^{226}_{88}Ra -> ^{A}_{Z}Rn + ^{4}_{2}He}. \]
\textbf{Step 2 --- conservation of $A$:}
\[ 226 = A + 4 \quad\Rightarrow\quad A = 222. \]
\textbf{Step 3 --- conservation of $Z$:}
\[ 88 = Z + 2 \quad\Rightarrow\quad Z = 86. \]
\textbf{Step 4 --- identification.} $Z = 86$ is indeed radon, consistent with the statement. Hence
\[ \boxed{\ \ce{^{226}_{88}Ra -> ^{222}_{86}Rn + ^{4}_{2}He}\ } \]
\textbf{Check:} $226 = 222 + 4$ \checkmark\ and $88 = 86 + 2$ \checkmark. The daughter nucleus has $N = 222 - 86 = 136$ neutrons and lies closer to the valley of stability than radium-226.
\end{exoresolu}

\begin{exoresolu}[Identifying an unknown particle]
When nitrogen-14 in the upper atmosphere is struck by a neutron, carbon-14 is produced together with a second particle. Identify it.
\tcblower
\textbf{Step 1:} $\displaystyle \ce{^{14}_{7}N + ^{1}_{0}n -> ^{14}_{6}C + ^{A}_{Z}X}$.\par
\textbf{Step 2 (balance $A$):} $14 + 1 = 14 + A \Rightarrow A = 1$.\par
\textbf{Step 3 (balance $Z$):} $7 + 0 = 6 + Z \Rightarrow Z = 1$.\par
\textbf{Step 4:} ${}^{1}_{1}X$ is a proton, ${}^{1}_{1}p$ (a hydrogen nucleus). The full equation is
\[ \ce{^{14}_{7}N + ^{1}_{0}n -> ^{14}_{6}C + ^{1}_{1}p}. \]
This is precisely how the carbon-14 used in archaeological dating is continuously created in the atmosphere by cosmic rays.
\end{exoresolu}

\begin{attention}[Frequent examination traps]
\begin{itemize}
\item Writing $\beta^{-}$ as ${}^{0}_{1}e$: the electron has charge $-e$, so its subscript is $-1$, not $+1$. With ${}^{0}_{-1}e$, the charge balance works automatically.
\item Believing that $A$ is conserved ``because mass is conserved'': in nuclear physics mass is \emph{not} conserved (Section 3). What is conserved is the \emph{number of nucleons} and the \emph{charge}.
\item Forgetting the emitted particle altogether when balancing: always check whether an $\alpha$, $\beta$ or $\gamma$ appears on the right-hand side.
\item Confusing the daughter element: after balancing $Z$, look up the symbol --- never keep the parent's symbol with a new $Z$.
\item Ignoring $\gamma$ emission: a $\gamma$ photon (${}^{0}_{0}\gamma$) carries away energy but changes neither $A$ nor $Z$. It often accompanies $\alpha$ or $\beta$ decay but must be included if the question mentions it.
\end{itemize}
\end{attention}

\begin{savaistu}[Radioactivity all around us]
Radioactivity is not an exotic laboratory curiosity. The potassium-40 in bananas and in the volcanic soils of the Mount Cameroon region, the radon gas seeping from granitic rocks, and the cosmic rays that generate carbon-14 high above the Sahel all make natural radioactivity part of everyday life. The same carbon-14, absorbed by plants and then by animals, allows scientists to date fossils and archaeological remains up to about $50\,000$ years old --- a direct application of the nuclear physics of this chapter.
\end{savaistu}

\begin{aretenir}[Key points of this section]
\begin{itemize}
\item The nucleus contains $Z$ protons and $N = A - Z$ neutrons, written ${}^{A}_{Z}X$; nucleons $=$ protons $+$ neutrons.
\item Isotopes: same $Z$, different $N$ (e.g.\ $\ensuremath{\mathrm{^{12}_{6}C}}$ / $\ensuremath{\mathrm{^{14}_{6}C}}$, $\ensuremath{\mathrm{^{235}_{92}U}}$ / $\ensuremath{\mathrm{^{238}_{92}U}}$). Isobars: same $A$. Isotones: same $N$.
\item Nuclear radius $R = r_0 A^{1/3}$ with $r_0 \approx 1.2\ \mathrm{fm}$; nuclear density $\approx 2.3\times 10^{17}\ \mathrm{kg\,m^{-3}}$, the same for all nuclei.
\item The strong nuclear force is attractive, very short-ranged ($\sim 1\ \mathrm{fm}$), charge-independent, and about $100$ times stronger than Coulomb repulsion at contact.
\item Stability results from the strong-versus-Coulomb competition: $N \approx Z$ for light nuclei, $N > Z$ for heavy ones; no stable nucleus exists for $Z > 83$.
\item On the $N$--$Z$ chart, unstable nuclides lie off the valley of stability: $\beta^{-}$ above, $\beta^{+}$ below, $\alpha$ for the very heavy.
\item Radioactivity is the spontaneous disintegration of an unstable nucleus with emission of radiation.
\item Soddy's laws: in every nuclear equation, $\sum A$ and $\sum Z$ are conserved.
\end{itemize}
\end{aretenir}

\begin{exercice}[Practice]
\begin{enumerate}
\item Give the number of protons, neutrons and electrons of the neutral atom $\ensuremath{\mathrm{^{60}_{27}Co}}$ (a source used in radiotherapy).
\item Balance the following equations and name the emitted particle:
\begin{enumerate}
\item $\ensuremath{\mathrm{^{14}_{6}C \rightarrow  ^{14}_{7}N + \dots}}$
\item $\ensuremath{\mathrm{^{238}_{92}U \rightarrow  ^{234}_{90}Th + \dots}}$
\end{enumerate}
\item Explain, in terms of the two competing forces, why stable heavy nuclei need $N > Z$.
\end{enumerate}
\end{exercice}

\begin{corrige}[Exercise 1]
\begin{enumerate}
\item $Z = 27$ protons; $N = A - Z = 60 - 27 = 33$ neutrons; the neutral atom has $27$ electrons.
\item (a) Balance $A$: $14 = 14 + 0$; balance $Z$: $6 = 7 + Z \Rightarrow Z = -1$. The particle is ${}^{0}_{-1}e$: a $\beta^{-}$ emission, $\ce{^{14}_{6}C -> ^{14}_{7}N + ^{0}_{-1}e}$.\par
(b) Balance $A$: $238 = 234 + 4$; balance $Z$: $92 = 90 + 2$. The particle is ${}^{4}_{2}He$: an $\alpha$ emission, $\ce{^{238}_{92}U -> ^{234}_{90}Th + ^{4}_{2}He}$.
\item Each proton repels \emph{all} other protons (long-range Coulomb force), while the strong force binds only nearest neighbours. As $Z$ grows, extra neutrons add strong attraction without adding repulsion, so a larger neutron excess $N > Z$ is needed to hold the nucleus together.
\end{enumerate}
\end{corrige}

\coursec{Types of Radiation, Properties and Applications}

In the previous section we saw that some nuclei are \cle{unstable}: their proton-to-neutron ratio lies outside the band of stability, or they are simply too heavy for the strong nuclear force to hold them together. Such nuclei do not remain unstable forever. Sooner or later, each one \emph{spontaneously} transforms itself into a more stable configuration by emitting particles or energy. This spontaneous transformation is called \cle{radioactivity}, and the emitted particles or photons are called \cle{radiations}. In this section we study the nature of these radiations, their properties, how we detect them, how we protect ourselves from them, and how humanity has learned to put them to work.

\begin{definition}[Radioactivity]
\cle{Radioactivity} is the spontaneous and random transformation of an unstable nucleus into a more stable nucleus, accompanied by the emission of radiation ($\alpha$, $\beta$ or $\gamma$). The original nucleus is called the \cle{parent nucleus} and the nucleus produced is the \cle{daughter nucleus}. The process is:
\begin{itemize}
\item \textbf{spontaneous}: it cannot be triggered, accelerated or slowed down by physical or chemical means (temperature, pressure, catalysts have no effect);
\item \textbf{random}: it is impossible to predict \emph{which} nucleus will decay or \emph{when}; only statistical laws apply to a large sample;
\item \textbf{inevitable}: an unstable nucleus will always decay eventually.
\end{itemize}
\end{definition}

Every nuclear decay obeys two conservation rules, known as \cle{Soddy's laws} (or conservation laws of nuclear reactions):

\begin{propriete}[Soddy's conservation laws]
In any nuclear transformation:
\begin{enumerate}
\item the total \textbf{nucleon number} $A$ is conserved: $\sum A_{\text{before}} = \sum A_{\text{after}}$;
\item the total \textbf{charge number} $Z$ is conserved: $\sum Z_{\text{before}} = \sum Z_{\text{after}}$.
\end{enumerate}
These two laws are sufficient to complete any decay equation. The proof is not examinable, but their \emph{use} is required at the GCE.
\end{propriete}

\courssub{Alpha Radiation}

\textbf{Observation.} When a sample of uranium salt is placed near a photographic plate wrapped in thick black paper, the plate is still darkened. But a single sheet of ordinary paper placed between the source and the plate is enough to stop the effect almost completely. The radiation responsible is easily absorbed and strongly ionising: this is \cle{alpha radiation}.

\begin{definition}[Alpha radiation]
An \cle{alpha particle} ($\alpha$) is a \textbf{helium-4 nucleus} \ensuremath{\mathrm{^{4}_{2}He}}, i.e.\ a cluster of 2 protons and 2 neutrons. Its characteristics are:
\begin{itemize}
\item mass: $m_\alpha \approx 4\ \mathrm{u}$ (about $6.64 \times 10^{-27}\ \mathrm{kg}$);
\item charge: $q_\alpha = +2e = +3.20 \times 10^{-19}\ \mathrm{C}$;
\item typical kinetic energy: a few MeV, giving speeds of order $10^{7}\ \mathrm{m\,s^{-1}}$ (non-relativistic).
\end{itemize}
\end{definition}

Alpha decay occurs in very \textbf{heavy} nuclei ($Z > 82$ roughly), which are unstable simply because they contain too many nucleons. By ejecting an $\alpha$ particle, the nucleus loses 4 nucleons and 2 units of charge and moves closer to stability:
\[ \ce{^{A}_{Z}X -> ^{A-4}_{Z-2}Y + ^{4}_{2}He} \]

\begin{exemplebox}[Alpha decay of uranium-238]
\[ \ce{^{238}_{92}U -> ^{234}_{90}Th + ^{4}_{2}He} \]
Check Soddy's laws: nucleons $238 = 234 + 4$ \checkmark; charge $92 = 90 + 2$ \checkmark. Uranium-238 is the most abundant isotope of natural uranium, mined in several African countries.
\end{exemplebox}

\courssub{Beta-Minus Radiation}

\textbf{Observation.} Some radiations pass easily through paper but are stopped by a few millimetres of aluminium. They are deflected by a magnetic field \emph{in the opposite direction} to alpha particles, and much more strongly. These are \cle{beta-minus particles}.

\begin{definition}[Beta-minus radiation]
A \cle{beta-minus particle} ($\beta^-$) is an \textbf{electron} \ensuremath{\mathrm{^{0}_{-1}e}} emitted by the nucleus. It has:
\begin{itemize}
\item mass: $m_e = 9.11 \times 10^{-31}\ \mathrm{kg}$ (about $1/7300$ of an $\alpha$ particle);
\item charge: $q = -e = -1.60 \times 10^{-19}\ \mathrm{C}$;
\item speed: up to about $0.9\,c$ (relativistic speeds are common).
\end{itemize}
The electron is \emph{not} one of the atom's orbital electrons: it is \textbf{created inside the nucleus} when a neutron transforms into a proton, emitting also an \textbf{antineutrino} $\bar{\nu}_e$ (often omitted in GCE equations but required for energy-momentum conservation):
\[ \ensuremath{\mathrm{^{1}_{0}n \rightarrow  ^{1}_{1}p + ^{0}_{-1}e + \bar{\nu}_e}} \]
\end{definition}

Beta-minus decay occurs in nuclei with an \textbf{excess of neutrons} (below the band of stability). Converting a neutron into a proton increases $Z$ by 1 while $A$ stays constant:
\[ \ce{^{A}_{Z}X -> ^{A}_{Z+1}Y + ^{0}_{-1}e} \]

\begin{exemplebox}[Beta-minus decay of carbon-14]
\[ \ce{^{14}_{6}C -> ^{14}_{7}N + ^{0}_{-1}e} \]
Check: $14 = 14 + 0$ \checkmark; $6 = 7 - 1$ \checkmark. This decay is the basis of carbon-14 dating, discussed below.
\end{exemplebox}

\courssub{Beta-Plus Radiation}

Nuclei with an \textbf{excess of protons} (above the band of stability) can convert a proton into a neutron by emitting a \cle{positron} \ensuremath{\mathrm{^{0}_{+1}e}} — the antiparticle of the electron, with the same mass but charge $+e$, together with a \textbf{neutrino} $\nu_e$:
\[ \ensuremath{\mathrm{^{1}_{1}p \rightarrow  ^{1}_{0}n + ^{0}_{+1}e + \nu_e}} \qquad\text{so}\qquad \ce{^{A}_{Z}X -> ^{A}_{Z-1}Y + ^{0}_{+1}e} \]
Example: $\ce{^{22}_{11}Na -> ^{22}_{10}Ne + ^{0}_{+1}e}$. Positron emitters are used in medicine (PET scanners). At GCE level you must recognise the equation and apply Soddy's laws.

\courssub{Gamma Radiation}

\textbf{Observation.} After an $\alpha$ or $\beta$ decay, the daughter nucleus is often left in an \textbf{excited state} (denoted $^*$), like an atom whose electrons have been raised to a higher level. It gets rid of this surplus energy by emitting a very penetrating radiation that is \emph{not} deflected by any field: \cle{gamma radiation}.

\begin{definition}[Gamma radiation]
A \cle{gamma ray} ($\gamma$) is a \textbf{high-energy electromagnetic photon}, of the same nature as light and X-rays but far more energetic (frequencies above $10^{19}\ \mathrm{Hz}$, energies typically $0.1$ to several MeV). It has:
\begin{itemize}
\item \textbf{no mass} and \textbf{no charge};
\item speed $c = 3.00 \times 10^{8}\ \mathrm{m\,s^{-1}}$;
\item energy $E = hf$.
\end{itemize}
Gamma emission is the \textbf{de-excitation} of a nucleus:
\[ \ensuremath{\mathrm{^{A}_{Z}X^{*} \rightarrow  ^{A}_{Z}X + \gamma}} \]
Neither $A$ nor $Z$ changes: the nucleus keeps its identity, it only loses energy. Gamma emission usually \emph{accompanies} an $\alpha$ or $\beta$ decay.
\end{definition}

\begin{attention}[Common mistakes to avoid]
\begin{itemize}
\item \textbf{Gamma decay does NOT change $A$ or $Z$.} Writing $\ensuremath{\mathrm{^{A}_{Z}X \rightarrow  ^{A-4}_{Z-2}Y + \gamma}}$ is doubly wrong. Only $\alpha$ and $\beta$ decays change the nucleus; $\gamma$ emission is pure energy release.
\item \textbf{The alpha particle is NOT an electron.} It is a helium \emph{nucleus} \ensuremath{\mathrm{^{4}_{2}He}}, about 7300 times heavier than an electron and positively charged. Confusing $\alpha$ with $\beta^-$ in an equation destroys both conservation laws.
\item In $\beta^-$ decay, $A$ stays constant and $Z$ \emph{increases} by 1 — many candidates mistakenly decrease it.
\item In $\beta^+$ decay, $A$ stays constant and $Z$ \emph{decreases} by 1.
\end{itemize}
\end{attention}

\courssub{Comparative Properties of the Radiations}

Because of their different masses and charges, the three radiations interact very differently with matter. A heavy, doubly-charged $\alpha$ particle ploughs through atoms, ripping electrons off (ionising) at every step, and therefore loses its energy very quickly: it is \textbf{strongly ionising but weakly penetrating}. A light $\beta$ particle ionises less and travels further. A neutral $\gamma$ photon interacts only occasionally, so it ionises weakly but penetrates deeply.

\begin{propriete}[Comparative properties of $\alpha$, $\beta$, $\gamma$ radiations]
\begin{center}
\begin{tabularx}{\linewidth}{|l|X|X|X|}
\hline
\rowcolor{popblueL} \textbf{Property} & \textbf{Alpha} $\alpha$ & \textbf{Beta} $\beta^-$ & \textbf{Gamma} $\gamma$ \\
\hline
Nature & \ensuremath{\mathrm{^{4}_{2}He}} nucleus & electron \ensuremath{\mathrm{^{0}_{-1}e}} & EM photon \\
\hline
Mass & $\approx 4\ \mathrm{u}$ & $9.11\times10^{-31}\ \mathrm{kg}$ & $0$ \\
\hline
Charge & $+2e$ & $-e$ & $0$ \\
\hline
Ionising power & very strong ($\approx 10^4$ ion pairs per cm in air) & moderate ($\approx 10^2$) & weak ($\approx 1$) \\
\hline
Penetrating power & weak & moderate & very strong \\
\hline
Range in air & a few cm & a few m & hundreds of m \\
\hline
Stopped by & sheet of paper, skin & a few mm of aluminium & thick lead or concrete (only attenuated) \\
\hline
In $\vec{E}$ or $\vec{B}$ field & deflected, slightly & deflected strongly, opposite way & not deflected \\
\hline
\end{tabularx}
\end{center}
Memory aid: ionising power $\alpha > \beta > \gamma$; penetrating power $\gamma > \beta > \alpha$. The two orders are \emph{reversed}: the better a radiation ionises, the faster it is stopped.
\end{propriete}

\begin{popfigure}
\begin{tikzpicture}[scale=0.95]
% source
\fill[poporange] (0,0) circle (0.18);
\node[below] at (0,-0.25) {\small source};
% beams
\draw[->, very thick, poppurple] (0.2,0) -- (1.6,0) node[midway, above]{\small $\alpha$};
\draw[->, very thick, popblue] (0.2,-0.05) -- (3.4,-0.05) node[midway, above=2pt]{\small $\beta$};
\draw[->, very thick, popgreen] (0.2,-0.1) -- (6.6,-0.1) node[pos=0.62, below=2pt]{\small $\gamma$};
% paper
\fill[popgold!60] (1.7,-1.1) rectangle (1.95,1.1);
\node[below, align=center, text width=1.6cm] at (1.82,-1.15) {\small paper};
% aluminium
\fill[gray!50] (3.5,-1.1) rectangle (3.95,1.1);
\node[below, align=center, text width=1.8cm] at (3.72,-1.15) {\small aluminium (mm)};
% lead
\fill[popink!70] (5.2,-1.1) rectangle (6.05,1.1);
\node[below, align=center, text width=2.2cm] at (5.62,-1.15) {\small lead / concrete (cm)};
% gamma continues, attenuated
\draw[->, thick, popgreen, dashed] (6.05,-0.1) -- (7.6,-0.1) node[right]{\small attenuated};
\end{tikzpicture}
\legende{Penetrating power: $\alpha$ is stopped by paper, $\beta$ by a few millimetres of aluminium, while $\gamma$ is only attenuated by thick lead or concrete.}
\end{popfigure}

In a uniform magnetic field $\vec{B}$ directed into the page, the magnetic force $\vec{F} = q\,\vec{v}\wedge\vec{B}$ curves the trajectory of charged particles into a circular arc of radius $R = \dfrac{mv}{|q|B}$. The $\alpha$ particle (heavy, charge $+2e$) curves \emph{slightly} one way; the $\beta^-$ (light, charge $-e$) curves \emph{sharply} the opposite way; the $\gamma$ photon goes straight through.

\begin{popfigure}
\begin{tikzpicture}[scale=0.9]
% field region
\fill[popblueL] (0,-2.4) rectangle (7,2.4);
\foreach \x in {0.7,1.9,3.1,4.3,5.5,6.7}{
  \foreach \y in {-1.8,-0.6,0.6,1.8}{
    \node[popblue] at (\x,\y) {$\times$};}}
\node[popblue] at (6.4,2.1) {\small $\vec{B}$ into page};
% source
\fill[poporange] (-0.5,0) circle (0.15);
\node[left] at (-0.6,0) {\small source};
% alpha: slight upward curve
\draw[very thick, poppurple, ->] (-0.35,0) .. controls (2,0.35) and (4.5,0.9) .. (7,1.7) node[right]{\small $\alpha$};
% beta: strong downward curve
\draw[very thick, popblue, ->] (-0.35,0) .. controls (1.2,-0.9) and (2.2,-1.9) .. (3.1,-2.4) node[below right]{\small $\beta^-$};
% gamma: straight
\draw[very thick, popgreen, ->] (-0.35,0) -- (7,0) node[right]{\small $\gamma$};
\end{tikzpicture}
\legende{Paths of $\alpha$, $\beta^-$ and $\gamma$ radiations in a uniform magnetic field. The curvature radius $R = mv/(|q|B)$ is large for the heavy $\alpha$, small for the light $\beta^-$, and infinite for the neutral $\gamma$.}
\end{popfigure}

\begin{methode}[Predicting the daughter nucleus with Soddy's laws]
To complete a decay equation $\ensuremath{\mathrm{^{A}_{Z}X \rightarrow  ^{A'}_{Z'}Y + \text{radiation}}}$:
\begin{enumerate}
\item \textbf{Identify the radiation} and write its symbol with its numbers: $\alpha = \ensuremath{\mathrm{^{4}_{2}He}}$, $\beta^- = \ensuremath{\mathrm{^{0}_{-1}e}}$, $\beta^+ = \ensuremath{\mathrm{^{0}_{+1}e}}$, $\gamma = \ensuremath{\mathrm{^{0}_{0}\gamma}}$.
\item \textbf{Conserve nucleons}: $A' = A - A_{\text{radiation}}$.
\item \textbf{Conserve charge}: $Z' = Z - Z_{\text{radiation}}$ (careful with the \emph{negative} charge of $\beta^-$: $Z' = Z - (-1) = Z+1$).
\item \textbf{Identify the element} from $Z'$ using the periodic table, and \textbf{verify} both conservation laws on the final equation.
\end{enumerate}
\end{methode}

\begin{exoresolu}[Completing decay equations (GCE style)]
\textbf{Statement.} (a) Radium-226 $\ensuremath{\mathrm{^{226}_{88}Ra}}$ is an $\alpha$ emitter. Write its decay equation, given that the daughter is radon (Rn). (b) Cobalt-60 $\ensuremath{\mathrm{^{60}_{27}Co}}$ emits a $\beta^-$ particle and then a $\gamma$ photon. Write both equations (daughter: nickel, Ni).
\tcblower
\textbf{Solution.} (a) Alpha emission removes 4 nucleons and 2 charges:
\[ \ce{^{226}_{88}Ra -> ^{226-4}_{88-2}Rn + ^{4}_{2}He} \quad\Longrightarrow\quad \ce{^{226}_{88}Ra -> ^{222}_{86}Rn + ^{4}_{2}He} \]
Check: $226 = 222 + 4$ \checkmark; $88 = 86 + 2$ \checkmark.

(b) Beta-minus emission keeps $A$ constant and increases $Z$ by 1:
\[ \ce{^{60}_{27}Co -> ^{60}_{28}Ni^{*} + ^{0}_{-1}e} \]
The nickel nucleus is left in an excited state and de-excites:
\[ \ensuremath{\mathrm{^{60}_{28}Ni^{*} \rightarrow  ^{60}_{28}Ni + \gamma}} \]
Check: $60 = 60 + 0$ \checkmark; $27 = 28 - 1$ \checkmark. This is exactly the decay used in cobalt-60 radiotherapy.
\end{exoresolu}

\courssub{Detection of Radiations}

Radiations are invisible, but their \textbf{ionising power} betrays them. Three detectors you must know qualitatively:

\begin{itemize}
\item \textbf{Geiger--Müller (GM) counter.} The radiation enters a gas-filled tube through a thin window (mica for $\alpha$) and ionises the gas; each ionisation triggers a brief electrical discharge which is counted (audible ``clicks'' or a digital count rate in counts per second). It detects $\beta$ and $\gamma$ readily; $\alpha$ only if the window is very thin.
\item \textbf{Cloud chamber (Wilson chamber).} Air saturated with alcohol vapour is cooled so the vapour is supersaturated. Ions created along a particle's path act as condensation nuclei: a visible \emph{track of droplets} appears. Alpha particles leave short, thick, straight tracks; $\beta$ particles leave long, thin, wandering tracks; $\gamma$ rays leave almost no track (they are detected only through secondary electrons).
\item \textbf{Photographic film.} Ionising radiation darkens a photographic emulsion exactly like light. Film badges worn by radiographers and hospital staff measure the accumulated dose over weeks.
\end{itemize}

\courssub{Biological Effects and Radiation Protection}

When radiation ionises the molecules of living cells, it can damage DNA, kill cells or trigger cancers. The danger depends on the radiation: \textbf{outside} the body, $\alpha$ particles are stopped by the dead outer layer of skin and are almost harmless, while $\gamma$ rays cross the whole body. \textbf{Inside} the body (inhaled or swallowed source), the situation reverses: the intensely ionising $\alpha$ particles deposit all their energy in a tiny volume of tissue and become the most dangerous.

\begin{aretenir}[The three rules of radiation protection]
\begin{enumerate}
\item \textbf{Distance}: keep as far as possible from the source — the received dose decreases rapidly with distance (roughly as $1/d^2$ for a point source). Laboratory sources are handled with tongs, never with fingers.
\item \textbf{Shielding}: place an absorber between source and person — paper or clothing for $\alpha$, aluminium or plastic for $\beta$, thick lead or concrete for $\gamma$.
\item \textbf{Time}: minimise the \textbf{exposure time}; the dose received is proportional to the duration of exposure.
\end{enumerate}
\end{aretenir}

\courssub{Applications of Radioactivity}

Far from being only a hazard, radioactivity is a powerful tool when used in a controlled way.

\textbf{Medicine.}
\begin{itemize}
\item \textbf{Radiotherapy}: the penetrating $\gamma$ rays of cobalt-60 are aimed at tumours to destroy cancerous cells (cancer treatment units in reference hospitals, e.g.\ in Yaoundé and Douala).
\item \textbf{Radioactive tracers}: a small amount of a short-lived emitter (e.g.\ iodine-131 for the thyroid, technetium-99m for imaging) is introduced into the body; its radiation, detected from outside, reveals how an organ functions.
\item \textbf{Sterilisation}: $\gamma$ rays from cobalt-60 sterilise surgical instruments and syringes sealed in their packaging — no heat, no chemicals.
\end{itemize}

\textbf{Agriculture and food.}
\begin{itemize}
\item \textbf{Food irradiation}: $\gamma$ rays kill bacteria, moulds and insects in stored food (grains, dried fish, spices), greatly extending shelf life without cooking the food.
\item \textbf{Pest control}: male insects sterilised by irradiation are released; mating with them produces no offspring, reducing the pest population.
\item Tracers also follow the uptake of fertilisers by crops, helping optimise their use.
\end{itemize}

\textbf{Industry.}
\begin{itemize}
\item \textbf{Thickness gauging}: a $\beta$ source on one side of a rolling mill and a detector on the other; any change in the thickness of a metal or paper sheet changes the count rate, and the rollers are adjusted automatically.
\item \textbf{Leak detection}: a radioactive tracer injected into a pipeline reveals leaks by appearing at the surface above the fault.
\end{itemize}

\textbf{Carbon-14 dating.} Living organisms continually exchange carbon with the atmosphere, so the proportion of radioactive $\ce{^{14}C}$ in their tissues stays constant while they live. At death, the exchange stops and the $\ce{^{14}C}$ content decreases by $\beta^-$ decay with a half-life $T_{1/2} \approx 5730$ years. Measuring the remaining activity of a wooden tool, a bone or a piece of charcoal and comparing it with that of living matter gives the time elapsed since death:
\[ N = N_0\left(\frac{1}{2}\right)^{t/T_{1/2}} \]
This links radioactivity directly to the law of radioactive decay and half-life studied elsewhere in this chapter — a favourite GCE combination.

\begin{savaistu}[Did you know?]
Radioactivity was discovered in 1896 by Henri Becquerel, who noticed that uranium salts fogged photographic plates kept in a dark drawer. Marie and Pierre Curie then isolated two new radioactive elements, polonium and radium. Rutherford later separated the emissions into ``$\alpha$'' and ``$\beta$'' rays using a magnetic field, and Villard identified the undeflected ``$\gamma$'' rays — the very classification you have just studied.
\end{savaistu}

\begin{attention}[Exam tip — what the GCE Board loves to ask]
GCE Paper 2 questions very frequently ask you to:
\begin{itemize}
\item \textbf{complete a decay equation} (apply Soddy's laws and \emph{show the check} — it earns marks);
\item \textbf{identify a radiation from its properties}: ``stopped by paper'' $\Rightarrow \alpha$; ``deflected opposite to $\alpha$ and more strongly'' $\Rightarrow \beta^-$; ``undeflected, needs lead shielding'' $\Rightarrow \gamma$;
\item \textbf{state one application and one hazard} of a given radioisotope.
\end{itemize}
Learn the comparative table by heart: it answers most of these questions in one line.
\end{attention}

\begin{exercice}[Practice]
\begin{enumerate}
\item Polonium-210 $\ensuremath{\mathrm{^{210}_{84}Po}}$ decays by $\alpha$ emission into lead (Pb). Write the equation.
\item A radiation passes through paper but is stopped by 3 mm of aluminium and is deflected towards the positive plate of an electric field. Identify it and justify.
\item Explain why an $\alpha$ source is safe outside the body but dangerous if swallowed.
\item State one medical and one industrial application of radioactivity, naming the radiation used in each case.
\end{enumerate}
\end{exercice}

\begin{corrige}[Exercise 1]
\textbf{1.} $\ce{^{210}_{84}Po -> ^{206}_{82}Pb + ^{4}_{2}He}$. Check: $210 = 206 + 4$; $84 = 82 + 2$.

\textbf{2.} It is $\beta^-$ radiation: it penetrates paper but not a few mm of aluminium (moderate penetrating power), and its deflection \emph{towards the positive plate} shows a \emph{negative} charge.

\textbf{3.} Outside the body, $\alpha$ particles are stopped by the dead outer layer of skin before reaching living cells. If the source is swallowed or inhaled, the $\alpha$ particles are emitted directly inside living tissue; their very high ionising power then destroys cells in a small volume, causing severe damage.

\textbf{4.} Medical: cobalt-60 radiotherapy, using penetrating $\gamma$ rays to destroy tumours (or iodine-131 tracer, $\beta^-$/$\gamma$). Industrial: thickness gauging of metal sheets using a $\beta$ source and detector (or pipeline leak detection with a tracer).
\end{corrige}

\coursec{Mass Defect and Nuclear Energy}

In the previous section we studied the radiations emitted by unstable nuclei and their applications. A natural question now arises: \emph{where does the energy carried away by these radiations come from?} The answer lies in one of the most profound discoveries of the twentieth century: mass and energy are two faces of the same reality. By carefully \emph{weighing} nuclei, physicists discovered that a bound nucleus is slightly \emph{lighter} than its separated nucleons, and that the ``missing mass'' corresponds exactly to the energy that holds the nucleus together. This is the key to understanding radioactivity, nuclear power and the energy of the stars.

\courssub{The Atomic Mass Unit and the Masses of the Nucleons}

Nuclear masses are extremely small in kilograms, so physicists use a more convenient unit.

\begin{definition}[Atomic mass unit]
The \cle{atomic mass unit}, symbol $\mathrm{u}$, is defined as one twelfth of the mass of a neutral atom of carbon-12, \ensuremath{\mathrm{^{12}_{6}C}}:
\[
1\ \mathrm{u} = \frac{m(\ensuremath{\mathrm{^{12}_{6}C}})}{12} = 1.66054\times 10^{-27}\ \mathrm{kg}.
\]
\end{definition}

The masses of the three particles we shall need are:

\begin{center}
\begin{tabularx}{\linewidth}{|l|X|X|}
\hline
\rowcolor{popblueL} \textbf{Particle} & \textbf{Mass in u} & \textbf{Mass in kg} \\
\hline
Proton & $m_p = 1.00728\ \mathrm{u}$ & $1.6726\times 10^{-27}\ \mathrm{kg}$ \\
\hline
Neutron & $m_n = 1.00866\ \mathrm{u}$ & $1.6749\times 10^{-27}\ \mathrm{kg}$ \\
\hline
Electron & $m_e = 0.00055\ \mathrm{u}$ & $9.11\times 10^{-31}\ \mathrm{kg}$ \\
\hline
\end{tabularx}
\end{center}

Note that the neutron is slightly heavier than the proton: this small difference has important consequences (a free neutron is unstable and undergoes $\beta^-$ decay into a proton, an electron and an antineutrino).

\courssub{Observation: The Nucleus Weighs Less Than Its Parts}

Consider the helium-4 nucleus, \ensuremath{\mathrm{^{4}_{2}He}} (the $\alpha$ particle). It contains $Z = 2$ protons and $A - Z = 2$ neutrons. Let us add the masses of the separated nucleons:
\[
2m_p + 2m_n = 2(1.00728) + 2(1.00866) = 4.03188\ \mathrm{u}.
\]
Now the measured mass of the helium-4 \emph{nucleus} is $m_{\mathrm{nucleus}} = 4.00151\ \mathrm{u}$. The bound nucleus is \emph{lighter} by about $0.03037\ \mathrm{u}$! This is a general experimental fact:

\begin{attention}[A surprising but universal observation]
For \emph{every} stable nucleus, the mass of the nucleus is strictly less than the sum of the masses of its separated nucleons. Mass has ``disappeared'' during the formation of the nucleus. As we shall see, it has in fact been converted into energy and carried away.
\end{attention}

\begin{popfigure}
\begin{center}
\begin{tikzpicture}[scale=1.0]
% Left pan: separated nucleons
\draw[popline, thick] (-4.6,0) -- (-1.4,0);
\draw[popline, thick] (-3,0) -- (-3,1.6);
\draw[popline, thick] (-4.4,1.6) -- (-1.6,1.6);
\draw[popblue, thick] (-4.4,1.6) -- (-4.4,2.0) -- (-1.6,2.0) -- (-1.6,1.6);
\shade[ball color=popink] (-4.0,2.35) circle (0.28);
\shade[ball color=popink] (-3.35,2.35) circle (0.28);
\shade[ball color=popmuted] (-2.65,2.35) circle (0.28);
\shade[ball color=popmuted] (-2.0,2.35) circle (0.28);
\node[align=center, font=\small] at (-3,3.0) {separated nucleons};
\node[align=center, font=\small] at (-3,1.1) {$2m_p + 2m_n = 4.03188\ \mathrm{u}$};
% Fulcrum
\draw[popline, thick, fill=popblueL] (0,0) -- (0.35,0.9) -- (-0.35,0.9) -- cycle;
% Beam tilted
\draw[popdark, very thick] (-1.4,1.25) -- (1.4,0.55);
% Right pan: bound nucleus
\draw[popline, thick] (1.4,0) -- (4.6,0);
\draw[popline, thick] (3,0) -- (3,0.55);
\draw[popline, thick] (1.6,0.55) -- (4.4,0.55);
\draw[poporange, thick] (1.6,0.55) -- (1.6,0.95) -- (4.4,0.95) -- (4.4,0.55);
\shade[ball color=popink] (2.7,1.3) circle (0.28);
\shade[ball color=popmuted] (3.3,1.3) circle (0.28);
\shade[ball color=popmuted] (2.85,1.75) circle (0.28);
\shade[ball color=popink] (3.2,1.75) circle (0.28);
\node[align=center, font=\small] at (3,2.35) {bound nucleus \ensuremath{\mathrm{^{4}_{2}He}}};
\node[align=center, font=\small] at (3,0.05) {$m_{\mathrm{nucleus}} = 4.00151\ \mathrm{u}$};
% Delta m arrow
\draw[->, poppurple, very thick] (0.2,2.6) -- (0.2,1.5);
\node[align=center, font=\small, poppurple] at (1.6,2.9) {missing mass $\Delta m$};
\end{tikzpicture}
\legende{The mass balance: four separated nucleons (2 protons + 2 neutrons) are heavier than the bound helium-4 nucleus. The difference $\Delta m$ is the mass defect.}
\end{center}
\end{popfigure}

\courssub{Mass Defect: Definition}

\begin{definition}[Mass defect]
The \cle{mass defect} $\Delta m$ of a nucleus \ensuremath{\mathrm{^{A}_{Z}X}} is the difference between the total mass of its $Z$ protons and $(A-Z)$ neutrons taken separately, and the measured mass of the nucleus:
\[\boxed{\Delta m = \left[ Z\,m_p + (A - Z)\,m_n \right] - m_{\mathrm{nucleus}}}\]
The mass defect is always \emph{positive} for a bound nucleus. It is expressed in u or in kg.
\end{definition}

\begin{attention}[Nucleus mass or atom mass?]
Data tables often give the mass of the \emph{neutral atom} (nucleus + electrons), not of the bare nucleus. If you use atomic masses, you must use $Z\,m(\ensuremath{\mathrm{^{1}_{1}H}})$ (hydrogen atom = proton + electron) instead of $Z\,m_p$, so that the electron masses cancel. Mixing atomic masses with proton masses is a classic source of error at the GCE.
\end{attention}

\courssub{Einstein's Mass--Energy Equivalence}

In 1905, Einstein showed from the theory of special relativity that mass is a form of energy.

\begin{propriete}[Einstein's mass--energy relation]
Any body of mass $m$ possesses, simply because of its mass, an energy
\[\boxed{E = m c^{2}}\]
where $c = 3.00\times 10^{8}\ \mathrm{m\,s^{-1}}$ is the speed of light in vacuum. Conversely, whenever a system releases energy $\Delta E$, its mass decreases by $\Delta m = \Delta E / c^{2}$.
\end{propriete}

Because $c^2 \approx 9\times 10^{16}\ \mathrm{m^2\,s^{-2}}$ is enormous, a tiny mass corresponds to a huge energy: $1\ \mathrm{kg}$ of matter ``contains'' $9\times 10^{16}\ \mathrm{J}$, roughly the energy produced by a large power station in several years.

\textbf{The conversion factor 1 u $\leftrightarrow$ 931.5 MeV.} Let us compute the energy equivalent of one atomic mass unit, using the precise values $1\ \mathrm{u} = 1.66054\times 10^{-27}\ \mathrm{kg}$ and $c = 2.998\times 10^{8}\ \mathrm{m\,s^{-1}}$:
\[
E = (1\ \mathrm{u})\,c^2 = 1.66054\times 10^{-27} \times (2.998\times 10^{8})^2 = 1.4924\times 10^{-10}\ \mathrm{J}.
\]
Converting to electron-volts ($1\ \mathrm{eV} = 1.60\times 10^{-19}\ \mathrm{J}$):
\[
E = \frac{1.4924\times 10^{-10}}{1.602\times 10^{-19}} = 9.315\times 10^{8}\ \mathrm{eV} = 931.5\ \mathrm{MeV}.
\]

\begin{aretenir}[The key conversion]
\[\boxed{1\ \mathrm{u} \times c^{2} \approx 931.5\ \mathrm{MeV}}\]
This factor lets you convert a mass defect in u directly into an energy in MeV, without passing through kilograms and joules. Memorise it.
\end{aretenir}

\courssub{Binding Energy of the Nucleus}

To separate a nucleus into its individual nucleons, one must supply energy to overcome the strong nuclear attraction. Conversely, when the nucleus forms from separated nucleons, that energy is released (carried away by a photon, for instance), which is exactly why the nucleus ends up lighter.

\begin{definition}[Binding energy]
The \cle{binding energy} $E_l$ of a nucleus is the energy that must be supplied to separate it completely into its constituent nucleons at rest. It equals the energy equivalent of the mass defect:
\[\boxed{E_l = \Delta m \cdot c^{2} = \left[ Z m_p + (A-Z)m_n - m_{\mathrm{nucleus}} \right] c^{2}}\]
It is expressed in joules (SI) or, more conveniently, in MeV.
\end{definition}

\begin{exemplebox}[Binding energy of helium-4]
For \ensuremath{\mathrm{^{4}_{2}He}}, we found $\Delta m = 4.03188 - 4.00151 = 0.03037\ \mathrm{u}$. Hence
\[
E_l = 0.03037 \times 931.5 = 28.3\ \mathrm{MeV} = 28.3\times 10^{6}\times 1.60\times 10^{-19} = 4.53\times 10^{-12}\ \mathrm{J}.
\]
About 28 MeV is needed to tear a helium nucleus apart: this is why $\alpha$ particles are so stable.
\end{exemplebox}

\courssub{Binding Energy per Nucleon and the Aston Curve}

Is a heavy nucleus more stable than a light one simply because its total binding energy is larger? No: a heavy nucleus also has more nucleons to bind. The fair comparison is the \emph{average} binding energy of each nucleon.

\begin{propriete}[Binding energy per nucleon measures stability]
The \cle{binding energy per nucleon} is
\[\boxed{\frac{E_l}{A}}\]
The \emph{larger} $E_l/A$, the \emph{more stable} the nucleus, because more energy per nucleon is needed to dismantle it. The maximum of $E_l/A$ (about $8.8\ \mathrm{MeV}$ per nucleon) occurs near \emph{iron-56}, \ensuremath{\mathrm{^{56}_{26}Fe}}, which is therefore one of the most stable nuclei in nature.
\end{propriete}

The graph of $E_l/A$ against the mass number $A$ is called the \cle{Aston curve}.

\begin{popfigure}
\begin{center}
\begin{tikzpicture}
\begin{axis}[
    width=12cm, height=8cm,
    xlabel={Mass number $A$},
    ylabel={$E_l/A$ (MeV per nucleon)},
    xmin=0, xmax=250, ymin=0, ymax=10,
    xtick={0,50,100,150,200,250},
    ytick={0,2,4,6,8,10},
    grid=both, grid style={popline, very thin},
    axis line style={popdark},
    tick label style={font=\small},
    label style={font=\small},
]
\addplot[popblue, very thick, smooth] coordinates {
(2,1.1) (4,7.1) (6,5.3) (12,7.7) (16,8.0) (20,8.0) (28,8.4) (40,8.6)
(56,8.8) (75,8.7) (100,8.6) (120,8.5) (140,8.4) (160,8.2) (180,8.1)
(200,7.9) (220,7.8) (238,7.6) (250,7.5)
};
\addplot[only marks, mark=*, poporange, mark size=2.5pt] coordinates {(56,8.8)};
\node[poporange, font=\small, anchor=south west] at (axis cs:56,8.8) {\ce{^{56}Fe}};
\addplot[only marks, mark=*, poppurple, mark size=2.5pt] coordinates {(238,7.6)};
\node[poppurple, font=\small, anchor=south west] at (axis cs:200,7.35) {\ce{^{235}U} region};
\addplot[only marks, mark=*, popgreen, mark size=2.5pt] coordinates {(2,1.1)};
\node[popgreen, font=\small, anchor=south west] at (axis cs:4,1.3) {\ce{^{2}H}};
\draw[dashed, popline] (axis cs:56,0) -- (axis cs:56,8.8);
\draw[dashed, popline] (axis cs:0,8.8) -- (axis cs:56,8.8);
\end{axis}
\end{tikzpicture}
\legende{The Aston curve: binding energy per nucleon versus mass number. The maximum ($\approx 8.8$ MeV/nucleon) is near iron-56. Very light nuclei (left) can gain stability by \emph{fusion}; very heavy nuclei (right) by \emph{fission}.}
\end{center}
\end{popfigure}

\textbf{Consequences of the shape of the curve:}
\begin{itemize}
\item \emph{Heavy nuclei} ($A \gtrsim 200$, e.g.\ uranium-235) sit on the descending part. If such a nucleus \emph{splits} into two medium nuclei (fission), the products lie higher on the curve: they are more stable, and the difference in binding energy is released.
\item \emph{Light nuclei} ($A \lesssim 20$, e.g.\ \ce{^{2}H}, \ce{^{3}H}) sit on the steep rising part. If two light nuclei \emph{merge} (fusion), the product is more stable and energy is released. This powers the Sun.
\item \emph{Medium nuclei} near iron are at the top: neither fission nor fusion of iron releases energy. Iron is the ``nuclear ash'' of stars.
\end{itemize}

\courssub{Energy Balance of a Nuclear Reaction}

Any nuclear reaction (radioactive decay, fission, fusion, bombardment) conserves the total mass--energy. If the products are lighter than the reactants, the missing mass has been converted into kinetic energy of the products and into radiation.

\begin{propriete}[Energy released by a nuclear reaction]
The \cle{energy balance} (or $Q$-value) of a nuclear reaction is
\[\boxed{Q = \left( m_{\text{reactants}} - m_{\text{products}} \right) c^{2}}\]
\begin{itemize}
\item If $Q > 0$: the reaction is \cle{exoenergetic} (exothermic); it releases energy and can occur spontaneously.
\item If $Q < 0$: the reaction is \cle{endoenergetic} (endothermic); it can only occur if the missing energy is supplied, for example as kinetic energy of an incoming projectile.
\end{itemize}
(The proof is a direct application of conservation of total energy $E = mc^2 + E_k$; it is instructive but not required as a formal proof at the GCE.)
\end{propriete}

\begin{methode}[Computing the energy released by a nuclear reaction]
\begin{enumerate}
\item Write the balanced nuclear equation (conservation of $A$ and $Z$).
\item Look up the masses (in u) of all reactants and all products.
\item Compute $\Delta m = m_{\text{reactants}} - m_{\text{products}}$ in u.
\item Convert to MeV: $Q\ (\mathrm{MeV}) = \Delta m\ (\mathrm{u}) \times 931.5$.
\item Convert to joules if required: $Q\ (\mathrm{J}) = Q\ (\mathrm{MeV}) \times 10^{6} \times 1.60\times 10^{-19}$.
\item Conclude: $Q > 0$ means energy is released; $Q < 0$ means energy must be supplied.
\end{enumerate}
\end{methode}

\begin{exoresolu}[Energy released in the $\alpha$ decay of radium-226]
The radium-226 nucleus decays according to
\[
\ce{^{226}_{88}Ra -> ^{222}_{86}Rn + ^{4}_{2}He}.
\]
Given: $m(\ensuremath{\mathrm{^{226}_{88}Ra}}) = 225.9771\ \mathrm{u}$; $m(\ensuremath{\mathrm{^{222}_{86}Rn}}) = 221.9703\ \mathrm{u}$; $m(\ensuremath{\mathrm{^{4}_{2}He}}) = 4.0015\ \mathrm{u}$. Calculate the energy released in MeV and in joules.
\tcblower
\textbf{Step 1 -- Mass defect of the reaction:}
\[
\Delta m = m_{\text{reactants}} - m_{\text{products}} = 225.9771 - (221.9703 + 4.0015) = 0.0053\ \mathrm{u}.
\]
\textbf{Step 2 -- Conversion to MeV:}
\[
Q = 0.0053 \times 931.5 = 4.9\ \mathrm{MeV}.
\]
\textbf{Step 3 -- Conversion to joules:}
\[
Q = 4.9\times 10^{6} \times 1.60\times 10^{-19} = 7.9\times 10^{-13}\ \mathrm{J}.
\]
\textbf{Conclusion:} $Q > 0$, so the decay is exoenergetic: radium-226 decays spontaneously. The energy appears essentially as the kinetic energy of the $\alpha$ particle (the radon nucleus recoils slightly). Although $7.9\times 10^{-13}\ \mathrm{J}$ per decay seems tiny, one gram of radium contains about $2.7\times 10^{21}$ nuclei, so the energy released per gram is macroscopic.
\end{exoresolu}

\begin{exoresolu}[Binding energy per nucleon of carbon-12]
The mass of the carbon-12 nucleus is $m_{\mathrm{nucleus}} = 11.99671\ \mathrm{u}$ (obtained by subtracting the six electron masses from the atomic mass, exactly $12\ \mathrm{u}$). Calculate the binding energy per nucleon of \ensuremath{\mathrm{^{12}_{6}C}}.
\tcblower
\textbf{Step 1 -- Mass defect:} carbon-12 has $Z = 6$ protons and $A - Z = 6$ neutrons:
\[
\Delta m = \left[6(1.00728) + 6(1.00866)\right] - 11.99671 = 12.09564 - 11.99671 = 0.09893\ \mathrm{u}.
\]
\textbf{Step 2 -- Binding energy:}
\[
E_l = 0.09893 \times 931.5 = 92.2\ \mathrm{MeV}.
\]
\textbf{Step 3 -- Per nucleon:}
\[
\frac{E_l}{A} = \frac{92.2}{12} = 7.7\ \mathrm{MeV/nucleon}.
\]
This is below the iron-56 maximum of $8.8\ \mathrm{MeV/nucleon}$, consistent with the position of carbon on the rising part of the Aston curve.
\end{exoresolu}

\begin{attention}[Common mistakes to avoid]
\begin{itemize}
\item \textbf{Forgetting the MeV $\to$ J conversion.} The SI unit of energy is the joule. If a question asks for the energy in joules, multiply the value in MeV by $1.60\times 10^{-13}$ (since $1\ \mathrm{MeV} = 10^{6}\times 1.60\times 10^{-19}\ \mathrm{J}$). Leaving the answer in MeV costs marks.
\item \textbf{Mixing atomic and nuclear masses.} Either use \emph{nuclear} masses with $m_p$ everywhere, or use \emph{atomic} masses with the mass of the hydrogen atom $m(\ensuremath{\mathrm{^{1}_{1}H}})$ everywhere. Never mix the two conventions in the same calculation.
\item \textbf{Sign of $\Delta m$.} The mass defect is reactants \emph{minus} products. A negative $Q$ does not mean ``energy destroyed''; it means the reaction needs an energy input.
\item \textbf{Confusing $E_l$ and $E_l/A$.} A large total binding energy does \emph{not} imply great stability: uranium has a huge $E_l$ but is unstable because its $E_l/A$ is low. Always compare binding energies \emph{per nucleon}.
\item \textbf{Rounding too early.} Mass defects are tiny differences between large numbers; keep at least four decimal places in the masses (in u) until the final step.
\end{itemize}
\end{attention}

\begin{savaistu}[Why the Sun shines]
Deep in the Sun's core, four hydrogen nuclei fuse (through a chain of reactions) into one helium-4 nucleus. The helium nucleus is lighter than the four protons by about $0.028\ \mathrm{u}$, so each fusion releases about $26\ \mathrm{MeV}$. The Sun converts roughly 600 million tonnes of hydrogen per second, and ``burns'' about 4 million tonnes of \emph{mass} into pure energy every second, in perfect agreement with $E = mc^2$. The light that warms the farms of the North West Region left the Sun as mass defect!
\end{savaistu}

\begin{exercice}[Uranium-235 and the Aston curve]
Given $m(\ensuremath{\mathrm{^{235}_{92}U}}) = 234.9934\ \mathrm{u}$ (nucleus), $m_p = 1.00728\ \mathrm{u}$, $m_n = 1.00866\ \mathrm{u}$.
\begin{enumerate}
\item Calculate the mass defect and the binding energy of the uranium-235 nucleus in MeV.
\item Deduce its binding energy per nucleon, and compare with iron-56 ($8.8\ \mathrm{MeV/nucleon}$). What do you conclude about the possibility of releasing energy by fission of uranium-235?
\end{enumerate}
\end{exercice}

\begin{corrige}[Exercise 1]
\textbf{1.} Uranium-235 has $Z = 92$ protons and $A - Z = 143$ neutrons:
\[
\Delta m = 92(1.00728) + 143(1.00866) - 234.9934 = 92.6698 + 144.2384 - 234.9934 = 1.9148\ \mathrm{u}.
\]
\[
E_l = 1.9148 \times 931.5 \approx 1783.6\ \mathrm{MeV} \approx 1.78\ \mathrm{GeV}.
\]
\textbf{2.} Binding energy per nucleon:
\[
\frac{E_l}{A} = \frac{1783.6}{235} \approx 7.6\ \mathrm{MeV/nucleon}.
\]
This is well below the $8.8\ \mathrm{MeV/nucleon}$ of iron-56: uranium-235 lies on the descending branch of the Aston curve. If it splits into two medium-mass fragments (each with $E_l/A \approx 8.5\ \mathrm{MeV/nucleon}$), each nucleon gains about $0.9\ \mathrm{MeV}$ of binding, so the fission of one nucleus releases roughly $235 \times 0.9 \approx 200\ \mathrm{MeV}$: fission is exoenergetic, as we shall study in the next section.
\end{corrige}

\begin{aretenir}[Key points of this section]
\begin{itemize}
\item Masses: $m_p = 1.00728\ \mathrm{u}$, $m_n = 1.00866\ \mathrm{u}$; $1\ \mathrm{u} = 1.66054\times 10^{-27}\ \mathrm{kg}$.
\item Mass defect: $\Delta m = [Zm_p + (A-Z)m_n] - m_{\mathrm{nucleus}} > 0$.
\item Einstein: $E = mc^2$; binding energy $E_l = \Delta m\,c^2$; conversion $1\ \mathrm{u} \leftrightarrow 931.5\ \mathrm{MeV}$.
\item Stability is measured by $E_l/A$; maximum $\approx 8.8\ \mathrm{MeV/nucleon}$ at iron-56 (Aston curve).
\item Energy balance: $Q = (m_{\text{reactants}} - m_{\text{products}})c^2$; $Q > 0$ exoenergetic, $Q < 0$ endoenergetic.
\end{itemize}
\end{aretenir}

We now understand \emph{why} certain nuclei can release energy: heavy nuclei by splitting, light nuclei by merging. The next section examines these two processes in detail: nuclear fission, harnessed in power reactors, and nuclear fusion, the energy source of the stars.

\coursec{Nuclear Fission and Nuclear Fusion}

In the previous section we established Einstein's mass--energy relation $E = \Delta m\,c^{2}$ and the \cle{Aston curve} of binding energy per nucleon. We saw that the curve rises steeply for light nuclei, reaches a maximum near iron ($A \approx 56$, about $8.8\ \mathrm{MeV}$ per nucleon), and then decreases gently for heavy nuclei. This shape has a spectacular practical consequence: there are \emph{two} ways of rearranging nucleons so that the products sit higher on the curve (more tightly bound), and therefore release energy. A very heavy nucleus can \emph{split} into medium-sized fragments --- this is \cle{fission}. Two very light nuclei can \emph{join} into a heavier one --- this is \cle{fusion}. These two processes power nuclear reactors, the Sun, and the most destructive weapons ever built. This section studies them in turn and then compares them.

\courssub{Nuclear fission}

\textbf{Observation.} In 1938, Otto Hahn and Fritz Strassmann bombarded uranium with neutrons and were astonished to find barium --- a much lighter element --- among the products. Lise Meitner and Otto Frisch correctly interpreted the result: the uranium nucleus had been \emph{split in two}. How can a slow, electrically neutral particle break apart a nucleus containing 92 protons? Precisely because it is neutral, the neutron is not repelled by the nuclear charge, so it can approach and be absorbed; the resulting compound nucleus \ensuremath{\mathrm{^{236}_{92}U}} is highly unstable and deforms like a liquid drop until it breaks into two pieces.

\begin{definition}[Nuclear fission]
\cle{Nuclear fission} is a nuclear reaction in which a \textbf{heavy nucleus} ($A \gtrsim 230$) splits into \textbf{two lighter nuclei} (the \emph{fission fragments}), with the emission of a few \textbf{neutrons} and the release of a large amount of \textbf{energy}. Fission of \ensuremath{\mathrm{^{235}_{92}U}} is usually \emph{induced} by the capture of a slow (\emph{thermal}) neutron, of kinetic energy of order $0.025\ \mathrm{eV}$. (A few very heavy nuclei also undergo \emph{spontaneous} fission, without any projectile, but this is rare.)
\end{definition}

A typical induced fission of uranium-235 is:
\[ \ce{^{235}_{92}U + ^{1}_{0}n -> ^{94}_{38}Sr + ^{140}_{54}Xe + 2^{1}_{0}n} \]
Another frequently observed channel is:
\[ \ce{^{235}_{92}U + ^{1}_{0}n -> ^{92}_{36}Kr + ^{141}_{56}Ba + 3^{1}_{0}n} \]

\begin{attention}[Fission is not unique]
Uranium-235 can split in many different ways (more than thirty fragment pairs are known). \textbf{All} fission equations must conserve charge number $Z$ and nucleon number $A$. Check the first equation: $235 + 1 = 94 + 140 + 2 = 236$ \checkmark\ and $92 + 0 = 38 + 54 + 0 = 92$ \checkmark. The number of emitted neutrons depends on the channel (usually 2 or 3). Note also that the fragments are generally \textbf{radioactive}: being neutron-rich, they decay by $\beta^{-}$ emission until they reach stability --- this is why used reactor fuel remains dangerous for a very long time.
\end{attention}

\begin{popfigure}
\begin{center}
\begin{tikzpicture}[scale=1.0, every node/.style={font=\small}]
% neutron
\fill[poporange] (-3.2,0) circle (0.14);
\node[above] at (-3.2,0.2) {$n$ (slow)};
\draw[->, thick, poporange] (-3.0,0) -- (-1.6,0);
% U-235 nucleus
\shade[ball color=popblue] (-0.6,0) circle (0.85);
\node[white] at (-0.6,0) {\ce{^{235}U}};
\draw[->, thick, popink] (0.4,0) -- (1.4,0);
% fragments
\shade[ball color=popgreen] (2.6,0.9) circle (0.55);
\node[white] at (2.6,0.9) {\ce{^{94}Sr}};
\shade[ball color=poppurple] (2.6,-0.9) circle (0.65);
\node[white] at (2.6,-0.9) {\ce{^{140}Xe}};
% emitted neutrons
\fill[poporange] (4.2,0.25) circle (0.12);
\fill[poporange] (4.2,-0.25) circle (0.12);
\draw[->, thick, poporange] (3.4,0.1) -- (4.05,0.25);
\draw[->, thick, poporange] (3.4,-0.1) -- (4.05,-0.25);
\node[right] at (4.35,0) {$2n$};
% energy
\draw[->, very thick, popgold] (1.9,-1.8) -- (3.6,-1.8);
\node[below, popgold] at (2.75,-1.85) {energy $\approx 200\ \mathrm{MeV}$};
\end{tikzpicture}
\legende{Induced fission of uranium-235 by a slow neutron: two fragments, two or three neutrons and about $200\ \mathrm{MeV}$ of energy are released.}
\end{center}
\end{popfigure}

\textbf{Why does fission release energy?} The fragments (Sr, Xe, Kr, Ba, \dots) have $A$ in the range 90--145, a region where the binding energy per nucleon ($\approx 8.5\ \mathrm{MeV}$) is \emph{greater} than that of uranium ($\approx 7.6\ \mathrm{MeV}$). The products are therefore more tightly bound: the total mass of the products is \emph{less} than the initial mass, and the difference appears as kinetic energy of the fragments and neutrons, plus $\gamma$ radiation. A detailed mass balance gives about $\SI{200}{MeV}$ per fission --- roughly $3.2 \times 10^{-11}\ \mathrm{J}$ per nucleus, which is enormous compared with the few electron-volts released per atom in a chemical reaction.

\courssub{The chain reaction and critical mass}

\textbf{Observation.} Each fission needs \emph{one} neutron but \emph{produces two or three}. If these neutrons in turn strike other \ce{^{235}U} nuclei, the number of fissions multiplies generation after generation: $1 \to 3 \to 9 \to 27 \to \dots$ This self-sustaining process is the \cle{chain reaction}.

\begin{definition}[Multiplication factor and critical mass]
The \cle{multiplication factor} $k$ is the ratio of the number of fissions in one generation to the number in the previous generation. The \cle{critical mass} is the minimum mass of fissile material needed for a self-sustaining chain reaction ($k = 1$): below it, too many neutrons escape through the surface or are captured uselessly, and the reaction dies out.
\end{definition}

\begin{center}
\begin{tabularx}{\linewidth}{|l|X|X|}
\hline
\rowcolor{popblueL} Regime & Condition & Consequence \\
\hline
Subcritical & $k < 1$ & Reaction dies out. \\
\hline
Critical & $k = 1$ & Steady, controlled power (nuclear reactor). \\
\hline
Supercritical & $k > 1$ & Explosive growth (atomic bomb). \\
\hline
\end{tabularx}
\end{center}

\begin{popfigure}
\begin{center}
\begin{tikzpicture}[scale=0.95, every node/.style={font=\small}]
% generation 1
\shade[ball color=popblue] (0,0) circle (0.45);
\node[white] at (0,0) {\ce{^{235}U}};
\fill[poporange] (-1.6,0) circle (0.11);
\draw[->, thick, poporange] (-1.45,0) -- (-0.5,0);
% arrows to generation 2
\draw[->, thick, popink] (0.35,0.3) -- (1.6,1.3);
\draw[->, thick, popink] (0.45,0) -- (1.6,0);
\draw[->, thick, popink] (0.35,-0.3) -- (1.6,-1.3);
% generation 2
\shade[ball color=popblue] (2.3,1.5) circle (0.4);
\shade[ball color=popblue] (2.3,0) circle (0.4);
\shade[ball color=popblue] (2.3,-1.5) circle (0.4);
% arrows to generation 3
\foreach \y in {1.5,0,-1.5}{
  \draw[->, thick, popink] (2.65,\y+0.15) -- (4.0,\y+0.75);
  \draw[->, thick, popink] (2.7,\y) -- (4.0,\y);
  \draw[->, thick, popink] (2.65,\y-0.15) -- (4.0,\y-0.75);
}
% generation 3
\foreach \y in {2.25,1.5,0.75,0,-0.75,-1.5,-2.25}{
  \shade[ball color=popgreen] (4.6,\y) circle (0.28);
}
% labels
\node[below] at (0,-0.7) {generation 1};
\node[below] at (2.3,-2.1) {generation 2};
\node[below] at (4.6,-2.6) {generation 3};
\node at (6.1,0) {$1 \to 3 \to 9 \to \dots$};
\end{tikzpicture}
\legende{Chain reaction: neutrons from each fission trigger further fissions, so the number of reactions grows geometrically.}
\end{center}
\end{popfigure}

In an \textbf{atomic bomb}, two subcritical masses of \ce{^{235}U} (or plutonium-239) are slammed together to form a highly supercritical mass: the chain reaction runs away in about a microsecond, releasing an enormous amount of energy. In a \textbf{nuclear reactor}, the goal is exactly the opposite: keep $k = 1$ permanently so that energy is produced at a constant, controlled rate.

\textbf{Main parts of a pressurised water reactor (PWR).}
\begin{itemize}
\item \textbf{Fuel:} rods of uranium slightly enriched in \ce{^{235}U}, where fission takes place.
\item \textbf{Moderator} (ordinary water): slows the fast fission neutrons down to thermal speeds, because slow neutrons are far more likely to induce further fission of \ce{^{235}U}.
\item \textbf{Control rods} (cadmium or boron): absorb neutrons. Pushing them in lowers $k$; withdrawing them raises $k$. They keep the reactor exactly critical and allow emergency shutdown.
\item \textbf{Coolant:} the pressurised water also carries the heat away to a steam generator; the steam drives a turbine coupled to an alternator, exactly as in a thermal power station.
\item \textbf{Containment:} a thick steel and concrete vessel that protects the environment from radiation and radioactive leaks.
\end{itemize}

\begin{savaistu}[Nuclear energy and Africa]
Only South Africa currently operates a commercial nuclear power station (Koeberg, near Cape Town), but several African countries, including Ghana, Nigeria and Egypt, are developing nuclear programmes, and many operate small \emph{research reactors} used for medicine, agriculture and materials analysis. Cameroon, whose electricity demand grows every year and whose grid relies heavily on hydroelectric dams vulnerable to drought, takes part in international training and follows International Atomic Energy Agency (IAEA) standards for radiation safety and radioactive waste management. Any future nuclear option would require strict solutions for the long-term storage of highly radioactive fission products.
\end{savaistu}

\courssub{Nuclear fusion}

\textbf{Observation.} The Sun has been radiating about $4 \times 10^{26}\ \mathrm{W}$ for billions of years. No chemical or fission fuel could last so long. Its energy comes from the opposite process: light nuclei (hydrogen) \emph{fusing} into heavier ones (helium).

\begin{definition}[Nuclear fusion]
\cle{Nuclear fusion} is a nuclear reaction in which \textbf{two light nuclei unite} to form a \textbf{heavier nucleus}, with the emission of a particle (often a neutron) and the release of \textbf{energy}. Example (the deuterium--tritium reaction):
\[ \ce{^{2}_{1}H + ^{3}_{1}H -> ^{4}_{2}He + ^{1}_{0}n} \]
\end{definition}

\begin{popfigure}
\begin{center}
\begin{tikzpicture}[scale=1.0, every node/.style={font=\small}]
\shade[ball color=popblue] (-2.6,0.7) circle (0.4);
\node[white] at (-2.6,0.7) {\ce{^{2}H}};
\shade[ball color=poppurple] (-2.6,-0.7) circle (0.45);
\node[white] at (-2.6,-0.7) {\ce{^{3}H}};
\draw[->, thick, popink] (-2.1,0.6) -- (-0.9,0.15);
\draw[->, thick, popink] (-2.1,-0.6) -- (-0.9,-0.15);
\node at (-1.5,0.55) {$T \approx 10^{8}\ \mathrm{K}$};
\draw[->, very thick, popdark] (-0.6,0) -- (0.6,0);
\shade[ball color=popgreen] (1.6,0) circle (0.55);
\node[white] at (1.6,0) {\ce{^{4}He}};
\fill[poporange] (3.4,0) circle (0.13);
\draw[->, thick, poporange] (2.3,0) -- (3.2,0);
\node[right] at (3.55,0) {$n$};
\draw[->, very thick, popgold] (1.0,-1.1) -- (2.8,-1.1);
\node[below, popgold] at (1.9,-1.15) {$17.6\ \mathrm{MeV}$};
\end{tikzpicture}
\legende{Fusion of deuterium and tritium into helium-4, releasing a neutron and $17.6\ \mathrm{MeV}$.}
\end{center}
\end{popfigure}

\textbf{Why is fusion so difficult to achieve?} Both nuclei are positively charged, so they repel each other electrically (Coulomb repulsion). To bring them close enough ($\approx 10^{-15}\ \mathrm{m}$) for the strong nuclear force to bind them, they must collide with enormous kinetic energy. This requires temperatures of about $10^{8}\ \mathrm{K}$, at which matter is a \emph{plasma} (fully ionised gas) that no ordinary container can hold. In \textbf{stars}, gravity provides the confinement and the core temperature ignites the fusion of hydrogen into helium. In the \textbf{hydrogen bomb}, a small fission bomb is used as a match to reach the ignition temperature --- an \emph{uncontrolled} fusion. \emph{Controlled} fusion on Earth is the goal of the international \textbf{ITER} project (France), which confines the plasma with powerful magnetic fields; it aims to demonstrate that fusion can produce more energy than it consumes, but commercial fusion power remains a long-term hope.

\begin{propriete}[Fission and fusion both release energy]
Both the fission of a heavy nucleus and the fusion of light nuclei release energy because, in each case, the \textbf{products lie higher on the Aston curve} (greater binding energy per nucleon) than the reactants. The total mass decreases and the mass defect appears as energy, $E = \Delta m\,c^{2}$. (Not a proof to reproduce in the examination, but you must be able to \emph{justify} it from the curve.)
\end{propriete}

\begin{methode}[Energy released by a fission or fusion reaction]
\begin{enumerate}
\item Write the balanced nuclear equation (conserve $A$ and $Z$).
\item Compute the mass defect: $\Delta m = m_{\text{reactants}} - m_{\text{products}}$ (use the atomic or nuclear masses as given in the question).
\item Convert to energy: $E = \Delta m\,c^{2}$, or directly $E\ (\mathrm{MeV}) = \Delta m\ (\mathrm{u}) \times 931.5\ \mathrm{MeV/u}$.
\item Quote the result with its unit; a positive $\Delta m$ means energy is \emph{released}.
\end{enumerate}
\end{methode}

\begin{exoresolu}[Energy released by the D--T fusion reaction]
Given $m(\ce{^{2}H}) = 2.0141\ \mathrm{u}$, $m(\ce{^{3}H}) = 3.0160\ \mathrm{u}$, $m(\ce{^{4}He}) = 4.0026\ \mathrm{u}$, $m_{n} = 1.0087\ \mathrm{u}$ and $1\ \mathrm{u} \equiv 931.5\ \mathrm{MeV}$, calculate the energy released in $\ce{^{2}H + ^{3}H -> ^{4}He + n}$.
\tcblower
\textbf{Solution.}
Mass of reactants: $2.0141 + 3.0160 = 5.0301\ \mathrm{u}$.
Mass of products: $4.0026 + 1.0087 = 5.0113\ \mathrm{u}$.
Mass defect:
\[ \Delta m = 5.0301 - 5.0113 = 0.0188\ \mathrm{u}. \]
Energy released:
\[ E = 0.0188 \times 931.5 \approx 17.5\ \mathrm{MeV} \approx 2.8 \times 10^{-12}\ \mathrm{J}. \]
This is smaller than the $\approx 200\ \mathrm{MeV}$ of one fission, but per \emph{kilogram of fuel} fusion releases several times more energy, because the reacting nuclei are so light.
\end{exoresolu}

\begin{exoresolu}[Balancing a fission equation]
Complete the fission equation $\ensuremath{\mathrm{^{235}_{92}U + ^{1}_{0}n \rightarrow  ^{92}_{36}Kr + ^{141}_{56}Ba + x\,^{1}_{0}n}}$ by finding $x$.
\tcblower
\textbf{Solution.} Conservation of nucleon number:
\[ 235 + 1 = 92 + 141 + x \quad\Rightarrow\quad x = 236 - 233 = 3. \]
Conservation of charge: $92 = 36 + 56$ \checkmark. Hence \textbf{three neutrons} are emitted.
\end{exoresolu}

\courssub{Fission versus fusion: a comparison}

\begin{center}
\begin{tabularx}{\linewidth}{|l|X|X|}
\hline
\rowcolor{popblueL} Criterion & Fission & Fusion \\
\hline
Fuel & \ce{^{235}U}, \ce{^{239}Pu}: scarce, mined & \ce{^{2}H} from sea water: almost unlimited \\
\hline
Energy per kg of fuel & Very high & Even higher (several times) \\
\hline
Waste & Long-lived radioactive fission products & Little radioactive waste (activated structures) \\
\hline
Safety & Risk of runaway ($k>1$), meltdown & No chain reaction; reaction stops if confinement fails \\
\hline
Conditions & Achieved at ordinary temperatures & Needs $T \approx 10^{8}\ \mathrm{K}$: not yet industrial \\
\hline
Uses & Power reactors, atomic bomb & Stars, hydrogen bomb; ITER research \\
\hline
\end{tabularx}
\end{center}

\begin{attention}[Common mistakes]
\begin{itemize}
\item Saying that \textbf{fusion} occurs in today's nuclear power plants. All operating reactors use \textbf{fission}; controlled fusion is still experimental (ITER).
\item Confusing the \textbf{critical mass} with the \emph{total mass of fuel}: the critical mass is the \emph{minimum} mass for a self-sustaining chain reaction, not the amount loaded into a reactor.
\item Forgetting that fission equations have several possible fragment pairs --- never ``correct'' the given fragments; only find the number of neutrons.
\item Mixing up the two reactions: fission = heavy nucleus \emph{splits}; fusion = light nuclei \emph{join}.
\item Writing $E = \Delta m\,c^{2}$ with $\Delta m$ in atomic mass units and $c$ in SI units without converting: either convert $\Delta m$ to kilograms, or use the shortcut $1\ \mathrm{u} \equiv 931.5\ \mathrm{MeV}$.
\end{itemize}
\end{attention}

\begin{methode}[Exam tip: finding the number of emitted neutrons]
In a fission equation with unknown neutron number $x$, apply conservation of nucleon number: $x = (A_{\text{fuel}} + 1) - (A_{\text{fragment 1}} + A_{\text{fragment 2}})$, then verify that charge is conserved. This is a frequent two-mark question at the GCE.
\end{methode}

\begin{aretenir}[Key points]
\begin{itemize}
\item \textbf{Fission}: a heavy nucleus (\ce{^{235}U}) captures a slow neutron and splits into two fragments plus 2--3 neutrons, releasing $\approx 200\ \mathrm{MeV}$.
\item The emitted neutrons sustain a \textbf{chain reaction}; $k = 1$ in a reactor (controlled), $k > 1$ in a bomb (uncontrolled); a \textbf{critical mass} is required.
\item A PWR uses fuel rods, a moderator, control rods, a coolant and a containment vessel.
\item \textbf{Fusion}: two light nuclei (\ce{^{2}H} + \ce{^{3}H}) unite at $T \approx 10^{8}\ \mathrm{K}$; it powers the stars and the hydrogen bomb; ITER seeks controlled fusion.
\item Both processes release energy because the products are higher on the Aston curve; compute $E = \Delta m\,c^{2}$ with $\Delta m$ in u and $1\ \mathrm{u} \equiv 931.5\ \mathrm{MeV}$.
\end{itemize}
\end{aretenir}

\begin{exercice}[Energy from uranium fission]
One fission of \ce{^{235}U} releases on average $200\ \mathrm{MeV}$.
\begin{enumerate}
\item Express this energy in joules.
\item Calculate the energy released by the complete fission of $1.0\ \mathrm{g}$ of \ce{^{235}U} (molar mass $235\ \mathrm{g\,mol^{-1}}$, $N_{A} = 6.02 \times 10^{23}\ \mathrm{mol^{-1}}$).
\item Compare with the combustion of petrol ($\approx 4.5 \times 10^{7}\ \mathrm{J\,kg^{-1}}$) and comment.
\end{enumerate}
\end{exercice}

\begin{corrige}[Exercise 1]
\textbf{1.} $E = 200 \times 10^{6} \times 1.60 \times 10^{-19} = 3.2 \times 10^{-11}\ \mathrm{J}$.

\textbf{2.} Number of nuclei in $1.0\ \mathrm{g}$:
\[ N = \frac{1.0}{235} \times 6.02 \times 10^{23} = 2.56 \times 10^{21}\ \text{nuclei}. \]
Total energy: $E_{\text{tot}} = 2.56 \times 10^{21} \times 3.2 \times 10^{-11} \approx 8.2 \times 10^{10}\ \mathrm{J}$.

\textbf{3.} Burning $1\ \mathrm{g}$ of petrol releases about $4.5 \times 10^{4}\ \mathrm{J}$. Fission of the same mass of uranium releases roughly \textbf{two million times more} energy, which is why nuclear fuel is so concentrated --- and why its handling demands extreme care.
\end{corrige}

\begin{exercice}[Energy released in a fission reaction]
Consider the fission $\ce{^{235}_{92}U + ^{1}_{0}n -> ^{94}_{38}Sr + ^{140}_{54}Xe + 2^{1}_{0}n}$. Given $m(\ce{^{235}U}) = 235.0439\ \mathrm{u}$, $m(\ce{^{94}Sr}) = 93.9154\ \mathrm{u}$, $m(\ce{^{140}Xe}) = 139.9252\ \mathrm{u}$, $m_{n} = 1.0087\ \mathrm{u}$ and $1\ \mathrm{u} \equiv 931.5\ \mathrm{MeV}$:
\begin{enumerate}
\item Calculate the mass defect $\Delta m$ for this reaction.
\item Deduce the energy released per fission, in MeV, and compare with the average value of $200\ \mathrm{MeV}$.
\end{enumerate}
\end{exercice}

\begin{corrige}[Exercise 2]
\textbf{1.} Mass of reactants: $235.0439 + 1.0087 = 236.0526\ \mathrm{u}$.
Mass of products: $93.9154 + 139.9252 + 2 \times 1.0087 = 235.8580\ \mathrm{u}$.
\[ \Delta m = 236.0526 - 235.8580 = 0.1946\ \mathrm{u}. \]

\textbf{2.} $E = 0.1946 \times 931.5 \approx 181\ \mathrm{MeV}$. This is of the same order as the average $200\ \mathrm{MeV}$; the exact value depends on the fission channel, and part of the energy is carried away later by the radioactive decay of the fragments.
\end{corrige}

\newpage

\coursec{Integration activity}

\courssub{Aims of the activity}
This integration activity places you in the role of energy consultants for \textbf{ENEO} (Energy of Cameroon), the national electricity utility. You will mobilise the full toolkit of nuclear physics developed in this chapter — nuclear equations, mass defect calculations, binding energy curves, and radiation safety principles — to evaluate two nuclear options for diversifying Cameroon's energy mix, currently dominated by hydroelectricity. By the end of this activity, you should be able to:
\begin{itemize}
    \item Balance nuclear reactions for fission and fusion using conservation of nucleon number $A$ and charge (proton) number $Z$.
    \item Calculate the energy released per reaction and per unit mass of fuel using the mass--energy equivalence $\Delta E = \Delta m c^2$ and the conversion $1\,\mathrm{u} = 931.5\,\mathrm{MeV}/c^2$.
    \item Interpret the Aston curve (binding energy per nucleon versus mass number) to explain qualitatively why both very heavy and very light nuclei can release energy.
    \item Synthesise technical, safety, environmental and economic arguments into a concise policy memo suitable for decision-makers.
\end{itemize}

\courssub{Situation}
Cameroon's electricity demand is growing steadily, by several per cent per year. While hydroelectric projects (Lom Pangar, Nachtigal, Memve'ele) provide a low-carbon base, seasonal variability of river flows and the need for firm baseload power lead planners to ask whether nuclear energy could one day play a role. You are part of a multidisciplinary team of young physicists and engineers seconded to \textbf{ENEO}'s Strategic Planning Division. Your mandate: produce a one-page technical memo comparing the viability of a \textbf{Pressurised Water Reactor (PWR) fission plant} fuelled by enriched uranium versus a hypothetical \textbf{Deuterium--Tritium (D--T) fusion pilot plant} for the Yaoundé region (peak demand $\sim 400\,\mathrm{MW}$).

The following nuclear data are provided (atomic masses in unified atomic mass units, u):
\begin{center}
\begin{tabularx}{\linewidth}{|l|X|}\hline
\rowcolor{popblueL}
\textbf{Particle / Nucleus} & \textbf{Atomic mass (u)} \\ \hline
Neutron ($^{1}_{0}n$) & $1.008\,665$ \\
Uranium-235 ($^{235}_{92}\mathrm{U}$) & $235.043\,930$ \\
Krypton-95 ($^{95}_{36}\mathrm{Kr}$) & $94.939\,710$ \\
Barium-140 ($^{140}_{56}\mathrm{Ba}$) & $139.910\,581$ \\
Deuterium ($^{2}_{1}\mathrm{H}$) & $2.014\,102$ \\
Tritium ($^{3}_{1}\mathrm{H}$) & $3.016\,049$ \\
Helium-4 ($^{4}_{2}\mathrm{He}$) & $4.002\,603$ \\ \hline
\end{tabularx}
\end{center}
\textbf{Constants:} $1\,\mathrm{u} = 931.5\,\mathrm{MeV}/c^2$; $N_A = 6.022 \times 10^{23}\,\mathrm{mol^{-1}}$; $1\,\mathrm{eV} = 1.602 \times 10^{-19}\,\mathrm{J}$, so $1\,\mathrm{MeV} = 1.602 \times 10^{-13}\,\mathrm{J}$.

\courssub{Tasks}
\begin{enumerate}
    \item \textbf{Balance the nuclear equations.}
    \begin{enumerate}
        \item Write the balanced equation for the induced fission of $^{235}_{92}\mathrm{U}$ by a thermal neutron producing $^{95}_{36}\mathrm{Kr}$, $^{140}_{56}\mathrm{Ba}$ and $x$ neutrons. Determine $x$.
        \item Write the balanced equation for the fusion of deuterium and tritium ($^{2}_{1}\mathrm{H} + ^{3}_{1}\mathrm{H} \rightarrow ^{4}_{2}\mathrm{He} + n$). Verify the conservation laws.
    \end{enumerate}
    \item \textbf{Energy yield calculations.}
    \begin{enumerate}
        \item Calculate the mass defect $\Delta m$ (in u) and the energy $Q$ (in MeV) released \emph{per reaction} for both the fission and the fusion processes above.
        \item Calculate the energy released \emph{per kilogram of fuel} for each process. For fission, assume pure $^{235}\mathrm{U}$ fuel. For fusion, assume an equimolar D--T mixture (effective molar mass $\approx 5.03\,\mathrm{g\,mol^{-1}}$).
        \item Compare the specific energy densities (in $\mathrm{J\,kg^{-1}}$) and comment on the orders of magnitude.
    \end{enumerate}
    \item \textbf{Aston curve justification.}
    Using the general shape of the binding energy per nucleon ($B/A$) curve (Aston curve), explain \emph{why} both reactions are exothermic. Sketch a qualitative $B/A$ versus $A$ graph indicating the approximate positions of the initial and final nuclei for each reaction. Show that the products lie higher on the curve (more tightly bound) than the reactants.
    \item \textbf{Policy memo (one page max).}
    Draft a memo to the Director General of ENEO recommending an energy strategy. Your memo must include:
    \begin{itemize}
        \item A clear recommendation (Fission now / Fusion R\&D / Neither / Hybrid).
        \item Quantitative support from Tasks 2 and 3.
        \item Radiation safety considerations: nature of waste (half-lives, activity), shielding requirements, and proliferation risk.
        \item Cameroonian context: uranium availability in the Central African region, cooling water constraints (Sanaga river), grid stability, regulatory capacity, and public acceptance.
    \end{itemize}
    Groups will present and defend their memo in a 10-minute plenary session.
\end{enumerate}

\courssub{Model answer}

\textbf{Task 1: Balancing the nuclear equations}

\begin{methode}[Balancing any nuclear equation]
\begin{enumerate}
    \item Write the reaction with unknown coefficients.
    \item Conserve the \cle{nucleon number} $A$ (top numbers): sum of $A$ on the left $=$ sum of $A$ on the right.
    \item Conserve the \cle{charge number} $Z$ (bottom numbers): sum of $Z$ on the left $=$ sum of $Z$ on the right.
    \item Solve for the unknown(s) and rewrite the final balanced equation.
\end{enumerate}
\end{methode}

\begin{enumerate}
    \item \textbf{Fission:} Conservation of nucleon number $A$ and proton number $Z$:
    \[
    ^{235}_{92}\mathrm{U} + ^{1}_{0}n \rightarrow ^{95}_{36}\mathrm{Kr} + ^{140}_{56}\mathrm{Ba} + x\,^{1}_{0}n
    \]
    \begin{align*}
        A:&\quad 235 + 1 = 95 + 140 + x \implies 236 = 235 + x \implies \boxed{x = 1} \\
        Z:&\quad 92 + 0 = 36 + 56 + 0 \implies 92 = 92 \quad \text{(conserved)}
    \end{align*}
    Balanced equation:
    \begin{center}
    \ce{^{235}_{92}U + ^{1}_{0}n -> ^{95}_{36}Kr + ^{140}_{56}Ba + ^{1}_{0}n}
    \end{center}
    \begin{attention}[Common mistake]
    Many candidates assume $x=2$ or $3$ because \emph{on average} a $^{235}\mathrm{U}$ fission event releases 2--3 neutrons. Here the specific fragment pair ($^{95}\mathrm{Kr}$, $^{140}\mathrm{Ba}$) fixes $x=1$ by nucleon conservation. Always balance the equation \emph{as given}, not the statistical average.
    \end{attention}

    \item \textbf{Fusion:}
    \[
    ^{2}_{1}\mathrm{H} + ^{3}_{1}\mathrm{H} \rightarrow ^{4}_{2}\mathrm{He} + ^{1}_{0}n
    \]
    Check: $A$: $2+3 = 4+1 = 5$; $Z$: $1+1 = 2+0 = 2$. Both conserved.
\end{enumerate}

\textbf{Task 2: Energy yield calculations}

\begin{aretenir}[Recipe for a $Q$-value]
\begin{itemize}
    \item $\Delta m = m_{\text{initial}} - m_{\text{final}}$ (mass that ``disappears'').
    \item $Q = \Delta m \times 931.5\,\mathrm{MeV/u}$ if masses are in u, or $Q = \Delta m\, c^2$ in SI.
    \item If $\Delta m > 0$, the reaction is \cle{exothermic} (releases energy).
\end{itemize}
\end{aretenir}

\begin{enumerate}
    \item \textbf{Mass defect and $Q$-value per reaction.}

    \textbf{Fission:}
    \begin{align*}
        m_{\text{initial}} &= m(^{235}\mathrm{U}) + m_n = 235.043\,930 + 1.008\,665 = 236.052\,595\,\mathrm{u} \\
        m_{\text{final}} &= m(^{95}\mathrm{Kr}) + m(^{140}\mathrm{Ba}) + m_n \\
        &= 94.939\,710 + 139.910\,581 + 1.008\,665 = 235.858\,956\,\mathrm{u} \\
        \Delta m_{\text{fis}} &= m_{\text{initial}} - m_{\text{final}} = 236.052\,595 - 235.858\,956 = 0.193\,639\,\mathrm{u} \\
        Q_{\text{fis}} &= \Delta m_{\text{fis}} \times 931.5\,\mathrm{MeV/u} = 0.193\,639 \times 931.5 \approx \boxed{180.4\,\mathrm{MeV}}
    \end{align*}

    \textbf{Fusion:}
    \begin{align*}
        m_{\text{initial}} &= m(^{2}\mathrm{H}) + m(^{3}\mathrm{H}) = 2.014\,102 + 3.016\,049 = 5.030\,151\,\mathrm{u} \\
        m_{\text{final}} &= m(^{4}\mathrm{He}) + m_n = 4.002\,603 + 1.008\,665 = 5.011\,268\,\mathrm{u} \\
        \Delta m_{\text{fus}} &= 5.030\,151 - 5.011\,268 = 0.018\,883\,\mathrm{u} \\
        Q_{\text{fus}} &= 0.018\,883 \times 931.5 \approx \boxed{17.59\,\mathrm{MeV}}
    \end{align*}

    \textbf{Comment.} One fission releases about ten times more energy than one fusion \emph{per reaction} — but a uranium nucleus is about 47 times heavier than a D--T pair. Per unit \emph{mass} of fuel, fusion will win: this is what we check next.

    \item \textbf{Energy per kilogram of fuel.}

    \textbf{Fission (pure $^{235}\mathrm{U}$):}
    Molar mass $M_{\mathrm{U}} \approx 235\,\mathrm{g\,mol^{-1}} = 0.235\,\mathrm{kg\,mol^{-1}}$.
    Number of nuclei per kg: $N = \dfrac{N_A}{M_{\mathrm{U}}} = \dfrac{6.022 \times 10^{23}}{0.235} = 2.563 \times 10^{24}\,\mathrm{kg^{-1}}$.
    \begin{align*}
        E_{\text{fis, kg}} &= N \times Q_{\text{fis}} \times 1.602 \times 10^{-13}\,\mathrm{J/MeV} \\
        &= 2.563 \times 10^{24} \times 180.4 \times 1.602 \times 10^{-13} \\
        &\approx \boxed{7.41 \times 10^{13}\,\mathrm{J\,kg^{-1}}}
    \end{align*}

    \textbf{Fusion (D--T mixture, $M_{\text{mix}} \approx 5.03\,\mathrm{g\,mol^{-1}}$):}
    One reaction consumes one D and one T nucleus, i.e. one ``mole of reactions'' consumes $5.03\,\mathrm{g}$ of equimolar mixture.
    Number of reactions per kg: $N' = \dfrac{N_A}{0.00503} = \dfrac{6.022 \times 10^{23}}{5.03 \times 10^{-3}} = 1.197 \times 10^{26}\,\mathrm{kg^{-1}}$.
    \begin{align*}
        E_{\text{fus, kg}} &= N' \times Q_{\text{fus}} \times 1.602 \times 10^{-13} \\
        &= 1.197 \times 10^{26} \times 17.59 \times 1.602 \times 10^{-13} \\
        &\approx \boxed{3.37 \times 10^{14}\,\mathrm{J\,kg^{-1}}}
    \end{align*}

    \item \textbf{Comparison:}
    \[
    \frac{E_{\text{fus, kg}}}{E_{\text{fis, kg}}} \approx \frac{3.37 \times 10^{14}}{7.41 \times 10^{13}} \approx 4.6
    \]
    Fusion releases about \textbf{4.6 times more energy per unit mass of fuel} than fission. Both are roughly $10^6$--$10^7$ times more energy-dense than chemical combustion (petrol $\approx 4.6 \times 10^{7}\,\mathrm{J\,kg^{-1}}$; coal $\approx 3 \times 10^{7}\,\mathrm{J\,kg^{-1}}$).
\end{enumerate}

\textbf{Task 3: Aston curve justification}

\begin{aretenir}[The Aston Curve (Binding Energy per Nucleon)]
The curve of $B/A$ versus $A$ peaks near $^{56}\mathrm{Fe}$ ($B/A \approx 8.8\,\mathrm{MeV}$ per nucleon), the most tightly bound region of the nuclear chart.
\begin{itemize}
    \item \textbf{Fission:} a heavy nucleus ($A \approx 235$, $B/A \approx 7.6\,\mathrm{MeV}$) splits into medium-mass fragments ($A \approx 95$ and $140$, $B/A \approx 8.5\,\mathrm{MeV}$). The products are more tightly bound $\Rightarrow$ energy is released.
    \item \textbf{Fusion:} very light nuclei ($A=2,3$; $B/A \approx 1.1$ and $2.8\,\mathrm{MeV}$) fuse into $^{4}\mathrm{He}$ ($B/A \approx 7.1\,\mathrm{MeV}$). The product is much more tightly bound $\Rightarrow$ large energy release.
\end{itemize}
In both cases the reaction moves the system \textbf{towards the peak at $A \approx 56$}, increasing the average binding energy per nucleon. The mass defect $\Delta m$ corresponds exactly to this gain in total binding energy: $\Delta m c^2 = B_{\text{products}} - B_{\text{reactants}}$.
\end{aretenir}

\begin{popfigure}
\begin{tikzpicture}[scale=1.0]
    \draw[->] (0,0) -- (13,0) node[right, font=\small] {$A$};
    \draw[->] (0,0) -- (0,6.2) node[above, font=\small] {$B/A$ (MeV)};
    % x-axis ticks: x = A/20
    \foreach \a/\lab in {2.8/56, 5/{100}, 7/{140}, 10/{200}, 11.75/235}
        \draw (\a,0.06) -- (\a,-0.06) node[below, font=\tiny] {\lab};
    \foreach \y in {2,4,6,8}
        \draw (0.06,\y*0.6) -- (-0.06,\y*0.6) node[left, font=\tiny] {\y};
    % Aston curve, x = A/20, y = (B/A)*0.6
    \draw[popblue, thick, smooth] plot coordinates {
        (0.1,0.66) (0.15,1.68) (0.2,4.26) (0.5,4.4) (1,4.6) (1.5,4.9) (2.8,5.28) (4.75,5.16) (7,5.04) (10,4.8) (11.75,4.56)
    };
    % Points
    \fill[poppink] (0.125,1.2) circle (1.6pt) node[above, font=\tiny, poppink] {$^{2,3}\mathrm{H}$};
    \fill[popgreen] (0.2,4.26) circle (1.6pt) node[above right, font=\tiny, popgreen] {$^{4}\mathrm{He}$};
    \fill[popdark] (2.8,5.28) circle (1.6pt) node[above, font=\tiny, popdark] {$^{56}\mathrm{Fe}$};
    \fill[poporange] (5.9,5.1) circle (1.6pt) node[above, font=\tiny, poporange] {$^{95}\mathrm{Kr},\,^{140}\mathrm{Ba}$};
    \fill[poppurple] (11.75,4.56) circle (1.6pt) node[above, font=\tiny, poppurple] {$^{235}\mathrm{U}$};
    % Arrows
    \draw[poppink, ->, thick] (0.35,1.5) .. controls (0.5,3) .. (0.35,4.1);
    \node[poppink, font=\tiny] at (1.1,2.6) {Fusion};
    \draw[poppurple, ->, thick] (11.5,4.75) .. controls (9,5.35) .. (6.3,5.3);
    \node[poppurple, font=\tiny] at (9,5.6) {Fission};
\end{tikzpicture}
\legende{Qualitative Aston curve (horizontal scale: $x = A/20$). Fusion (light $\to$ medium) and fission (heavy $\to$ medium) both climb towards the iron peak, i.e. towards more tightly bound nuclei, releasing energy.}
\end{popfigure}

\textbf{Task 4: Policy Memo — Model Response}

\begin{center}
\fbox{\begin{minipage}{0.92\linewidth}
\textbf{MEMORANDUM}\\
\textbf{TO:} Director General, ENEO\\
\textbf{FROM:} Nuclear Physics Advisory Group, Upper Sixth Cohort\\
\textbf{SUBJECT:} Nuclear Energy Option for Yaoundé Baseload — Fission vs. Fusion

\textbf{1. Executive Recommendation}
We recommend a \textbf{phased strategy}: a \textbf{feasibility study for a Small Modular Reactor (SMR) fission plant} as a long-term option, while devoting a small fraction of the energy R\&D budget to \textbf{international fusion partnerships}. A domestic fusion pilot plant is not realistic in the foreseeable future: no fusion reactor anywhere in the world yet produces net electricity.

\textbf{2. Quantitative Energy Density Argument}
Our calculations (Task 2) give:
\begin{itemize}
    \item Fission ($^{235}\mathrm{U}$): $7.4 \times 10^{13}\,\mathrm{J\,kg^{-1}}$.
    \item Fusion (D--T): $3.4 \times 10^{14}\,\mathrm{J\,kg^{-1}}$, i.e. $\approx 4.6$ times more per kg.
\end{itemize}
For a $400\,\mathrm{MW}_e$ plant at $\sim 33\%$ thermal efficiency, the thermal power is $\approx 1.2 \times 10^{9}\,\mathrm{W}$, i.e. $\approx 1.0 \times 10^{14}\,\mathrm{J}$ per day. Fuel consumption (assuming complete burn-up):
\begin{itemize}
    \item Fission: $\dfrac{1.0 \times 10^{14}}{7.4 \times 10^{13}} \approx 1.4\,\mathrm{kg}$ of $^{235}\mathrm{U}$ per day.
    \item Fusion: $\dfrac{1.0 \times 10^{14}}{3.4 \times 10^{14}} \approx 0.3\,\mathrm{kg}$ of D--T mixture per day.
\end{itemize}
Either way, a few hundred kilograms of fuel per year replace millions of tonnes of fossil fuel. The decisive difference is \textbf{technological maturity}: PWRs are proven, licensed technology with hundreds of units operating worldwide; D--T fusion remains at the experimental stage (research tokamaks; first net-energy demonstrations only very recent and not yet at plant scale).

\textbf{3. Radiation Safety \& Waste Profile}
\begin{center}
\begin{tabularx}{\linewidth}{|l|X|X|}\hline
\rowcolor{popblueL}
\textbf{Aspect} & \textbf{Fission (PWR)} & \textbf{Fusion (D--T)} \\ \hline
\textbf{High-level waste} & Spent fuel contains long-lived actinides (e.g. $^{239}\mathrm{Pu}$, $T_{1/2} \approx 24\,000$ yr) and fission products ($^{137}\mathrm{Cs}$, $^{90}\mathrm{Sr}$, $T_{1/2} \approx 30$ yr). Requires deep geological storage. & No actinides. Structural materials activated by $14.1\,\mathrm{MeV}$ neutrons; with low-activation alloys, waste decays to safe levels in about a century. \\ \hline
\textbf{Radioactive inventory} & Large (core + storage pool). Residual decay heat requires active cooling long after shutdown. & Moderate (tritium inventory of order kilograms, activated structures). No runaway decay-heat scenario. \\ \hline
\textbf{Proliferation} & Enrichment and reprocessing are dual-use technologies; international safeguards mandatory. & No fissile material; tritium handling is the main concern. Lower proliferation profile. \\ \hline
\textbf{Shielding} & Thick concrete and steel against $\gamma$ rays and neutrons. & Thick blanket/shield against $14.1\,\mathrm{MeV}$ neutrons; comparable footprint. \\ \hline
\end{tabularx}
\end{center}

\textbf{4. Cameroonian Context}
\begin{itemize}
    \item \textbf{Fuel security:} uranium deposits are known in the Central African region (e.g. Bakouma in neighbouring CAR); a regional supply partnership would be needed since Cameroon has no enrichment industry. Deuterium is extractable from water; tritium must be bred from lithium — a technology not yet demonstrated at scale.
    \item \textbf{Cooling:} a $\sim 1.2\,\mathrm{GW}_{th}$ plant needs a few tens of $\mathrm{m^3\,s^{-1}}$ of cooling water; the Sanaga is a major river, but siting studies must respect ecological flow requirements, especially in the dry season.
    \item \textbf{Grid:} Cameroon's interconnected grid is of the order of $1$--$2\,\mathrm{GW}$. A single large $1000\,\mathrm{MW}$ unit would violate the N-1 security criterion (loss of the largest unit must not collapse the grid). \textbf{SMRs (50--300 MW)} fit far better.
    \item \textbf{Regulatory:} Cameroon has a radiation protection framework and a national radiation protection agency, but nuclear \emph{licensing} capacity (following the IAEA Milestones Approach) would take 10--15 years to build.
    \item \textbf{Public acceptance:} a transparent waste-management plan and an independent regulator are preconditions for any programme.
\end{itemize}

\textbf{5. Conclusion}
Fission via the SMR route is the only nuclear option that could, in principle, be licensed and built within 15--20 years; fusion should be followed through international cooperation and local training of plasma physicists. \textbf{Neither} option removes the urgent priorities: grid reinforcement, completion of the hydroelectric programme, and solar PV integration.

\textit{End of Memo}
\end{minipage}}
\end{center}

\begin{attention}[Examiner's traps in the plenary defence]
\begin{itemize}
    \item \textbf{Atomic vs. nuclear masses:} the given masses are \emph{atomic} masses (they include the electrons). In the fission and fusion equations above, the electron numbers balance automatically (92 e$^-$ on each side for fission; 2 e$^-$ on each side for fusion), so atomic masses can be used directly. In $\beta^+$ decay or electron capture this is \emph{not} automatic — always check electron balance.
    \item \textbf{Per kg of \emph{what}?} Fission: per kg of $^{235}\mathrm{U}$. Fusion: per kg of the \emph{D--T mixture}. Do not quote per kg of tritium alone.
    \item \textbf{Aston curve axes:} the $x$-axis is the mass number $A$, not the atomic number $Z$; the peak is at $A \approx 56$ (Fe/Ni region).
    \item \textbf{Units:} energy per reaction in MeV; energy per kg in $\mathrm{J\,kg^{-1}}$. Convert with $1\,\mathrm{MeV} = 1.602 \times 10^{-13}\,\mathrm{J}$.
    \item \textbf{Sign of $\Delta m$:} define $\Delta m = m_{\text{initial}} - m_{\text{final}}$ so that an exothermic reaction gives $\Delta m > 0$ and $Q > 0$. Mixing up the order flips the sign and costs marks.
\end{itemize}
\end{attention}

\newpage

\coursec{Methods and worked examples}

\coursec{Atomic Nucleus and Nuclear Stability}

\begin{savaistu}[Discovery of the Nucleus]
In 1911, Ernest Rutherford, Hans Geiger and Ernest Marsden performed the famous gold foil experiment at Manchester. They observed that while most $\alpha$ particles passed through a thin gold foil with little deflection, a few were scattered at large angles ($> 90^\circ$). Rutherford concluded: \emph{``It was quite the most incredible event that has ever happened to me in my life. It was almost as incredible as if you fired a 15-inch shell at a piece of tissue paper and it came back and hit you.''} This led to the nuclear model: a tiny, dense, positively charged nucleus containing almost all the mass, surrounded by electrons.
\end{savaistu}

\begin{definition}[Nuclear Composition and Notation]
An atomic nucleus is characterised by:
\begin{itemize}
\item \textbf{Proton number (Atomic number) $Z$:} Number of protons. Determines the chemical element.
\item \textbf{Neutron number $N$:} Number of neutrons.
\item \textbf{Nucleon number (Mass number) $A = Z + N$:} Total number of nucleons (protons + neutrons).
\end{itemize}
A nuclide is denoted $\ensuremath{\mathrm{^{A}_{Z}X}}$ (e.g., $\ensuremath{\mathrm{^{238}_{92}U}}$). \textbf{Isotopes} have same $Z$, different $N$ (and $A$). \textbf{Isotones} have same $N$. \textbf{Isobars} have same $A$.
\end{definition}

\begin{definition}[Unified Atomic Mass Unit]
The unified atomic mass unit (\textbf{u}) is defined as $\frac{1}{12}$ the mass of a neutral $\ensuremath{\mathrm{^{12}_{6}C}}$ atom.
\[
1\ \text{u} = \SI{1.660539e-27}{kg} = \SI{931.494}{MeV/c^2}
\]
Masses of fundamental particles:
$m_p = \SI{1.007276}{u}$, $m_n = \SI{1.008665}{u}$, $m_e = \SI{0.00054858}{u}$.
Atomic mass $M_{\text{atom}}$ includes $Z$ electrons: $M_{\text{atom}} \approx M_{\text{nucleus}} + Z m_e - B_e/c^2$ (electron binding energy $B_e$ negligible).
\end{definition}

\begin{propriete}[Nuclear Stability and the $N$-$Z$ Curve]
\begin{itemize}
\item \textbf{Strong Nuclear Force:} Short-range ($\approx 1-3\ \text{fm}$), attractive, charge-independent (acts $p-p$, $n-n$, $p-n$), saturates. Overcomes Coulomb repulsion between protons.
\item \textbf{Stability Belt:} For light nuclei ($A < 40$), stable nuclei have $N \approx Z$. For heavier nuclei, Coulomb repulsion requires excess neutrons ($N > Z$) to provide more strong force without adding repulsion. The line of stability curves upward.
\item \textbf{Even-Even preference:} Nuclei with even $Z$ and even $N$ are most stable (pairing effect). Odd-odd nuclei are rare (only 4 stable: $\ensuremath{\mathrm{^{2}_{1}H}}$, $\ensuremath{\mathrm{^{6}_{3}Li}}$, $\ensuremath{\mathrm{^{10}_{5}B}}$, $\ensuremath{\mathrm{^{14}_{7}N}}$).
\item \textbf{Magic Numbers:} $Z$ or $N = 2, 8, 20, 28, 50, 82, 126$ confer extra stability (shell model).
\end{itemize}
\end{propriete}

\begin{experience}[The Segre Chart (N vs Z)]
Plot $N$ vs $Z$ for known nuclides. Stable nuclides (black dots) form a narrow ``valley of stability''. Neutron-rich nuclei lie above $\to$ $\beta^-$ decay. Proton-rich nuclei lie below $\to$ $\beta^+$ decay or Electron Capture (EC). Very heavy nuclei ($Z > 82$) $\to$ $\alpha$ decay or spontaneous fission. Technetium ($Z=43$) and Promethium ($Z=61$) have no stable isotopes.
\end{experience}

\begin{propriete}[Binding Energy and Mass Defect]
The mass of a nucleus is \textbf{less} than the sum of its free constituent nucleons. The difference is the \textbf{mass defect} $\Delta m$.
\[
\Delta m = \left[ Z m_p + (A-Z) m_n \right] - M_{\text{nucleus}}
\]
Using atomic masses (includes electrons, cancels out if $Z$ conserved):
\[
\Delta m = \left[ Z M(\ensuremath{\mathrm{^{1}_{1}H}}) + (A-Z) m_n \right] - M(\ensuremath{\mathrm{^{A}_{Z}X}})
\]
The \textbf{Binding Energy} $E_b$ is the energy equivalent (Einstein):
\[
E_b = \Delta m \, c^2
\]
\textbf{Binding Energy per Nucleon} $E_b/A$ measures stability. Maximum $\approx \SI{8.8}{MeV}$ at $\ensuremath{\mathrm{^{56}_{26}Fe}}$ (most stable nucleus).
\end{propriete}

\begin{aretenir}[Binding Energy Curve Features]
\begin{itemize}
\item Rises steeply for light nuclei $\to$ \textbf{Fusion} (combining light nuclei) releases energy (products more tightly bound).
\item Broad maximum around $A=56$ ($\ce{Fe}$, $\ce{Ni}$).
\item Falls slowly for heavy nuclei $\to$ \textbf{Fission} (splitting heavy nuclei) releases energy (fragments more tightly bound).
\item $\ensuremath{\mathrm{^{4}_{2}He}}$ ($\alpha$ particle) has exceptionally high $E_b/A \approx \SI{7.07}{MeV}$ $\to$ explains $\alpha$ decay prevalence in heavy nuclei.
\end{itemize}
\end{aretenir}

\begin{exoresolu}[Binding Energy of Helium-4 and Alpha Decay Energy]
Given atomic masses: $\ensuremath{\mathrm{^{4}_{2}He}} = 4.002603\ \text{u}$, $\ensuremath{\mathrm{^{1}_{1}H}} = 1.007825\ \text{u}$, $m_n = 1.008665\ \text{u}$, $\ensuremath{\mathrm{^{238}_{92}U}} = 238.050788\ \text{u}$, $\ensuremath{\mathrm{^{234}_{90}Th}} = 234.043601\ \text{u}$.
\begin{enumerate}
\item Calculate the binding energy per nucleon of $\ensuremath{\mathrm{^{4}_{2}He}}$.
\item Calculate the $Q$-value for the $\alpha$-decay: $\ce{^{238}_{92}U -> ^{234}_{90}Th + ^{4}_{2}He}$.
\item Determine the kinetic energy of the $\alpha$ particle ($T_\alpha$) assuming the Uranium nucleus was at rest.
\end{enumerate}
\tcblower
\textbf{Detailed Solution}

\textbf{1. Binding Energy of $\ensuremath{\mathrm{^{4}_{2}He}}$:}
Constituents: $2$ protons + $2$ neutrons $\equiv 2$ Hydrogen atoms + $2$ neutrons.
Mass of constituents $= 2(1.007825) + 2(1.008665) = 2.015650 + 2.017330 = 4.032980\ \text{u}$.
Mass of $\ce{^{4}He}$ atom $= 4.002603\ \text{u}$.
Mass defect $\Delta m = 4.032980 - 4.002603 = 0.030377\ \text{u}$.
$E_b = 0.030377 \times 931.5\ \text{MeV} = 28.30\ \text{MeV}$.
$E_b/A = 28.30 / 4 = \SI{7.07}{MeV/nucleon}$.
(Note: Very high for light nucleus, explains $\alpha$ stability).

\textbf{2. $Q$-value for $\ce{^{238}U}$ $\alpha$-decay:}
$Q = [M(\ce{^{238}U}) - M(\ce{^{234}Th}) - M(\ce{^{4}He})] c^2$.
$\Delta m = 238.050788 - 234.043601 - 4.002603 = 0.004584\ \text{u}$.
$Q = 0.004584 \times 931.5 = \SI{4.270}{MeV}$.
(Total kinetic energy shared by $\alpha$ and Th).

\textbf{3. Kinetic Energy of $\alpha$ particle ($T_\alpha$):}
Conservation of momentum: $p_\alpha = p_{\text{Th}}$ (magnitudes equal, opposite directions).
$T = p^2 / 2m \Rightarrow T \propto 1/m$.
$T_\alpha / T_{\text{Th}} = m_{\text{Th}} / m_\alpha \approx 234 / 4 = 58.5$.
$Q = T_\alpha + T_{\text{Th}} = T_\alpha (1 + m_\alpha/m_{\text{Th}}) = T_\alpha (1 + 4/234) = T_\alpha (238/234)$.
$T_\alpha = Q \times \frac{234}{238} = 4.270 \times 0.9832 = \SI{4.198}{MeV}$.
$T_{\text{Th}} = Q - T_\alpha = \SI{0.072}{MeV}$.
\textbf{Check:} Typical $\alpha$ energy for $\ce{^{238}U}$ is $\SI{4.20}{MeV}$. Matches.
\end{exoresolu}

\coursec{Types of Radiation, Properties and Applications}

\begin{propriete}[Modes of Radioactive Decay]
An unstable nucleus decays to reach a more stable configuration. Conservation laws: $A$, $Z$, energy-momentum, charge, lepton number.
\begin{itemize}
\item \textbf{$\alpha$ decay:} $\ce{^{A}_{Z}X -> ^{A-4}_{Z-2}Y + ^{4}_{2}He}$. Occurs in very heavy nuclei ($Z > 82$). Discrete energy spectrum.
\item \textbf{$\beta^-$ decay:} $\ensuremath{\mathrm{^{A}_{Z}X \rightarrow  ^{A}_{Z+1}Y + ^{0}_{-1}e + \bar{\nu}_e}}$. Neutron-rich. $n \to p + e^- + \bar{\nu}_e$. Continuous energy spectrum (shared with neutrino).
\item \textbf{$\beta^+$ decay:} $\ensuremath{\mathrm{^{A}_{Z}X \rightarrow  ^{A}_{Z-1}Y + ^{0}_{+1}e + \nu_e}}$. Proton-rich. $p \to n + e^+ + \nu_e$. Requires $Q > 2m_e c^2 = 1.022\ \text{MeV}$.
\item \textbf{Electron Capture (EC):} $\ensuremath{\mathrm{^{A}_{Z}X + ^{0}_{-1}e \rightarrow  ^{A}_{Z-1}Y + \nu_e}}$. Inner shell electron captured. Alternative to $\beta^+$ if $Q < 1.022\ \text{MeV}$.
\item \textbf{$\gamma$ decay:} $\ensuremath{\mathrm{^{A}_{Z}X^* \rightarrow  ^{A}_{Z}X + \gamma}}$. Excited nucleus de-excites. No change in $A$ or $Z$. Often follows $\alpha$ or $\beta$ decay leaving daughter in metastable state (e.g., $\ce{^{99m}Tc}$).
\end{itemize}
\end{propriete}

\begin{aretenir}[Properties of Radiations]
\begin{center}
\begin{tabularx}{\linewidth}{|l|X|X|X|}\hline
\rowcolor{popblueL}
\textbf{Property} & \textbf{$\alpha$ particle} & \textbf{$\beta^-$ particle} & \textbf{$\gamma$ ray} \\ \hline
\textbf{Nature} & $\ensuremath{\mathrm{^{4}_{2}He}}$ nucleus & Electron ($\ensuremath{\mathrm{^{0}_{-1}e}}$) & High-energy photon \\ \hline
\textbf{Charge} & $+2e$ & $-e$ & $0$ \\ \hline
\textbf{Mass} & $\approx 4\ \text{u}$ ($6.64 \times 10^{-27}\ \text{kg}$) & $m_e = 9.11 \times 10^{-31}\ \text{kg}$ & $0$ (energy $E=h\nu$) \\ \hline
\textbf{Speed} & $\sim 0.05c - 0.1c$ & Up to $\sim 0.99c$ (relativistic) & $c$ \\ \hline
\textbf{Ionising Power} & Very High (dense track) & Medium & Low (indirect) \\ \hline
\textbf{Range in Air} & Few cm ($\approx 5-8\ \text{cm}$) & $\sim 1\ \text{m} / \text{MeV}$ & Very large (km) \\ \hline
\textbf{Penetration} & Stopped by paper / skin & Stopped by few mm Al / perspex & Requires thick Pb / concrete \\ \hline
\textbf{Deflection in E-field} & Towards negative plate (small) & Towards positive plate (large) & None \\ \hline
\textbf{Deflection in B-field} & Perpendicular to $\vec{v}, \vec{B}$ (LH rule) & Opposite direction to $\alpha$ (tighter) & None \\ \hline
\end{tabularx}
\end{center}
\end{aretenir}

\begin{methode}[Identifying Radiation using Electric and Magnetic Fields]
\begin{enumerate}
\item \textbf{Setup:} Source $\to$ narrow beam (collimator) $\to$ uniform field region $\to$ detector (photographic plate, GM tube, or fluorescent screen). Vacuum or low pressure for $\alpha/\beta$.
\item \textbf{Electric Field ($\vec{E}$):} Force $\vec{F} = q\vec{E}$. Parabolic path if $v$ constant.
\begin{itemize}
\item $\alpha$ ($q=+2e$): Deflected \textbf{towards negative plate}, small curvature (large mass $m \approx 4u$).
\item $\beta^-$ ($q=-e$): Deflected \textbf{towards positive plate}, large curvature (small mass $m_e$).
\item $\gamma$ ($q=0$): \textbf{No deflection}, straight line.
\end{itemize}
\item \textbf{Magnetic Field ($\vec{B}$, into page $\otimes$):} Force $\vec{F} = q(\vec{v} \wedge \vec{B})$. Circular path $r = mv/(qB)$. Use Fleming's Left Hand Rule (Motor Rule) for positive charge; reverse for negative.
\begin{itemize}
\item $\alpha$: Deflected \textbf{upwards} (if beam right, B into page), radius $r_\alpha = \frac{m_\alpha v}{2eB}$.
\item $\beta^-$: Deflected \textbf{downwards}, radius $r_\beta = \frac{m_e v}{eB} \ll r_\alpha$ (for same $v$). Relativistic correction often needed for $\beta$: $r = \frac{\gamma m v}{qB}$.
\item $\gamma$: No deflection.
\end{itemize}
\item \textbf{Key distinction:} $\alpha$ and $\beta$ deflect in \textbf{opposite directions} in both fields. $\gamma$ never deflects. $\beta$ deflection $\gg$ $\alpha$ deflection for same energy (or same velocity).
\end{enumerate}
\end{methode}

\begin{exoresolu}[Deflection in a Magnetic Field]
A mixed beam of $\alpha$ particles and $\beta^-$ particles, both with kinetic energy \SI{5.00}{MeV}, enters a uniform magnetic field of flux density \SI{0.200}{T} perpendicular to their velocity.
\begin{enumerate}
\item Calculate the velocity of the $\alpha$ particle (mass $m_\alpha = 6.64 \times 10^{-27}\ \text{kg}$) and the $\beta$ particle (mass $m_e = 9.11 \times 10^{-31}\ \text{kg}$). Note: Check if relativistic correction is needed for the $\beta$ particle.
\item Determine the radius of curvature for each particle in the magnetic field.
\item Sketch the paths on a diagram (beam entering horizontally, $\vec{B}$ into page).
\end{enumerate}
\tcblower
\textbf{Detailed Solution}

\textbf{1. Velocities:}
Kinetic energy $K = \SI{5.00}{MeV} = 5.00 \times 10^6 \times 1.60 \times 10^{-19} = \SI{8.00e-13}{J}$.

\textit{$\alpha$ particle (Non-relativistic check):}
$K = \frac{1}{2} m_\alpha v_\alpha^2 \Rightarrow v_\alpha = \sqrt{\frac{2K}{m_\alpha}} = \sqrt{\frac{2 \times 8.00 \times 10^{-13}}{6.64 \times 10^{-27}}} = \sqrt{2.41 \times 10^{14}} = \SI{1.55e7}{m.s^{-1}}$.
$v_\alpha / c = 1.55 \times 10^7 / 3.00 \times 10^8 = 0.052$ ($5.2\% c$). Non-relativistic approximation valid.

\textit{$\beta$ particle (Relativistic check):}
Classical $v = \sqrt{2K/m_e} = \sqrt{2 \times 8.00 \times 10^{-13} / 9.11 \times 10^{-31}} = \SI{1.33e9}{m.s^{-1}} > c$. \textbf{Impossible. Must use relativity.}
Total Energy $E = K + m_0 c^2 = 5.00\ \text{MeV} + 0.511\ \text{MeV} = 5.511\ \text{MeV}$.
$E = \gamma m_0 c^2 \Rightarrow \gamma = \frac{5.511}{0.511} = 10.78$.
$\gamma = \frac{1}{\sqrt{1 - v^2/c^2}} \Rightarrow v = c \sqrt{1 - 1/\gamma^2} = 3.00 \times 10^8 \sqrt{1 - 1/116.2} = \SI{2.987e8}{m.s^{-1}} \approx 0.996c$.

\textbf{2. Radius of curvature $r = \frac{p}{qB} = \frac{\gamma m v}{qB}$ (Relativistic momentum).}

\textit{$\alpha$ particle:} $q_\alpha = +2e = 3.20 \times 10^{-19}\ \text{C}$.
$p_\alpha = m_\alpha v_\alpha = 6.64 \times 10^{-27} \times 1.55 \times 10^7 = 1.03 \times 10^{-19}\ \text{kg m/s}$.
$r_\alpha = \frac{1.03 \times 10^{-19}}{3.20 \times 10^{-19} \times 0.200} = \SI{1.61}{m}$.

\textit{$\beta$ particle:} $q_\beta = -e = 1.60 \times 10^{-19}\ \text{C}$.
Relativistic momentum $p_\beta = \frac{1}{c}\sqrt{E^2 - (m_0 c^2)^2}$.
$E = 5.511\ \text{MeV} = 8.82 \times 10^{-13}\ \text{J}$. $m_0 c^2 = 0.511\ \text{MeV} = 8.18 \times 10^{-14}\ \text{J}$.
$p_\beta c = \sqrt{(8.82 \times 10^{-13})^2 - (8.18 \times 10^{-14})^2} \approx 8.78 \times 10^{-13}\ \text{J}$.
$p_\beta = 2.93 \times 10^{-21}\ \text{kg m/s}$.
$r_\beta = \frac{2.93 \times 10^{-21}}{1.60 \times 10^{-19} \times 0.200} = \SI{0.0915}{m} = \SI{9.15}{cm}$.

\textbf{3. Sketch Description:}
\begin{popfigure}
\begin{tikzpicture}[scale=0.8]
\draw[->, thick] (-3,0) -- (4,0) node[right] {Beam direction};
\draw[popblue, thick] (-3,0) .. controls (-1,1.5) and (1,1.5) .. (3,1.6) node[right, black] {$\alpha$ ($r \approx 1.6$ m)};
\draw[poporange, thick] (-3,0) .. controls (-1,-0.5) and (0,-0.5) .. (1,-0.5) node[right, black] {$\beta$ ($r \approx 9$ cm)};
\draw[popgreen, thick, dashed] (-3,0) -- (4,0) node[right, black] {$\gamma$ (undeflected)};
% B field symbol (dots for into page, skip y=0 to avoid overlapping beam line)
\foreach \x in {-2,-1,0,1,2,3} {
  \foreach \y in {-1.5,-1,-0.5,0.5,1,1.5} {
    \fill[popmuted] (\x,\y) circle (0.05);
  }
}
\node at (0, -2) {$\vec{B}$ into page ($\otimes$)};
\end{tikzpicture}
\legende{Paths in magnetic field (qualitative sketch). $\alpha$ curves slightly up (large $r$). $\beta$ curves sharply down (small $r$). $\gamma$ straight.}
\end{popfigure}
\textbf{Note:} $\alpha$ deflects up (Fleming LH rule: current = velocity, field = into page $\to$ force up). $\beta$ has negative charge $\to$ force down.
\end{exoresolu}

\begin{savaistu}[Applications in Cameroon Context]
\begin{itemize}
\item \textbf{Radiotherapy:} $\ce{^{60}Co}$ (gamma, $\SI{1.17}{MeV}$ \& $\SI{1.33}{MeV}$) units in Yaoundé and Douala hospitals for cancer treatment. Being replaced by Linear Accelerators (Linacs).
\item \textbf{Industrial Gauges:} $\beta$ sources (e.g., $\ce{^{90}Sr}$) for paper/plastic thickness control in factories (e.g., paper mill in Edéa). $\gamma$ sources ($\ce{^{137}Cs}$, $\ce{^{60}Co}$) for level gauges in oil tanks (SONARA refinery, Limbe) and density gauges.
\item \textbf{Tracers:} $\ce{^{131}I}$ (thyroid diagnosis), $\ce{^{99m}Tc}$ (SPECT imaging - see worked example later). $\ce{^{3}H}$, $\ce{^{14}C}$ in hydrology (tracing groundwater in Sahel regions).
\item \textbf{Smoke Detectors:} $\ce{^{241}Am}$ ($\alpha$ source) ionises air; smoke reduces current $\to$ alarm.
\item \textbf{Sterilisation:} Gamma irradiation ($\ce{^{60}Co}$) of medical equipment, food preservation (potatoes, spices).
\end{itemize}
\end{savaistu}

\coursec{Mass Defect and Nuclear Energy}

\begin{propriete}[Radioactive Decay Law]
The decay of a radioactive sample is a random process at the level of single nuclei, but statistically predictable for large numbers.
\begin{itemize}
\item \textbf{Decay Constant $\lambda$:} Probability per unit time that a given nucleus will decay. Unit: $\si{s^{-1}}$.
\item \textbf{Activity $A$:} Rate of decay $A = -\frac{dN}{dt} = \lambda N$. Unit: Becquerel (\si{Bq} = \si{s^{-1}}). Old unit: Curie (\si{Ci} = \SI{3.7e10}{Bq}).
\item \textbf{Differential Equation:} $\frac{dN}{dt} = -\lambda N$.
\item \textbf{Solution:} $N(t) = N_0 e^{-\lambda t}$ and $A(t) = A_0 e^{-\lambda t}$.
\item \textbf{Half-Life $T_{1/2}$:} Time for half the nuclei to decay. $N(T_{1/2}) = N_0/2$.
\[
T_{1/2} = \frac{\ln 2}{\lambda} \approx \frac{0.693}{\lambda}
\]
\item \textbf{Fraction remaining after time $t$:} $\frac{N}{N_0} = e^{-\lambda t} = \left(\frac{1}{2}\right)^{t/T_{1/2}}$.
\end{itemize}
\end{propriete}

\begin{methode}[Radioactive Decay Law and Half-Life Calculations]
\begin{enumerate}
\item \textbf{Identify the knowns:} Initial number of nuclei $N_0$ (or initial activity $A_0$), half-life $T_{1/2}$ (or decay constant $\lambda$), and time elapsed $t$.
\item \textbf{Convert half-life to decay constant:} $\lambda = \dfrac{\ln 2}{T_{1/2}}$. Ensure consistent time units (seconds, years, etc.).
\item \textbf{Apply the decay law:} Number of remaining nuclei $N = N_0 e^{-\lambda t}$. Activity $A = A_0 e^{-\lambda t} = \lambda N$.
\item \textbf{Calculate fraction remaining/decayed:} Fraction remaining $= e^{-\lambda t} = \left(\frac{1}{2}\right)^{t/T_{1/2}}$. Fraction decayed $= 1 - e^{-\lambda t}$.
\item \textbf{Check units:} Activity in Becquerel (\si{Bq} = \si{s^{-1}}), time in seconds (or consistent unit), $\lambda$ in \si{s^{-1}}.
\item \textbf{Reflex:} For integer numbers of half-lives ($t = n T_{1/2}$), use $N = N_0 / 2^n$ directly for speed.
\end{enumerate}
\end{methode}

\begin{exoresolu}[Activity of a Cobalt-60 Source]
A radiotherapy unit in a Yaoundé hospital uses a $\ensuremath{\mathrm{^{60}_{27}Co}}$ source with an initial activity of \SI{3.7e12}{Bq} (\SI{100}{Ci}). The half-life of $\ce{^{60}Co}$ is \SI{5.27}{years}.
\begin{enumerate}
\item Calculate the decay constant $\lambda$ in \si{s^{-1}}.
\item Determine the activity after \SI{3.00}{years}.
\item After how long will the activity fall to \SI{1.00e11}{Bq}?
\end{enumerate}
\tcblower
\textbf{Detailed Solution}

\textbf{1. Decay constant $\lambda$:}
Convert half-life to seconds:
$T_{1/2} = 5.27\ \text{years} \times 365.25\ \frac{\text{days}}{\text{year}} \times 24\ \frac{\text{h}}{\text{day}} \times 3600\ \frac{\text{s}}{\text{h}} = 1.663 \times 10^8\ \text{s}$.
\[
\lambda = \frac{\ln 2}{T_{1/2}} = \frac{0.693}{1.663 \times 10^8\ \text{s}} = \SI{4.17e-9}{s^{-1}}
\]

\textbf{2. Activity after $t = 3.00$ years:}
Convert $t$ to seconds: $t = 3.00 \times 3.15576 \times 10^7\ \text{s} = 9.467 \times 10^7\ \text{s}$.
Use $A = A_0 e^{-\lambda t}$:
\[
A = 3.7 \times 10^{12} \times \exp\left(-4.17 \times 10^{-9} \times 9.467 \times 10^7\right)
\]
\[
\lambda t = 0.395 \quad \Rightarrow \quad e^{-0.395} = 0.674
\]
\[
A = 3.7 \times 10^{12} \times 0.674 = \SI{2.49e12}{Bq}
\]
\textit{Alternative (using years directly):} $\lambda_y = \ln 2 / 5.27 = 0.1315\ \text{year}^{-1}$. $A = 3.7 \times 10^{12} \times e^{-0.1315 \times 3} = \SI{2.49e12}{Bq}$.

\textbf{3. Time for $A = \SI{1.00e11}{Bq}$:}
\[
\frac{A}{A_0} = e^{-\lambda t} \quad \Rightarrow \quad \frac{1.00 \times 10^{11}}{3.7 \times 10^{12}} = 0.0270 = e^{-\lambda t}
\]
\[
-\lambda t = \ln(0.0270) = -3.612 \quad \Rightarrow \quad t = \frac{3.612}{\lambda}
\]
Using $\lambda = 0.1315\ \text{year}^{-1}$:
\[
t = \frac{3.612}{0.1315} = \SI{27.5}{years}
\]
\textbf{Answer:} $\lambda = \SI{4.17e-9}{s^{-1}}$; $A(3\ \text{yr}) = \SI{2.49e12}{Bq}$; $t = \SI{27.5}{years}$.
\end{exoresolu}

\begin{methode}[Balancing Nuclear Equations and Decay Series]
\begin{enumerate}
\item \textbf{Conservation Laws:} Total nucleon number $A$ (mass number) and total proton number $Z$ (atomic number) are conserved on both sides of the arrow.
\item \textbf{Identify the radiation:}
\begin{itemize}
\item $\alpha$-decay: $A \to A-4$, $Z \to Z-2$ (emits $\ensuremath{\mathrm{^{4}_{2}He}}$).
\item $\beta^-$-decay: $A$ constant, $Z \to Z+1$ (emits $\ensuremath{\mathrm{^{0}_{-1}e}} + \bar{\nu}_e$).
\item $\beta^+$-decay / EC: $A$ constant, $Z \to Z-1$ (emits $\ensuremath{\mathrm{^{0}_{+1}e}} + \nu_e$ or captures $e^-$).
\item $\gamma$-decay: $A$ and $Z$ unchanged (energy release only).
\end{itemize}
\item \textbf{Step-by-step for series:} Write the parent nucleus. Apply decay 1 $\to$ daughter 1. Apply decay 2 $\to$ daughter 2, etc. Verify final stable isotope.
\item \textbf{Notation:} Use $\ensuremath{\mathrm{^{A}_{Z}X}}$ format. For electrons/positrons: $\ensuremath{\mathrm{^{0}_{-1}e}}$, $\ensuremath{\mathrm{^{0}_{+1}e}}$. For neutrinos: $\ensuremath{\mathrm{\nu_e}}$, $\ensuremath{\mathrm{\bar{\nu}_e}}$.
\end{enumerate}
\end{methode}

\begin{exoresolu}[Natural Radioactive Series: Uranium-238]
The $\ensuremath{\mathrm{^{238}_{92}U}}$ series ends at stable $\ensuremath{\mathrm{^{206}_{82}Pb}}$. The sequence of decays is: $\alpha, \beta, \beta, \alpha, \alpha, \alpha, \alpha, \beta, \beta, \alpha, \beta, \alpha, \beta$.
\begin{enumerate}
\item Write the full balanced nuclear equation for the first three decays.
\item Determine the total number of $\alpha$ and $\beta$ particles emitted in the complete series.
\item Identify the noble gas isotope produced in this series (Radon) and its half-life (\SI{3.82}{days}). Why is it a radiological hazard in homes?
\end{enumerate}
\tcblower
\textbf{Detailed Solution}

\textbf{1. First three decays:}
Start: $\ensuremath{\mathrm{^{238}_{92}U}}$
\begin{itemize}
\item Decay 1 ($\alpha$): $A=234, Z=90$ $\to$ $\ensuremath{\mathrm{^{234}_{90}Th}}$
$\ce{^{238}_{92}U -> ^{234}_{90}Th + ^{4}_{2}He}$
\item Decay 2 ($\beta^-$): $A=234, Z=91$ $\to$ $\ensuremath{\mathrm{^{234}_{91}Pa}}$
$\ensuremath{\mathrm{^{234}_{90}Th \rightarrow  ^{234}_{91}Pa + ^{0}_{-1}e + \bar{\nu}_e}}$
\item Decay 3 ($\beta^-$): $A=234, Z=92$ $\to$ $\ensuremath{\mathrm{^{234}_{92}U}}$
$\ensuremath{\mathrm{^{234}_{91}Pa \rightarrow  ^{234}_{92}U + ^{0}_{-1}e + \bar{\nu}_e}}$
\end{itemize}

\textbf{2. Total $\alpha$ and $\beta$ count:}
Let $n_\alpha$ = number of $\alpha$, $n_\beta$ = number of $\beta^-$.
Change in $A$: $\Delta A = 238 - 206 = 32$. Only $\alpha$ changes $A$ (by $-4$).
$4 n_\alpha = 32 \Rightarrow n_\alpha = 8$.
Change in $Z$: $\Delta Z = 92 - 82 = 10$. $\alpha$ decreases $Z$ by 2, $\beta^-$ increases $Z$ by 1.
$-2 n_\alpha + n_\beta = -10 \Rightarrow -2(8) + n_\beta = -10 \Rightarrow n_\beta = 6$.
\textbf{Check:} The given sequence has 8 $\alpha$ and 6 $\beta$ (count them: $\alpha\beta\beta\alpha\alpha\alpha\alpha\beta\beta\alpha\beta\alpha\beta$). Matches.

\textbf{3. Radon isotope:}
Look at the series steps. After 6 decays ($\alpha\beta\beta\alpha\alpha\alpha$), we reach $A=222, Z=86$ (Radon).
$\ensuremath{\mathrm{^{226}_{88}Ra \rightarrow [\alpha] ^{222}_{86}Rn}}$ (Radon-222).
Half-life $T_{1/2} = \SI{3.82}{days}$.
\textbf{Hazard:} $\ce{^{222}Rn}$ is a noble gas, chemically inert. It diffuses from soil/rocks (granite areas in West Cameroon, Adamawa) into basements and poorly ventilated rooms. Being a gas, it is inhaled. Its short-lived solid daughters ($\ce{^{218}Po}$, $\ce{^{214}Pb}$, $\ce{^{214}Bi}$, $\ce{^{214}Po}$) attach to dust aerosols, deposit in lungs, and emit high-energy $\alpha$ particles, causing significant radiation dose to bronchial epithelium $\to$ lung cancer risk. It is the second leading cause of lung cancer after smoking.
\end{exoresolu}

\begin{methode}[Radioactive Dating using Carbon-14]
\begin{enumerate}
\item \textbf{Principle:} Living organisms maintain equilibrium $\ce{^{14}C}/\ce{^{12}C}$ ratio ($\approx 1.3 \times 10^{-12}$) via $\ce{^{14}N(n,p)^{14}C}$ in atmosphere and photosynthesis/feeding. At death, intake stops. $\ce{^{14}C}$ decays ($\beta^-$, $T_{1/2} = 5730\ \text{yr}$). Ratio decreases.
\item \textbf{Formula:} $\frac{A}{A_0} = \frac{N}{N_0} = e^{-\lambda t} = \left(\frac{1}{2}\right)^{t/T_{1/2}}$.
\item \textbf{Solve for age $t$:} $t = -\frac{1}{\lambda} \ln\left(\frac{A}{A_0}\right) = -T_{1/2} \frac{\ln(A/A_0)}{\ln 2}$.
\item \textbf{Assumptions/Limits:} Constant atmospheric $\ce{^{14}C}$ production (calibration curves needed for precision), closed system (no contamination), range $\approx 100 - 50,000$ years (beyond: activity too low; below: not enough decay).
\item \textbf{Units:} $T_{1/2} = 5730\ \text{yr}$ (Libby half-life) or $5700\ \text{yr}$ (Cambridge). Use value given in question.
\end{enumerate}
\end{methode}

\begin{exoresolu}[Age of a Wooden Artefact]
A wooden mask from a Cameroonian museum collection shows a $\ce{^{14}C}$ activity of \SI{0.12}{Bq per gram of carbon}. The activity of living wood is \SI{0.25}{Bq/g}. The half-life of $\ce{^{14}C}$ is taken as \SI{5730}{years}.
\begin{enumerate}
\item Calculate the age of the mask.
\item The measured activity has an uncertainty of $\pm \SI{0.01}{Bq/g}$. Determine the uncertainty in the age.
\item Why can't this method date a dinosaur bone (\SI{65}{million years old})?
\end{enumerate}
\tcblower
\textbf{Detailed Solution}

\textbf{1. Age calculation:}
$\frac{A}{A_0} = \frac{0.12}{0.25} = 0.48$.
$\lambda = \frac{\ln 2}{5730} = 1.2097 \times 10^{-4}\ \text{yr}^{-1}$.
$t = -\frac{1}{\lambda} \ln(0.48) = -\frac{1}{1.2097 \times 10^{-4}} \times (-0.7340) = \SI{6068}{years}$.
Using half-life formula: $t = -5730 \times \frac{\ln 0.48}{\ln 2} = 5730 \times 1.058 = \SI{6062}{years}$.
\textbf{Age} $\approx \SI{6070}{years}$ (Neolithic period).

\textbf{2. Uncertainty propagation:}
$A = 0.12 \pm 0.01 \Rightarrow A_{\text{min}} = 0.11, A_{\text{max}} = 0.13$.
$t_{\text{max}} = -5730 \frac{\ln(0.11/0.25)}{\ln 2} = -5730 \frac{\ln 0.44}{\ln 2} = 6600\ \text{yr}$.
$t_{\text{min}} = -5730 \frac{\ln(0.13/0.25)}{\ln 2} = -5730 \frac{\ln 0.52}{\ln 2} = 5580\ \text{yr}$.
Uncertainty $\Delta t \approx \frac{6600 - 5580}{2} = \SI{510}{years}$.
Or via derivative: $\Delta t = \frac{T_{1/2}}{\ln 2} \frac{\Delta A}{A} = \frac{5730}{0.693} \times \frac{0.01}{0.12} = 8268 \times 0.0833 = \SI{689}{years}$.
(Asymmetric errors due to log; range method better). \textbf{Age} = $\SI{6070 \pm 510}{years}$.

\textbf{3. Dinosaur bone limit:}
$t = 65 \times 10^6\ \text{yr}$. Number of half-lives $n = 65 \times 10^6 / 5730 \approx 11340$.
Fraction remaining $= (1/2)^{11340} \approx 10^{-3414}$. Effectively zero.
Activity would be immeasurably small (far below background radiation). Carbon-14 dating is limited to $\approx 8-10$ half-lives ($\sim 50,000$ years). For millions of years, use $\ce{^{40}K}/\ce{^{40}Ar}$ ($T_{1/2}=1.25\times 10^9\ \text{yr}$) or $\ce{^{238}U}/\ce{^{206}Pb}$ on surrounding volcanic rock (zircon crystals).
\end{exoresolu}

\coursec{Nuclear Fission and Nuclear Fusion}

\begin{propriete}[Nuclear Fission]
\textbf{Discovery:} Hahn, Strassmann, Meitner, Frisch (1938-39). Neutron-induced fission of $\ce{^{235}U}$.
\textbf{Typical Reaction:}
\[
\ensuremath{\mathrm{^{235}_{92}U + ^{1}_{0}n \rightarrow  ^{A_1}_{Z_1}X + ^{A_2}_{Z_2}Y + k\,^{1}_{0}n + Q}}
\]
\textbf{Conservation:} $A_{\text{U}} + 1 = A_1 + A_2 + k$; $Z_{\text{U}} = Z_1 + Z_2$.
Fragments $X, Y$ are neutron-rich $\to$ undergo $\beta^-$ decay chains (fission products).
\textbf{Energy $Q$:} $Q = [M(\ce{^{235}U}) + m_n - M(X) - M(Y) - k m_n] c^2$.
Or using Binding Energy: $Q = E_b(X) + E_b(Y) - E_b(\ce{^{235}U})$. (Neutron has $E_b=0$).
\textbf{Typical values:} $Q \approx \SI{200}{MeV}$ per fission.
Distribution: Kinetic energy of fragments ($\sim 165\ \text{MeV}$), prompt $\gamma$ ($\sim 7\ \text{MeV}$), kinetic energy of neutrons ($\sim 5\ \text{MeV}$), $\beta/\gamma$ from fission products ($\sim 23\ \text{MeV}$), neutrinos ($\sim 10\ \text{MeV}$, lost).
\textbf{Chain Reaction:} Requires $k \ge 1$ neutron per fission to sustain. Critical mass, moderator (slow neutrons increase $\sigma_f$ for $\ce{^{235}U}$), control rods (absorb neutrons: B, Cd, Hf).
\end{propriete}

\begin{methode}[Energy Released in Nuclear Fission]
\begin{enumerate}
\item \textbf{Write balanced equation:} Check $A$ and $Z$ conservation. Identify fragments and number of neutrons $k$.
\item \textbf{Calculate $Q$-value:} Use atomic masses. $Q = (\sum m_{\text{initial}} - \sum m_{\text{final}}) c^2$. $1\ \text{u} = 931.5\ \text{MeV}$.
\item \textbf{Power calculation:} $P = \frac{\text{Energy}}{\text{time}} = R \times Q$, where $R$ is fission rate (fissions/s). Mass consumption rate $\dot{m} = R \times m_{\text{nucleus}}$.
\end{enumerate}
\end{methode}

\begin{exoresolu}[Fission of Uranium-235]
A thermal neutron induces fission in $\ensuremath{\mathrm{^{235}_{92}U}}$ producing $\ensuremath{\mathrm{^{141}_{56}Ba}}$, $\ensuremath{\mathrm{^{92}_{36}Kr}}$ and 3 neutrons.
Atomic masses: $\ce{^{235}U} = 235.043930\ \text{u}$, $\ce{^{141}Ba} = 140.914411\ \text{u}$, $\ce{^{92}Kr} = 91.926156\ \text{u}$, neutron = 1.008665 u.
\begin{enumerate}
\item Write the balanced nuclear equation.
\item Calculate the total energy released $Q$ in MeV.
\item If a reactor core produces \SI{3000}{MW} of thermal power, calculate the mass of $\ce{^{235}U}$ consumed (fissioned) per day. (Assume $200\ \text{MeV}$/fission for usable energy).
\end{enumerate}
\tcblower
\textbf{Detailed Solution}

\textbf{1. Balanced Equation:}
$\ensuremath{\mathrm{^{235}_{92}U + ^{1}_{0}n \rightarrow  ^{141}_{56}Ba + ^{92}_{36}Kr + 3\,^{1}_{0}n}}$
Check: $A: 235+1 = 141+92+3 = 236$. $Z: 92 = 56+36 = 92$. OK.

\textbf{2. Energy Released $Q$:}
Initial mass $m_i = M(\ce{^{235}U}) + m_n = 235.043930 + 1.008665 = 236.052595\ \text{u}$.
Final mass $m_f = M(\ce{^{141}Ba}) + M(\ce{^{92}Kr}) + 3 m_n = 140.914411 + 91.926156 + 3(1.008665) = 232.840567 + 3.025995 = 235.866562\ \text{u}$.
Mass defect $\Delta m = m_i - m_f = 236.052595 - 235.866562 = 0.186033\ \text{u}$.
$Q = 0.186033 \times 931.5 = \SI{173.3}{MeV}$.
(Note: This specific channel yields $\sim 173\ \text{MeV}$. Average over all channels is $\sim 200\ \text{MeV}$. The difference is in excitation energy of fragments / prompt gammas).

\textbf{3. Mass consumption rate:}
Power $P = \SI{3000}{MW} = 3.00 \times 10^9\ \text{J/s}$.
Energy per fission $E_f = 200\ \text{MeV} = 200 \times 1.60 \times 10^{-13} = 3.20 \times 10^{-11}\ \text{J}$.
Fissions per second $R = P / E_f = 3.00 \times 10^9 / 3.20 \times 10^{-11} = 9.375 \times 10^{19}\ \text{s}^{-1}$.
Seconds per day $= 86400\ \text{s}$.
Total fissions/day $= 9.375 \times 10^{19} \times 86400 = 8.10 \times 10^{24}$.
Mass of one $\ce{^{235}U}$ atom $= 235 \times 1.6605 \times 10^{-27} = 3.90 \times 10^{-25}\ \text{kg}$.
Mass consumed $= 8.10 \times 10^{24} \times 3.90 \times 10^{-25} = \SI{3.16}{kg/day}$.
\textbf{Result:} $\approx \SI{3.2}{kg}$ of $\ce{^{235}U}$ fissioned per day. (Total fuel load much higher due to enrichment, non-fission capture, etc.).
\end{exoresolu}

\begin{propriete}[Nuclear Fusion]
\textbf{Principle:} Light nuclei combine to form a heavier, more tightly bound nucleus. Energy released because $E_b/A$ increases.
\textbf{Conditions (Lawson Criterion):}
\begin{itemize}
\item High Temperature ($T > 10^7\ \text{K}$, $\sim 1\ \text{keV}$) for kinetic energy to overcome Coulomb barrier ($V_C \approx \frac{Z_1 Z_2 e^2}{4\pi\epsilon_0 r}$). Quantum tunnelling essential.
\item High Density/Pressure ($n$) for collision frequency.
\item Confinement Time ($\tau_E$): Energy must be retained long enough. Product $n\tau_E$ must exceed threshold (Lawson criterion).
\end{itemize}
\textbf{Main Reactions:}
\begin{itemize}
\item \textbf{p-p Chain (Stars like Sun):} $4\ \ensuremath{\mathrm{^{1}_{1}H}} \to \ensuremath{\mathrm{^{4}_{2}He}} + 2\ \ensuremath{\mathrm{^{0}_{+1}e}} + 2\nu_e + \gamma$ ($Q \approx 26.7\ \text{MeV}$). Slow (weak interaction).
\item \textbf{D-T (Tokamaks/ITER):} $\ce{^{2}_{1}H + ^{3}_{1}H -> ^{4}_{2}He + ^{1}_{0}n}$ ($Q = 17.6\ \text{MeV}$). Largest cross-section at lowest temperature ($\sim 10\ \text{keV}$).
\item \textbf{D-D:} $\ce{^{2}_{1}H + ^{2}_{1}H -> ^{3}_{1}H + ^{1}_{1}H}$ (4.0 MeV) or $\ensuremath{\mathrm{^{3}_{2}He + ^{1}_{0}n}}$ (3.3 MeV).
\end{itemize}
\textbf{Advantages:} No long-lived high-level waste (activation of vessel only), abundant fuel (Deuterium in seawater $\approx 0.015\%$, Tritium bred from Lithium: $\ce{^{6}Li + n -> ^{4}He + T}$), inherent safety (no chain reaction, meltdown impossible).
\end{propriete}

\begin{methode}[Energy Released in Nuclear Fusion]
\begin{enumerate}
\item \textbf{Write balanced equation:} Check $A$, $Z$ conservation.
\item \textbf{Calculate $Q$-value:} $Q = (\sum m_{\text{initial}} - \sum m_{\text{final}}) c^2$. Use atomic masses (electrons balance if $Z$ conserved).
\item \textbf{Kinetic Energy Partition:} For 2-body final state at rest, momentum conservation $\Rightarrow p_1 = p_2$. $T_1/T_2 = m_2/m_1$. $Q = T_1 + T_2$.
\item \textbf{Fuel Energy Density:} Moles of fuel $\to$ Number of reactions $\times Q$.
\end{enumerate}
\end{methode}

\begin{exoresolu}[Deuterium-Tritium Fusion Energy]
Consider the D-T fusion reaction: $\ce{^{2}_{1}H + ^{3}_{1}H -> ^{4}_{2}He + ^{1}_{0}n}$.
Atomic masses: $\ce{^{2}H} = 2.014102\ \text{u}$, $\ce{^{3}H} = 3.016049\ \text{u}$, $\ce{^{4}He} = 4.002603\ \text{u}$, neutron = 1.008665 u.
\begin{enumerate}
\item Verify the $Q$-value is \SI{17.6}{MeV}.
\item Calculate the kinetic energy of the neutron and the alpha particle (Helium-4 nucleus) assuming the reactants are at rest.
\item \SI{1}{g} of Deuterium fuses with \SI{1.5}{g} of Tritium (stoichiometric). Calculate the energy released in Joules and compare to \SI{1}{tonne} of coal (\SI{30}{GJ}).
\end{enumerate}
\tcblower
\textbf{Detailed Solution}

\textbf{1. $Q$-value verification:}
Initial mass $m_i = 2.014102 + 3.016049 = 5.030151\ \text{u}$.
Final mass $m_f = 4.002603 + 1.008665 = 5.011268\ \text{u}$.
$\Delta m = 5.030151 - 5.011268 = 0.018883\ \text{u}$.
$Q = 0.018883 \times 931.5 = \SI{17.59}{MeV} \approx \SI{17.6}{MeV}$. Verified.

\textbf{2. Kinetic Energy Partition:}
Momentum conservation: $p_n = p_\alpha$ (magnitudes).
$T = p^2 / 2m \Rightarrow T_n / T_\alpha = m_\alpha / m_n \approx 4 / 1 = 4$.
$Q = T_n + T_\alpha = T_n + T_n/4 = 1.25 T_n$.
$T_n = Q / 1.25 = 17.6 / 1.25 = \SI{14.08}{MeV}$.
$T_\alpha = Q - T_n = 17.6 - 14.08 = \SI{3.52}{MeV}$.
(Neutron carries 80\% of energy $\to$ major challenge for material shielding and heat extraction in blankets).

\textbf{3. Energy from 1g Deuterium + 1.5g Tritium:}
Molar mass D = 2 g/mol, T = 3 g/mol.
Moles of D $= 1/2 = 0.5\ \text{mol}$. Moles of T $= 1.5/3 = 0.5\ \text{mol}$.
Number of pairs $N = 0.5 \times N_A = 0.5 \times 6.022 \times 10^{23} = 3.011 \times 10^{23}$.
Total Energy $E = N \times Q = 3.011 \times 10^{23} \times 17.6 \times 1.60 \times 10^{-13}\ \text{J}$.
$E = 3.011 \times 10^{23} \times 2.816 \times 10^{-12} = 8.48 \times 10^{11}\ \text{J} = \SI{848}{GJ}$.
Comparison: $848\ \text{GJ} / 30\ \text{GJ/tonne} \approx 28.3$ tonnes of coal.
\textbf{Conclusion:} \SI{2.5}{g} of fusion fuel $\approx$ \SI{28}{tonnes} of coal. Energy density $\sim 10^7$ times chemical.
\end{exoresolu}

\begin{methode}[Comprehensive Nuclear Physics Problem Solving]
\begin{enumerate}
\item \textbf{Read carefully:} Identify the process (decay, fission, fusion, dating, stability).
\item \textbf{List Data:} Extract masses, half-lives, energies, constants from text or data sheet.
\item \textbf{Conservation Checks:} Nucleon number $A$, Proton number $Z$, Energy-Momentum, Charge.
\item \textbf{Select Formulas:}
\begin{itemize}
\item Decay: $N=N_0 e^{-\lambda t}$, $A=\lambda N$, $T_{1/2}=\ln 2/\lambda$.
\item Energy: $Q = \Delta m c^2$, $E_b = \Delta m c^2$, $T_\alpha = Q \frac{A_{\text{daughter}}}{A_{\text{daughter}}+4}$.
\item Dating: $t = -T_{1/2} \ln(A/A_0)/\ln 2$.
\end{itemize}
\item \textbf{Unit Consistency:} Convert all to SI (kg, s, J) or Nuclear (u, MeV) consistently. $1\ \text{u} = 931.5\ \text{MeV}/c^2 = 1.6605 \times 10^{-27}\ \text{kg}$.
\item \textbf{Significant Figures:} Match input data (usually 3 or 4 s.f.). State units clearly.
\item \textbf{Physical Comment:} Interpret result (e.g., "This isotope is stable", "Energy is released", "Age is within range").
\end{enumerate}
\end{methode}

\begin{exoresolu}[Analysis of a Medical Radioisotope: Technetium-99m]
Technetium-99m ($\ensuremath{\mathrm{^{99m}_{43}Tc}}$, $m$ = metastable) is the workhorse of nuclear medicine (SPECT imaging). It decays by isomeric transition (IT) to $\ensuremath{\mathrm{^{99}_{43}Tc}}$ (ground state) emitting a \SI{140}{keV} $\gamma$ ray. $T_{1/2} = \SI{6.01}{hours}$. The ground state $\ce{^{99}Tc}$ decays by $\beta^-$ to stable $\ensuremath{\mathrm{^{99}_{44}Ru}}$ with $T_{1/2} = \SI{2.13e5}{years}$.
A $\ce{^{99}Mo}/\ce{^{99m}Tc}$ generator produces $\ce{^{99m}Tc}$ via $\beta^-$ decay of $\ce{^{99}Mo}$ ($T_{1/2} = \SI{66.0}{hours}$). At $t=0$ (elution), the generator yields \SI{1.00}{GBq} of $\ce{^{99m}Tc}$.
\begin{enumerate}
\item Write the decay equations for $\ce{^{99}Mo}$, $\ce{^{99m}Tc}$ (IT), and $\ce{^{99}Tc}$ ($\beta^-$).
\item Calculate the activity of $\ce{^{99m}Tc}$ remaining \SI{12.0}{hours} after elution (ignore $\ce{^{99}Mo}$ breakthrough).
\item Why is the \SI{140}{keV} $\gamma$ ray ideal for gamma camera imaging? (Consider penetration in tissue and collimator efficiency).
\item The $\ce{^{99}Tc}$ daughter accumulates. Calculate its activity \SI{24}{hours} after elution. Is it a significant radiation hazard to the patient? (Assume no biological clearance).
\end{enumerate}
\tcblower
\textbf{Detailed Solution}

\textbf{1. Decay Equations:}
$\ensuremath{\mathrm{^{99}_{42}Mo \rightarrow  ^{99m}_{43}Tc + ^{0}_{-1}e + \bar{\nu}_e}}$ \quad ($T_{1/2}=66.0\ \text{h}$, $\beta^-$ decay)
$\ensuremath{\mathrm{^{99m}_{43}Tc \rightarrow  ^{99}_{43}Tc + \gamma}}$ \quad ($T_{1/2}=6.01\ \text{h}$, Isomeric Transition, $E_\gamma=140\ \text{keV}$)
$\ensuremath{\mathrm{^{99}_{43}Tc \rightarrow  ^{99}_{44}Ru + ^{0}_{-1}e + \bar{\nu}_e}}$ \quad ($T_{1/2}=2.13\times 10^5\ \text{yr}$, $\beta^-$ decay, $E_{\beta,\text{max}}=292\ \text{keV}$)

\textbf{2. Activity of $\ce{^{99m}Tc}$ at $t=12.0\ \text{h}$:}
$\lambda_{\text{Tc}} = \ln 2 / 6.01 = 0.1153\ \text{h}^{-1}$.
$A = A_0 e^{-\lambda t} = 1.00 \times e^{-0.1153 \times 12.0} = 1.00 \times e^{-1.384} = 1.00 \times 0.250 = \SI{0.250}{GBq} = \SI{250}{MBq}$.
(Exactly 2 half-lives: $12.0 / 6.01 \approx 2.00 \Rightarrow 1/4$ activity).

\textbf{3. Suitability of 140 keV $\gamma$:}
\begin{itemize}
\item \textbf{Penetration:} In soft tissue (water), linear attenuation coefficient $\mu \approx 0.15\ \text{cm}^{-1}$ at 140 keV. Half-value layer (HVL) $\approx 4.6\ \text{cm}$. Allows imaging of deep organs (heart, brain) with sufficient photons escaping the body.
\item \textbf{Collimation:} Lead collimators (parallel hole) require septa thick enough to stop $\gamma$. For 140 keV, septa $\approx 0.2-0.3\ \text{mm}$ Pb suffice (HVL in Pb $\approx 0.03\ \text{cm}$). Allows high geometric resolution (hole diameter $\sim 1-2\ \text{mm}$). Higher energy (e.g., 511 keV) needs thick septa $\to$ poor resolution/sensitivity. Lower energy (e.g., 80 keV) absorbed in tissue $\to$ poor depth penetration.
\item \textbf{Detection:} NaI(Tl) scintillators have peak efficiency $\sim 80-90\%$ at 140 keV.
\end{itemize}
\textbf{Conclusion:} Optimal compromise between tissue penetration and collimator/detector performance.

\textbf{4. Activity of $\ce{^{99}Tc}$ at $t=24\ \text{h}$:}
$\ce{^{99m}Tc}$ decays to $\ce{^{99}Tc}$. This is a Parent ($P$) $\to$ Daughter ($D$) decay where $\lambda_P \gg \lambda_D$ ($T_{1/2,P}=6\ \text{h} \ll T_{1/2,D}=2.13\times 10^5\ \text{yr}$).
Activity of daughter $A_D(t) = A_P(0) \frac{\lambda_D}{\lambda_D - \lambda_P} (e^{-\lambda_P t} - e^{-\lambda_D t}) + A_D(0) e^{-\lambda_D t}$.
Since $\lambda_D \approx 0$ (very long half-life), $e^{-\lambda_D t} \approx 1$, $\frac{\lambda_D}{\lambda_D - \lambda_P} \approx 0$.
$A_D(t) \approx A_P(0) (1 - e^{-\lambda_P t})$.
$A_P(0) = \SI{1.00}{GBq}$. $\lambda_P = 0.1153\ \text{h}^{-1}$. $t=24\ \text{h}$.
$e^{-\lambda_P t} = e^{-0.1153 \times 24} = e^{-2.767} = 0.063$.
$A_D(24\ \text{h}) \approx 1.00 \times (1 - 0.063) = \SI{0.937}{GBq} = \SI{937}{MBq}$.
\textbf{Hazard Assessment:}
$\ce{^{99}Tc}$ emits low-energy $\beta^-$ ($E_{\text{max}}=292\ \text{keV}$, $E_{\text{avg}} \approx 100\ \text{keV}$). Range in tissue $\approx 0.5\ \text{mm}$. No significant $\gamma$ emission.
Dose is highly localized. However, $\ce{^{99}Tc}$ chemistry (pertechnetate $\ce{TcO4-}$) concentrates in thyroid, stomach, salivary glands.
Effective dose coefficient (ingestion) $\approx 6.4 \times 10^{-10}\ \text{Sv/Bq}$.
Committed dose $\approx 9.37 \times 10^8 \times 6.4 \times 10^{-10} \approx \SI{0.6}{mSv}$.
This is low (background $\sim 2.4\ \text{mSv/yr}$, CT scan $\sim 10\ \text{mSv}$). \textbf{Not a significant acute hazard}, but contributes to long-term stochastic risk. Justifies use in diagnostics.
\end{exoresolu}

\begin{attention}[Common Mistakes in Radioactivity]
\begin{itemize}
\item \textbf{Atomic vs Nuclear Masses:} Always use \textbf{atomic masses} from data tables for $Q$-value calculations. They include electrons which cancel correctly if $Z$ is conserved. Using nuclear masses requires manual electron accounting.
\item \textbf{Relativistic $\beta$ particles:} Never use $K = \frac{1}{2}mv^2$ for $\beta$ particles above $\sim 0.1\ \text{MeV}$. Use $E = \gamma m_0 c^2$, $p = \gamma m_0 v$.
\item \textbf{Half-life units:} Ensure $\lambda$ and $t$ have consistent time units (both seconds, or both years). $\lambda = \ln 2 / T_{1/2}$.
\item \textbf{Decay Series Counting:} For $\alpha$ count: $\Delta A / 4$. For $\beta$ count: Solve $-2n_\alpha + n_\beta = \Delta Z$. Don't just count arrows in a diagram if not given.
\item \textbf{Activity vs Number of Nuclei:} $A = \lambda N$. Don't confuse $A$ (Bq) with $N$ (dimensionless count).
\item \textbf{Carbon Dating Limit:} $\approx 50,000$ years ($\sim 8-10$ half-lives). Beyond that, use other isotopes ($\ce{^{40}K}$, $\ce{^{238}U}$).
\end{itemize}
\end{attention}

\begin{exercice}[GCE Style Practice Questions]
\begin{itemize}
\item A sample of radioactive isotope has an initial activity of \SI{800}{Bq}. After \SI{12}{hours}, the activity is \SI{100}{Bq}. Calculate the half-life of the isotope.
\item Write the nuclear equation for the $\beta^+$ decay of $\ensuremath{\mathrm{^{22}_{11}Na}}$. The daughter nucleus is $\ensuremath{\mathrm{^{22}_{10}Ne}}$. Calculate the maximum kinetic energy of the positron given: $M(\ce{^{22}Na}) = 21.994437\ \text{u}$, $M(\ce{^{22}Ne}) = 21.991385\ \text{u}$. ($1\ \text{u} = 931.5\ \text{MeV}$).
\item In a fission reactor, \SI{1.00}{kg} of $\ce{^{235}U}$ undergoes fission. Assuming an average energy release of \SI{200}{MeV} per fission, calculate the total energy released in Joules.
\item Explain why high temperatures (of order $10^7\ \text{K}$) are necessary for nuclear fusion to occur. Explain the role of quantum tunnelling.
\item A rock sample contains $\ce{^{238}U}$ and $\ce{^{206}Pb}$ in a molar ratio of 1:3. Assuming no initial $\ce{^{206}Pb}$ and no leakage, calculate the age of the rock. ($T_{1/2}(\ce{^{238}U}) = \SI{4.47e9}{years}$).
\end{itemize}
\end{exercice}

\begin{corrige}[Exercise 1]
$A/A_0 = 100/800 = 1/8 = (1/2)^3$. So 3 half-lives have passed in 12 hours.
$T_{1/2} = 12 / 3 = \SI{4}{hours}$.
\end{corrige}

\begin{corrige}[Exercise 2]
Equation: $\ensuremath{\mathrm{^{22}_{11}Na \rightarrow  ^{22}_{10}Ne + ^{0}_{+1}e + \nu_e}}$.
$Q = [M(\ce{^{22}Na}) - M(\ce{^{22}Ne}) - 2m_e] c^2$. (Atomic mass of Na includes 11 e$^-$, Ne includes 10 e$^-$. Positron emission requires creation of $e^+$ mass, equivalent to losing 2 electron masses from atomic mass balance).
$\Delta m = 21.994437 - 21.991385 - 2(0.00054858) = 0.003051 - 0.001097 = 0.001954\ \text{u}$.
$Q = 0.001954 \times 931.5 = \SI{1.82}{MeV}$.
$K_{\text{max}}(e^+) \approx Q = \SI{1.82}{MeV}$ (neglecting recoil of heavy Ne nucleus).
\end{corrige}

\begin{corrige}[Exercise 3]
Number of atoms in 1 kg $\ce{^{235}U}$: $N = \frac{1000\ \text{g}}{235\ \text{g/mol}} \times 6.022 \times 10^{23} = 2.56 \times 10^{24}$.
Energy per fission $= 200 \times 1.60 \times 10^{-13} = 3.20 \times 10^{-11}\ \text{J}$.
Total Energy $= 2.56 \times 10^{24} \times 3.20 \times 10^{-11} = \SI{8.19e13}{J}$.
\end{corrige}

\begin{corrige}[Exercise 4]
High temperature gives nuclei high kinetic energy ($E_k \approx k_B T$) to approach within range of strong force ($\sim 1-3\ \text{fm}$), overcoming Coulomb repulsion $V_C = \frac{Z_1 Z_2 e^2}{4\pi\epsilon_0 r}$. At $T \sim 10^7\ \text{K}$, $k_B T \approx 1\ \text{keV}$, while $V_C \sim 100-1000\ \text{keV}$. Classical physics forbids fusion. \textbf{Quantum Tunnelling:} Wave nature of particles allows a finite probability to penetrate the classically forbidden barrier region. The Gamow factor gives the probability $\propto \exp(-2\pi Z_1 Z_2 e^2 / (4\pi\epsilon_0 \hbar v))$. Without tunnelling, fusion rates in stars would be zero.
\end{corrige}

\begin{corrige}[Exercise 5]
Molar ratio $\ce{^{206}Pb} : \ce{^{238}U} = 3:1$.
Let $N_0$ be initial U atoms. $N_U = N_0 e^{-\lambda t}$. $N_{Pb} = N_0 - N_U = N_0(1 - e^{-\lambda t})$.
Ratio $N_{Pb}/N_U = (1 - e^{-\lambda t}) / e^{-\lambda t} = e^{\lambda t} - 1 = 3$.
$e^{\lambda t} = 4 \Rightarrow \lambda t = \ln 4 = 2\ln 2$.
$t = 2 \ln 2 / \lambda = 2 T_{1/2} = 2 \times 4.47 \times 10^9 = \SI{8.94e9}{years}$.
\end{corrige}

\newpage

\coursec{Exercises}

\courssub{I check my knowledge}

\begin{exercice}[Multiple choice questions]
\begin{enumerate}
    \item Which of the following radiations has the highest ionising power?
    \begin{enumerate}
        \item $\alpha$-particles
        \item $\beta$-particles
        \item $\gamma$-rays
        \item X-rays
    \end{enumerate}

    \item A radioactive nucleus $^{A}_{Z}X$ decays by emitting an $\alpha$-particle. The resulting nucleus is:
    \begin{enumerate}
        \item $^{A-4}_{Z-2}Y$
        \item $^{A-4}_{Z+2}Y$
        \item $^{A}_{Z-1}Y$
        \item $^{A}_{Z+1}Y$
    \end{enumerate}

    \item The half-life of a radioactive isotope is 10 days. What fraction of the original sample remains after 30 days?
    \begin{enumerate}
        \item $\frac{1}{2}$
        \item $\frac{1}{4}$
        \item $\frac{1}{8}$
        \item $\frac{1}{16}$
    \end{enumerate}

    \item In a nuclear fission reaction, the total mass of the products is:
    \begin{enumerate}
        \item greater than the total mass of the reactants
        \item equal to the total mass of the reactants
        \item less than the total mass of the reactants
        \item zero
    \end{enumerate}
\end{enumerate}
\end{exercice}

\begin{exercice}[True or False with justification]
State whether each statement is True or False and justify your answer briefly.
\begin{enumerate}
    \item $\gamma$-rays are deflected by a magnetic field.
    \item The mass defect of a nucleus is the difference between the mass of the nucleus and the sum of the masses of its constituent nucleons.
    \item In $\beta^{-}$ decay, the atomic number $Z$ decreases by 1 while the mass number $A$ remains constant.
    \item Nuclear fusion requires very high temperatures to overcome the electrostatic repulsion between nuclei.
\end{enumerate}
\end{exercice}

\begin{exercice}[Definitions and direct calculations]
\begin{enumerate}
    \item Define the following terms: (a) Half-life ($T_{1/2}$), (b) Activity ($A$), (c) Binding energy per nucleon.
    \item The decay constant of a radioactive substance is $\lambda = \SI{0.02}{s^{-1}}$. Calculate its half-life.
    \item A sample has an initial activity of $\SI{800}{Bq}$. After 12 hours, its activity is $\SI{100}{Bq}$. Determine the half-life of the sample.
    \item Write the balanced nuclear equation for the $\alpha$-decay of $^{226}_{88}\text{Ra}$ and the $\beta^{-}$-decay of $^{14}_{6}\text{C}$.
\end{enumerate}
\end{exercice}

\begin{exercice}[Radiation identification and properties]
\begin{enumerate}
    \item A radiation from a radioactive source is stopped by a sheet of paper and is strongly deflected by a magnetic field in a direction opposite to that of $\beta$-particles. Identify the radiation and state its charge.
    \item Complete the following nuclear equations and identify the type of decay in each case:
\begin{align*}
        &\ce{^{238}_{92}U -> ^{234}_{90}Th + ^{4}_{2}He} \\
        &\ce{^{14}_{6}C -> ^{14}_{7}N + ^{0}_{-1}e} \\
        &\ce{^{214}_{83}Bi -> ^{214}_{84}Po + ^{0}_{-1}e}
    \end{align*}
    \item For each radiation identified in (ii), state two distinct properties.
\end{enumerate}
\end{exercice}

\courssub{I practise}

\begin{exercice}[Activity and decay law]
A radioactive source used in a hospital has a half-life of 6 hours. At 8:00 am, its activity is measured as $\SI{1.6e6}{Bq}$.
\begin{enumerate}
    \item Calculate the decay constant $\lambda$ in $\SI{}{s^{-1}}$.
    \item What will be the activity at 2:00 pm on the same day?
    \item How long will it take for the activity to fall to $\SI{100000}{Bq}$?
    \item Sketch a graph of $\ln(A)$ against time $t$, indicating the slope and intercept.
\end{enumerate}
\end{exercice}

\begin{exercice}[Mass defect and binding energy]
Given the following masses:
$m_p = \SI{1.007276}{u}$, $m_n = \SI{1.008665}{u}$, $m(^{4}_{2}\text{He}) = \SI{4.002603}{u}$, $1\,\text{u} = \SI{931.5}{MeV/c^2}$.
\begin{enumerate}
    \item Calculate the mass defect of the helium-4 nucleus in u and in kg.
    \item Calculate the total binding energy of the nucleus in MeV and in joules.
    \item Determine the binding energy per nucleon for helium-4.
    \item Comment on the stability of helium-4 relative to lighter nuclei using the binding energy per nucleon curve.
\end{enumerate}
\end{exercice}

\begin{exercice}[Nuclear fission energy calculation]
Consider the fission reaction induced by a thermal neutron:
\[\ce{^{235}_{92}U + ^{1}_{0}n -> ^{94}_{38}Sr + ^{140}_{54}Xe + 2^{1}_{0}n}\]
Atomic masses: $m(^{235}\text{U}) = \SI{235.04393}{u}$, $m(^{94}\text{Sr}) = \SI{93.91536}{u}$, $m(^{140}\text{Xe}) = \SI{139.92164}{u}$, $m_n = \SI{1.008665}{u}$.
\begin{enumerate}
    \item Verify that the nucleon number and proton number are conserved.
    \item Calculate the mass defect $\Delta m$ for this reaction in u.
    \item Calculate the energy released per fission in MeV and in joules.
    \item Calculate the energy released by the fission of $\SI{1}{g}$ of pure $^{235}\text{U}$.
    \item Compare this energy with the heat of combustion of $\SI{1}{g}$ of petrol ($\approx \SI{47}{MJ}$). How many grams of petrol would be equivalent to $\SI{1}{g}$ of $^{235}\text{U}$?
\end{enumerate}
\end{exercice}

\begin{exercice}[Carbon-14 dating]
The half-life of $^{14}\text{C}$ is 5730 years. A piece of wood from an archaeological site shows an activity of 25\% of that found in living wood.
\begin{enumerate}
    \item Explain the principle of carbon-14 dating, mentioning the role of cosmic rays and the carbon cycle.
    \item Calculate the age of the piece of wood.
    \item If the measurement of activity has an uncertainty of $\pm 2\%$, estimate the uncertainty in the calculated age.
    \item Why is carbon-14 dating not suitable for samples older than about 50,000 years?
\end{enumerate}
\end{exercice}

\courssub{I prepare for the GCE}

\begin{exercice}[{Structured essay: Nuclear Reactors and Energy in Cameroon] [20 marks}]
\begin{enumerate}
    \item Draw a labelled diagram of a typical thermal nuclear reactor (e.g., Pressurised Water Reactor) showing the following components: reactor core, fuel rods, moderator, control rods, coolant, heat exchanger (steam generator), turbine, generator, condenser, cooling tower. [6 marks]
    \item Explain the function of the moderator and the control rods in the reactor core. Why is it necessary to slow down neutrons? [4 marks]
    \item Describe the process of nuclear fission of $^{235}\text{U}$ initiated by a thermal neutron, including the concept of a chain reaction and critical mass. [4 marks]
    \item Discuss two advantages and two risks/disadvantages of generating electricity from nuclear power specifically in the context of Cameroon's energy needs and environment. [6 marks]
\end{enumerate}
\end{exercice}

\begin{exercice}[{Binding energy curve, fission and fusion] [15 marks}]
\begin{enumerate}
    \item Sketch the binding energy per nucleon ($B/A$) against nucleon number ($A$) curve (Aston curve). Indicate the approximate position of iron-56 ($^{56}\text{Fe}$), uranium-235 ($^{235}\text{U}$), and helium-4 ($^{4}\text{He}$). [4 marks]
    \item Using the curve, explain qualitatively why:
    \begin{enumerate}
        \item Energy is released when $^{235}\text{U}$ undergoes fission into two medium-mass nuclei.
        \item Energy is released when two light nuclei (e.g., deuterium and tritium) undergo fusion.
    \end{enumerate} [4 marks]
    \item For the fusion reaction $\ce{^{2}_{1}H + ^{3}_{1}H -> ^{4}_{2}He + ^{1}_{0}n}$, the masses are:
    $m(^{2}\text{H}) = \SI{2.014102}{u}$, $m(^{3}\text{H}) = \SI{3.016049}{u}$, $m(^{4}\text{He}) = \SI{4.002603}{u}$, $m_n = \SI{1.008665}{u}$.
    Calculate the energy released in this reaction in MeV. [4 marks]
    \item State two major difficulties in achieving controlled nuclear fusion on Earth for power generation. [3 marks]
\end{enumerate}
\end{exercice}

\begin{exercice}[{Comprehensive radioactivity problem] [15 marks}]
A radioactive series starts with nuclide $^{238}_{92}\text{U}$ and ends with stable $^{206}_{82}\text{Pb}$. The series involves a sequence of $\alpha$ and $\beta^{-}$ decays.
\begin{enumerate}
    \item Determine the total number of $\alpha$ particles and $\beta^{-}$ particles emitted in the complete decay of one atom of $^{238}\text{U}$ to $^{206}\text{Pb}$. [3 marks]
    \item The half-life of $^{238}\text{U}$ is $\SI{4.468e9}{years}$. Calculate its decay constant in $\SI{}{s^{-1}}$. [2 marks]
    \item A rock sample contains $\SI{1.00}{g}$ of $^{238}\text{U}$ and no lead initially. Assuming a closed system, calculate the mass of $^{206}\text{Pb}$ present after $\SI{4.468e9}{years}$ (one half-life). [4 marks]
    \item The decay series includes the gas $^{222}_{86}\text{Rn}$ (radon) with a half-life of 3.82 days. Explain why radon accumulation in basements is a health hazard and how it relates to the uranium content of the soil. [3 marks]
    \item In a cloud chamber photograph, a track is observed that is thick, straight, and short (a few cm in air). Another track is thin, tortuous, and much longer. Identify the radiations responsible for each track and explain the differences in appearance based on their interaction with matter. [3 marks]
\end{enumerate}
\end{exercice}

\newpage

\coursec{Solutions to the exercises}

\begin{corrige}[Exercise 1 — Multiple choice questions]
\begin{enumerate}
    \item \textbf{Answer: (a) $\alpha$-particles.}\par
    The ionising power of a radiation depends on its charge and its speed (or mass for a given energy): the greater the charge and the slower the particle, the more ion pairs it creates per centimetre of path. The $\alpha$-particle ($^{4}_{2}\mathrm{He}^{2+}$, charge $+2e$, mass $\approx 4\ \mathrm{u}$) is heavy and doubly charged, so it interacts very strongly with matter and produces about $10^{4}$ to $10^{5}$ ion pairs per cm in air. $\beta$-particles (charge $-e$, very light) ionise about 100 times less, and $\gamma$-rays and X-rays, being uncharged photons, are the least ionising.
    \item \textbf{Answer: (a) $^{A-4}_{Z-2}Y$.}\par
    In $\alpha$-decay the nucleus emits $^{4}_{2}\mathrm{He}$. Conservation of nucleon number gives $A' = A - 4$ and conservation of charge gives $Z' = Z - 2$:
    \[\ce{^{A}_{Z}X -> ^{A-4}_{Z-2}Y + ^{4}_{2}He}\]
    \item \textbf{Answer: (c) $\frac{1}{8}$.}\par
    30 days corresponds to $n = \dfrac{30}{10} = 3$ half-lives. The remaining fraction is
    \[\frac{N}{N_0} = \left(\frac{1}{2}\right)^{3} = \frac{1}{8}.\]
    \item \textbf{Answer: (c) less than the total mass of the reactants.}\par
    Fission releases energy, and by Einstein's relation $E = \Delta m\,c^{2}$ this energy corresponds to a mass defect $\Delta m > 0$: the products are lighter than the reactants, the ``missing'' mass having been converted into kinetic energy of the fragments and radiation.
\end{enumerate}
\end{corrige}

\begin{corrige}[Exercise 2 — True or False with justification]
\begin{enumerate}
    \item \textbf{False.} $\gamma$-rays are electromagnetic radiation (photons); they carry \emph{no electric charge}. The magnetic (Lorentz) force $\vec{F} = q\,\vec{v}\wedge\vec{B}$ is zero for $q = 0$, so $\gamma$-rays travel straight through a magnetic field. Only charged radiations ($\alpha$, $\beta$) are deflected.
    \item \textbf{True.} The mass defect is defined by
    \[\Delta m = \big[Z\,m_p + (A-Z)\,m_n\big] - m_{\text{nucleus}} > 0.\]
    The nucleus is lighter than the sum of its separated nucleons; the difference $\Delta m$ corresponds, through $E = \Delta m c^{2}$, to the binding energy released when the nucleus formed.
    \item \textbf{False.} In $\beta^{-}$ decay a neutron transforms into a proton inside the nucleus: $\ensuremath{\mathrm{^{1}_{0}n \rightarrow  ^{1}_{1}p + ^{0}_{-1}e + \bar{\nu}_e}}$. The mass number $A$ stays constant (correct), but the atomic number \emph{increases} by 1: $\ce{^{A}_{Z}X -> ^{A}_{Z+1}Y + ^{0}_{-1}e}$. (It is in $\beta^{+}$ decay that $Z$ decreases by 1.)
    \item \textbf{True.} Both nuclei are positively charged and repel each other by the Coulomb force. To bring them within the very short range of the strong nuclear force ($\sim 10^{-15}\ \mathrm{m}$) so that fusion can occur, they must collide with very high kinetic energy. This requires temperatures of the order of $10^{7}$--$10^{8}\ \mathrm{K}$, as in the core of the Sun or in a tokamak plasma.
\end{enumerate}
\end{corrige}

\begin{corrige}[Exercise 3 — Definitions and direct calculations]
\begin{enumerate}
    \item \textbf{Definitions.}
    \begin{enumerate}
        \item \textbf{Half-life} $T_{1/2}$: the time after which half of the radioactive nuclei initially present in a sample have disintegrated (equivalently, the time for the activity to fall to half its initial value).
        \item \textbf{Activity} $A$: the average number of disintegrations per unit time in a sample, $A = \lambda N = -\dfrac{dN}{dt}$; SI unit: the becquerel ($1\ \mathrm{Bq} = 1$ disintegration per second).
        \item \textbf{Binding energy per nucleon} $B/A$: the binding energy of a nucleus divided by its number of nucleons, $\dfrac{B}{A} = \dfrac{\Delta m\,c^{2}}{A}$. It measures the average energy needed to remove one nucleon; the larger it is, the more stable the nucleus.
    \end{enumerate}
    \item Half-life from $\lambda = \SI{0.02}{s^{-1}}$:
    \[T_{1/2} = \frac{\ln 2}{\lambda} = \frac{0.693}{0.02} \approx \SI{34.7}{s} \approx \SI{35}{s}.\]
    \item The activity fell from 800 Bq to 100 Bq, i.e.\ it was divided by $8 = 2^{3}$. Hence 12 hours correspond to 3 half-lives:
    \[T_{1/2} = \frac{12\ \mathrm{h}}{3} = \SI{4}{h}.\]
    \item Balanced equations (conservation of $A$ and $Z$):
    \[\ce{^{226}_{88}Ra -> ^{222}_{86}Rn + ^{4}_{2}He} \qquad (\alpha\text{-decay})\]
    \[\ce{^{14}_{6}C -> ^{14}_{7}N + ^{0}_{-1}e} \qquad (\beta^{-}\text{-decay})\]
    Check: for radium, $226 = 222 + 4$ and $88 = 86 + 2$; for carbon, $14 = 14 + 0$ and $6 = 7 - 1$. \checkmark
\end{enumerate}
\end{corrige}

\begin{corrige}[Exercise 4 — Radiation identification and properties]
\begin{enumerate}
    \item A radiation stopped by a single sheet of paper has very low penetrating power: this is characteristic of an \cle{$\alpha$-particle}. Its deflection is \emph{opposite} to that of $\beta$-particles (which are negative), confirming that it is \textbf{positively charged}: it is the helium nucleus $^{4}_{2}\mathrm{He}^{2+}$, charge $q = +2e = +\SI{3.2e-19}{C}$.
    \item Completed equations and decay types:
    \begin{itemize}
        \item $\ce{^{238}_{92}U -> ^{234}_{90}Th + ^{4}_{2}He}$: \textbf{$\alpha$-decay} ($A$ decreases by 4, $Z$ by 2).
        \item $\ce{^{14}_{6}C -> ^{14}_{7}N + ^{0}_{-1}e}$: \textbf{$\beta^{-}$-decay} ($A$ constant, $Z$ increases by 1).
        \item $\ce{^{214}_{83}Bi -> ^{214}_{84}Po + ^{0}_{-1}e}$: \textbf{$\beta^{-}$-decay} ($A$ constant, $Z$ increases by 1).
    \end{itemize}
    \item Two properties of each:
    \begin{itemize}
        \item \textbf{$\alpha$-particles}: (i) very high ionising power (thick, short, straight tracks in a cloud chamber); (ii) very low penetrating power — stopped by paper or a few cm of air.
        \item \textbf{$\beta^{-}$-particles}: (i) moderate ionising power, thin and tortuous tracks; (ii) moderate penetrating power — stopped by a few mm of aluminium; they are deflected by magnetic fields (negative charge).
    \end{itemize}
\end{enumerate}
\end{corrige}

\begin{corrige}[Exercise 5 — Activity and decay law]
\begin{enumerate}
    \item Decay constant:
    \[\lambda = \frac{\ln 2}{T_{1/2}} = \frac{0.693}{6 \times 3600\ \mathrm{s}} = \frac{0.693}{21600} \approx \SI{3.21e-5}{s^{-1}}.\]
    \item From 8:00 am to 2:00 pm, $\Delta t = 6\ \mathrm{h} = 1$ half-life. The activity is halved:
    \[A = \frac{A_0}{2} = \frac{\SI{1.6e6}{Bq}}{2} = \SI{8.0e5}{Bq}.\]
    \item Using the decay law $A = A_0 e^{-\lambda t}$ with $A = \SI{1.0e5}{Bq}$:
    \[t = \frac{1}{\lambda}\ln\frac{A_0}{A} = \frac{T_{1/2}}{\ln 2}\ln\frac{1.6\times10^{6}}{1.0\times10^{5}} = \frac{6\ \mathrm{h}}{0.693}\times\ln 16 = \frac{6}{0.693}\times 2.773 \approx \SI{24}{h}.\]
    (Equivalently: $1.6\times10^{6}/10^{5} = 16 = 2^{4}$, so 4 half-lives $= 24$ h.)
    \item Since $A = A_0 e^{-\lambda t}$, taking logarithms: $\ln A = \ln A_0 - \lambda t$. The graph of $\ln A$ versus $t$ is a \textbf{straight line} of slope $-\lambda$ and $y$-intercept $\ln A_0$.
    \begin{popfigure}
    \begin{tikzpicture}[scale=1.0]
        \draw[->, thick] (0,0) -- (6.5,0) node[below left] {$t$ (h)};
        \draw[->, thick] (0,0) -- (0,4) node[below left] {$\ln A$};
        \draw[very thick, popblue] (0,3.4) -- (6,0.4);
        \fill[poporange] (0,3.4) circle (2pt) node[left, black] {$\ln A_0$};
        \node[popdark, align=center] at (4.2,2.6) {slope $= -\lambda$};
    \end{tikzpicture}
    \legende{Graph of $\ln A$ against time: a straight line of slope $-\lambda$ and intercept $\ln A_0$.}
    \end{popfigure}
\end{enumerate}
\end{corrige}

\begin{corrige}[Exercise 6 — Mass defect and binding energy]
\begin{enumerate}
    \item The helium-4 nucleus contains $Z = 2$ protons and $A - Z = 2$ neutrons. Using \emph{atomic} masses (which automatically account for the electrons) avoids electron-counting errors:
    \begin{align*}
    \Delta m &= \big[2\,m(\ensuremath{\mathrm{^{1}_{1}H}}) + 2\,m_n\big] - m(\ensuremath{\mathrm{^{4}_{2}He}})\\
    &= \big[2(1.007825) + 2(1.008665)\big] - 4.002603\\
    &= 4.032980 - 4.002603 = \SI{0.030377}{u}.
    \end{align*}
    In kilograms, with $1\ \mathrm{u} = 1.6605\times10^{-27}\ \mathrm{kg}$:
    \[\Delta m = 0.030377 \times 1.6605\times10^{-27} \approx \SI{5.04e-29}{kg}.\]
    \item Total binding energy:
    \[B = \Delta m \times 931.5\ \mathrm{MeV} = 0.030377 \times 931.5 \approx \SI{28.30}{MeV}.\]
    In joules ($1\ \mathrm{MeV} = 1.602\times10^{-13}\ \mathrm{J}$):
    \[B = 28.30 \times 1.602\times10^{-13} \approx \SI{4.53e-12}{J}.\]
    (Check with $E = \Delta m c^{2}$: $5.04\times10^{-29}\times(3.00\times10^{8})^{2} \approx 4.54\times10^{-12}\ \mathrm{J}$. \checkmark)
    \item Binding energy per nucleon:
    \[\frac{B}{A} = \frac{28.30\ \mathrm{MeV}}{4} \approx \SI{7.07}{MeV/nucleon}.\]
    \item \textbf{Comment.} The value $B/A \approx 7.1\ \mathrm{MeV}$ for $^{4}\mathrm{He}$ is far higher than that of its light neighbours ($^{2}\mathrm{H}$: $\approx 1.1$ MeV; $^{3}\mathrm{He}$: $\approx 2.6$ MeV). On the Aston curve, helium-4 sits on a sharp local peak: it is exceptionally stable for such a light nucleus. This explains why $\alpha$-particles are emitted intact in radioactivity and why fusion of light nuclei into helium releases so much energy.
\end{enumerate}
\end{corrige}

\begin{corrige}[Exercise 7 — Nuclear fission energy calculation]
\begin{enumerate}
    \item \textbf{Conservation checks.}
    \begin{itemize}
        \item Nucleons: reactants $235 + 1 = 236$; products $94 + 140 + 2(1) = 236$. \checkmark
        \item Protons: reactants $92 + 0 = 92$; products $38 + 54 + 0 = 92$. \checkmark
    \end{itemize}
    \item Mass defect:
    \begin{align*}
    m_{\text{reactants}} &= 235.04393 + 1.008665 = 236.052595\ \mathrm{u}\\
    m_{\text{products}} &= 93.91536 + 139.92164 + 2(1.008665) = 235.854330\ \mathrm{u}\\
    \Delta m &= 236.052595 - 235.854330 = \SI{0.198265}{u}.
    \end{align*}
    \item Energy released per fission:
    \[E = \Delta m \times 931.5 = 0.198265 \times 931.5 \approx \SI{184.7}{MeV} \approx \SI{185}{MeV}.\]
    In joules:
    \[E = 184.7 \times 1.602\times10^{-13} \approx \SI{2.96e-11}{J}.\]
    \item Number of nuclei in 1 g of $^{235}\mathrm{U}$:
    \[N = \frac{m}{M}\,N_A = \frac{1.00\ \mathrm{g}}{235\ \mathrm{g\,mol^{-1}}}\times 6.022\times10^{23} \approx 2.56\times10^{21}\ \text{nuclei}.\]
    Total energy:
    \[E_{\text{tot}} = N \times E = 2.56\times10^{21}\times 2.96\times10^{-11} \approx \SI{7.6e10}{J} \approx \SI{76}{GJ}.\]
    \item Comparison with petrol (calorific value $\approx \SI{47}{MJ\,kg^{-1}} = \SI{4.7e4}{J\,g^{-1}}$):
    \[m_{\text{petrol}} = \frac{E_{\text{tot}}}{4.7\times10^{4}\ \mathrm{J\,g^{-1}}} = \frac{7.6\times10^{10}}{4.7\times10^{4}} \approx 1.6\times10^{6}\ \mathrm{g} \approx \SI{1.6}{t}.\]
    \textbf{Conclusion:} 1 g of $^{235}\mathrm{U}$ releases roughly the same energy as about \textbf{1.6 tonnes of petrol} — nuclear fuel is about a million times more concentrated in energy than chemical fuel.
\end{enumerate}
\end{corrige}

\begin{corrige}[Exercise 8 — Carbon-14 dating]
\begin{enumerate}
    \item \textbf{Principle.} Cosmic rays striking the upper atmosphere produce neutrons which convert nitrogen-14 into radioactive carbon-14: $\ce{^{14}_{7}N + ^{1}_{0}n -> ^{14}_{6}C + ^{1}_{1}H}$. The $^{14}\mathrm{C}$ oxidises to $\mathrm{CO}_2$, enters the carbon cycle (photosynthesis $\to$ plants $\to$ animals), so living organisms maintain a constant $^{14}\mathrm{C}/^{12}\mathrm{C}$ ratio and a constant activity per gram of carbon. When the organism dies, it stops exchanging carbon with the atmosphere; the $^{14}\mathrm{C}$ then decays with $T_{1/2} = 5730$ years. Measuring the remaining activity and applying $A = A_0 e^{-\lambda t}$ gives the time elapsed since death.
    \item The activity is $25\% = \frac{1}{4} = \left(\frac{1}{2}\right)^{2}$ of the living value, so two half-lives have elapsed:
    \[t = 2 \times 5730 = \SI{11460}{years} \approx \SI{1.15e4}{years}.\]
    \item From $A = A_0 e^{-\lambda t}$, $t = \dfrac{1}{\lambda}\ln\dfrac{A_0}{A}$. Differentiating with respect to $A$:
    \[\Delta t = \frac{1}{\lambda}\,\frac{\Delta A}{A} = \frac{T_{1/2}}{\ln 2}\times 0.02 = \frac{5730}{0.693}\times 0.02 \approx \SI{165}{years}.\]
    So the age is $t = (11460 \pm 165)$ years (from the activity uncertainty alone).
    \item After about 50,000 years, i.e.\ $\sim 9$ half-lives, the remaining fraction is $\left(\frac12\right)^{9} \approx 0.2\%$. The activity is then so weak that it is comparable to background radiation and contamination; the counting uncertainty becomes enormous and the measured ``age'' is unreliable. Other methods (K--Ar, U--Pb dating) are used for older samples.
\end{enumerate}
\end{corrige}

\begin{corrige}[Exercise 9 — Structured essay: Nuclear Reactors and Energy in Cameroon (20 marks)]
\textbf{Indicative mark scheme and model answer.}
\begin{enumerate}
    \item \textbf{Labelled diagram [6 marks]} — 1 mark per correctly placed and labelled group of components (core + fuel rods; moderator; control rods; coolant circuit; heat exchanger; turbine + generator + condenser + cooling tower).
    \begin{popfigure}
    \begin{tikzpicture}[scale=0.9, every node/.style={font=\small}]
        % Reactor vessel
        \draw[very thick, popblue] (0,0) rectangle (3,3.4);
        \node[popblue] at (1.5,3.7) {Reactor core (pressure vessel)};
        % Fuel rods
        \foreach \x in {0.5,1.0,1.5,2.0,2.5} \draw[thick, poporange] (\x,0.3) -- (\x,2.2);
        \node[poporange] at (1.5,-0.35) {Fuel rods ($^{235}$U)};
        % Control rods
        \foreach \x in {0.75,1.75} \draw[very thick, popdark] (\x,3.4) -- (\x,1.4);
        \node[popdark, align=center] at (4.6,3.1) {Control rods (B, Cd)};
        \draw[->, popdark] (3.6,2.95) -- (1.8,2.6);
        % Moderator
        \node[popgreen] at (1.5,2.75) {Moderator (water)};
        % Coolant loop
        \draw[very thick, poppink, ->] (3,2.4) -- (5.5,2.4);
        \draw[very thick, poppink, ->] (5.5,0.8) -- (3,0.8);
        \node[poppink, align=center] at (4.3,2.75) {hot coolant};
        \node[poppink, align=center] at (4.3,0.45) {cool coolant};
        % Heat exchanger
        \draw[very thick, poppurple] (5.5,0.4) rectangle (7.3,2.8);
        \node[poppurple, align=center] at (6.4,3.15) {Heat exchanger (steam generator)};
        % Turbine + generator
        \draw[very thick, popgold, ->] (7.3,2.4) -- (8.6,2.4);
        \draw[thick] (8.6,2.0) -- (9.4,2.8) -- (10.2,2.0) -- cycle;
        \node[align=center] at (9.4,1.7) {Turbine};
        \draw[thick] (10.5,2.0) rectangle (11.5,2.8);
        \node[align=center] at (11.0,3.1) {Generator};
        % Condenser + cooling tower
        \draw[thick, popblue] (8.6,0.4) rectangle (10.2,1.2);
        \node[popblue, align=center] at (9.4,0.1) {Condenser};
        \draw[very thick, popmuted] (11.0,0.2) -- (10.7,1.6) -- (11.9,1.6) -- (11.6,0.2);
        \node[popmuted, align=center] at (11.3,-0.15) {Cooling tower};
        \draw[->, thick] (9.4,2.0) -- (9.4,1.2);
        \draw[->, thick] (10.2,0.8) -- (10.7,0.8);
        \draw[->, thick, poppurple] (8.6,0.8) -- (6.9,0.8) -- (6.9,0.4);
    \end{tikzpicture}
    \legende{Schematic of a pressurised water reactor (PWR) and its power conversion circuit.}
    \end{popfigure}
    \item \textbf{Moderator and control rods [4 marks].}
    \begin{itemize}
        \item \emph{Moderator} (water, heavy water or graphite): slows the fast neutrons released by fission down to thermal energies ($\approx 0.025\ \mathrm{eV}$) by elastic collisions. \textbf{Why slow them?} The probability (cross-section) of $^{235}\mathrm{U}$ capturing a neutron and fissioning is hundreds of times greater for slow (thermal) neutrons than for fast ones; slowing them sustains the chain reaction efficiently. [2 marks]
        \item \emph{Control rods} (boron or cadmium, strong neutron absorbers): inserted or withdrawn to adjust the neutron population and hence the reaction rate; fully inserted, they shut the reactor down. [2 marks]
    \end{itemize}
    \item \textbf{Fission and chain reaction [4 marks].} A thermal neutron is absorbed by $^{235}\mathrm{U}$, forming unstable $^{236}\mathrm{U}$, which splits into two medium-mass fragments plus 2--3 neutrons and $\approx 200\ \mathrm{MeV}$ of energy, e.g.
    \[\ce{^{235}_{92}U + ^{1}_{0}n -> ^{94}_{38}Sr + ^{140}_{54}Xe + 2^{1}_{0}n}.\]
    The emitted neutrons, once moderated, can trigger further fissions: this is the \cle{chain reaction}. If on average exactly one neutron per fission causes a new fission, the reaction is self-sustaining at constant power (\emph{critical} state). The \cle{critical mass} is the minimum mass of fissile material needed so that neutron losses through the surface do not quench the reaction. [1 mark per concept: fission process, neutron multiplication, chain reaction, critical mass.]
    \item \textbf{Advantages and risks for Cameroon [6 marks]} — any two of each, well argued:
    \begin{itemize}
        \item \emph{Advantages:} (i) enormous energy density — a small quantity of fuel supplies large cities and industry, complementing hydroelectricity (which suffers in dry seasons); (ii) no $\mathrm{CO}_2$ emissions during operation, helping preserve Cameroon's forests and climate commitments; (iii) stable, continuous base-load power unlike seasonal hydro.
        \item \emph{Risks/disadvantages:} (i) radioactive waste remains dangerous for thousands of years and Cameroon has no disposal facility; (ii) risk of accidents and need for highly specialised personnel, strict regulation and large capital investment; (iii) reactors need reliable cooling water and siting away from seismic/volcanic zones (e.g.\ Mount Cameroon region).
    \end{itemize}
    \emph{Examiner's expectations:} a clear, correctly labelled diagram; precise vocabulary (thermal neutrons, cross-section, criticality); balanced discussion linked to the Cameroonian context, not a generic list.
\end{enumerate}
\end{corrige}

\begin{corrige}[Exercise 10 — Binding energy curve, fission and fusion (15 marks)]
\begin{enumerate}
    \item \textbf{Aston curve [4 marks]} — 1 mark for axes, 1 for shape (steep rise, broad maximum near Fe, slow decrease), 1 for the three positions, 1 for the He-4 peak.
    \begin{popfigure}
    \begin{tikzpicture}[scale=0.85]
        \draw[->, thick] (0,0) -- (11,0) node[below left] {$A$};
        \draw[->, thick] (0,0) -- (0,5) node[below left] {$B/A$ (MeV)};
        \draw[very thick, popblue, smooth] plot coordinates {(0.3,0.4) (0.8,2.9) (1.2,2.2) (2,3.4) (3.5,4.1) (5,4.35) (6.5,4.3) (8,4.1) (10,3.7)};
        \fill[poporange] (0.8,2.9) circle (2.5pt) node[above, black] {$^{4}$He};
        \fill[popgreen] (5,4.35) circle (2.5pt) node[above, black] {$^{56}$Fe};
        \fill[poppurple] (9.6,3.75) circle (2.5pt) node[above, black] {$^{235}$U};
        \node[popmuted] at (5,0.35) {$\approx 56$};
        \node[popmuted] at (9.6,0.35) {$\approx 235$};
    \end{tikzpicture}
    \legende{Binding energy per nucleon versus nucleon number (Aston curve).}
    \end{popfigure}
    \item \textbf{Qualitative explanation [4 marks].}
    \begin{enumerate}
        \item \emph{Fission:} $^{235}\mathrm{U}$ lies on the descending part of the curve ($B/A \approx 7.6\ \mathrm{MeV}$). When it splits into two medium-mass nuclei ($A \approx 90$--$140$), the products lie \emph{higher} on the curve ($B/A \approx 8.5\ \mathrm{MeV}$), i.e.\ they are more tightly bound. The increase in total binding energy is released as kinetic energy and radiation: $E = (B/A)_{\text{products}} - (B/A)_{\text{reactants}} > 0$. [2 marks]
        \item \emph{Fusion:} very light nuclei ($^{2}\mathrm{H}$, $^{3}\mathrm{H}$) lie very low on the curve ($B/A \approx 1$--$3\ \mathrm{MeV}$). Fusing them into $^{4}\mathrm{He}$ ($B/A \approx 7.1\ \mathrm{MeV}$) moves the nucleons \emph{up} the curve to a much more bound state, releasing the difference as energy. [2 marks]
    \end{enumerate}
    In both cases, energy is released because the products are \emph{more tightly bound} (greater $B/A$) than the reactants.
    \item \textbf{Energy of the D--T fusion [4 marks].}
    \begin{align*}
    m_{\text{reactants}} &= 2.014102 + 3.016049 = 5.030151\ \mathrm{u}\\
    m_{\text{products}} &= 4.002603 + 1.008665 = 5.011268\ \mathrm{u}\\
    \Delta m &= 5.030151 - 5.011268 = \SI{0.018883}{u}\\
    E &= 0.018883 \times 931.5 \approx \SI{17.6}{MeV}.
    \end{align*}
    \item \textbf{Difficulties of controlled fusion [3 marks]} — any two:
    \begin{itemize}
        \item Reaching and maintaining temperatures of $\sim 10^{8}\ \mathrm{K}$ so that nuclei can overcome Coulomb repulsion — no material container can hold such a plasma (magnetic or inertial confinement is needed).
        \item Maintaining a sufficient plasma density and confinement time (Lawson criterion) for more energy to be produced than is injected.
        \item The large energy input currently exceeds the energy output: a net energy gain in a continuously operating power reactor is not yet industrially achieved.
    \end{itemize}
\end{enumerate}
\end{corrige}

\begin{corrige}[Exercise 11 — Comprehensive radioactivity problem (15 marks)]
\begin{enumerate}
    \item \textbf{Number of $\alpha$ and $\beta^{-}$ decays [3 marks].} Let $n_\alpha$ and $n_\beta$ be the numbers of decays. Conservation of nucleon number:
    \[238 = 206 + 4n_\alpha \implies n_\alpha = \frac{238 - 206}{4} = 8.\]
    Conservation of charge (each $\alpha$ removes 2 units of $Z$, each $\beta^{-}$ adds 1):
    \[92 = 82 + 2n_\alpha - n_\beta \implies n_\beta = 82 + 16 - 92 = 6.\]
    \textbf{8 $\alpha$-particles and 6 $\beta^{-}$-particles} are emitted per atom.
    \item \textbf{Decay constant [2 marks].} Convert the half-life to seconds:
    \[T_{1/2} = 4.468\times10^{9}\ \mathrm{y} \times 365.25 \times 24 \times 3600\ \mathrm{s} \approx 1.410\times10^{17}\ \mathrm{s}.\]
    \[\lambda = \frac{\ln 2}{T_{1/2}} = \frac{0.693}{1.410\times10^{17}} \approx \SI{4.92e-18}{s^{-1}}.\]
    \item \textbf{Mass of lead after one half-life [4 marks].} After one half-life, half of the uranium atoms have decayed: the number of decayed atoms is
    \[N_{\text{decayed}} = \frac{1}{2}\times\frac{1.00\ \mathrm{g}}{238\ \mathrm{g\,mol^{-1}}}\times 6.022\times10^{23} \approx 1.265\times10^{21}\ \text{atoms}.\]
    Each decayed $^{238}\mathrm{U}$ atom ends as one $^{206}\mathrm{Pb}$ atom, so
    \[m_{\text{Pb}} = \frac{N_{\text{decayed}}}{N_A}\times 206\ \mathrm{g\,mol^{-1}} = \frac{1.00}{2}\times\frac{206}{238} \approx \SI{0.433}{g}.\]
    (The remaining uranium mass is $0.50\ \mathrm{g}$; the difference, $0.50 - 0.433 = 0.067\ \mathrm{g}$, is carried away mainly by the 8 emitted $\alpha$-particles per atom.)
    \item \textbf{Radon hazard [3 marks].} $^{222}\mathrm{Rn}$ is an intermediate member of the $^{238}\mathrm{U}$ series, so soils and rocks containing uranium continuously produce it. Being a chemically inert \emph{gas}, it diffuses out of the ground and accumulates in poorly ventilated basements. Inhaled radon and especially its solid radioactive daughters deposit in the lungs, where their $\alpha$-emissions — highly ionising over short ranges — damage lung tissue and increase the risk of lung cancer. The higher the uranium content of the local soil/rock, the greater the radon flux into buildings.
    \item \textbf{Cloud chamber tracks [3 marks].}
    \begin{itemize}
        \item \emph{Thick, straight, short track:} an \cle{$\alpha$-particle}. Heavy and doubly charged, it ionises intensely (dense track), loses energy rapidly (short range, a few cm in air) and is barely deviated by collisions (straight).
        \item \emph{Thin, tortuous, long track:} a \cle{$\beta^{-}$-particle} (electron). Very light and singly charged, it ionises weakly (thin track), keeps a long range, and is easily scattered by atomic electrons, giving an irregular, winding path.
    \end{itemize}
\end{enumerate}
\end{corrige}

\newpage

\coursec{Summary sheet}

\begin{popfigure}
\centering
\begin{tikzpicture}[
    level distance=3.5cm,
    level 1/.style={sibling angle=360/7},
    edge from parent/.style={draw, thick},
    every node/.style={draw, rounded corners, fill=popblueL, text width=2.5cm, align=center, font=\small}
]
\node[fill=popdark, text=white] {Radioactivity}
    child {node[fill=popgreen, align=center]{Atomic Nucleus\\Stability}}
    child {node[fill=poporange, align=center]{Types of\\Radiation}}
    child {node[fill=poppurple, align=center]{Properties\\ \& Applications}}
    child {node[fill=poppink, align=center]{Mass Defect\\ \& Energy}}
    child {node[fill=popgold, align=center]{Nuclear\\Fission}}
    child {node[fill=popblueL, align=center]{Nuclear\\Fusion}}
    child {node[fill=popmuted, align=center]{Key Equations\\ \& Safety}};
\end{tikzpicture}
\legende{Mind map of the Radioactivity chapter}
\end{popfigure}

\begin{bilan}
\textbf{Key Definitions}\par
\begin{itemize}
\item \textbf{Nucleus}: Composed of protons (Z) and neutrons (N), mass number A = Z + N.
\item \textbf{Isotopes}: Same Z, different N.
\item \textbf{Radioactivity}: Spontaneous disintegration of unstable nuclei with emission of radiation.
\item \textbf{Activity (A)}: Number of decays per second, unit becquerel (Bq). $A = \lambda N$.
\item \textbf{Decay constant ($\lambda$)}: Probability per unit time of decay.
\item \textbf{Half-life ($T_{1/2}$)}: Time for activity to halve. $T_{1/2} = \frac{\ln 2}{\lambda}$.
\item \textbf{Mass defect ($\Delta m$)}: Difference between mass of nucleons and actual nuclear mass. $\Delta m = Zm_p + Nm_n - M_{\text{nuc}}$.
\item \textbf{Binding energy ($E_b$)}: Energy equivalent of mass defect, $E_b = \Delta m c^2$.
\item \textbf{Nuclear fission}: Splitting of heavy nucleus into lighter fragments with energy release.
\item \textbf{Nuclear fusion}: Combining light nuclei to form heavier nucleus with energy release.
\end{itemize}

\textbf{Laws, Formulas \& Theorems}\par
\begin{itemize}
\item Radioactive decay law: $N = N_0 e^{-\lambda t}$, $A = A_0 e^{-\lambda t}$.
\item Half-life relation: $T_{1/2} = \frac{\ln 2}{\lambda}$.
\item Mass-energy equivalence: $E = \Delta m c^2$; $1\ \mathrm{u} = 931.5\ \mathrm{MeV}$.
\item Binding energy per nucleon: $\frac{E_b}{A}$; peak at $^{56}\mathrm{Fe}$.
\item Fission energy: $Q = [M_{\text{initial}} - \sum M_{\text{final}}]c^2$.
\item Fusion energy: similar Q-value.
\item Radiation types: $\alpha$ (He nucleus), $\beta^-$ (electron), $\beta^+$ (positron), $\gamma$ (photon).
\item Decay equations: $\ce{^{A}_{Z}X -> ^{A-4}_{Z-2}Y + ^{4}_{2}He}$ ($\alpha$), $\ensuremath{\mathrm{^{A}_{Z}X \rightarrow  ^{A}_{Z+1}Y + ^{0}_{-1}e + \bar{\nu}_e}}$ ($\beta^-$), $\ensuremath{\mathrm{^{A}_{Z}X \rightarrow  ^{A}_{Z-1}Y + ^{0}_{+1}e + \nu_e}}$ ($\beta^+$), $\gamma$ emission no change in A or Z.
\end{itemize}

\textbf{Methods}\par
\begin{enumerate}
\item \textbf{Decay problems}: Identify $N_0$, $\lambda$ or $T_{1/2}$, time $t$; use $N = N_0 e^{-\lambda t}$.
\item \textbf{Half-life from data}: Plot $\ln A$ vs $t$; slope $= -\lambda$; $T_{1/2} = \ln 2 / \lambda$.
\item \textbf{Mass defect}: Compute $\Delta m = Zm_p + Nm_n - M_{\text{atom}} + Zm_e$ (using atomic masses).
\item \textbf{Binding energy}: $E_b = \Delta m \times 931.5\ \mathrm{MeV}$.
\item \textbf{Fission/Fusion Q-value}: $Q = (\sum m_{\text{initial}} - \sum m_{\text{final}})c^2$.
\item \textbf{Radiation identification}: Use penetration (paper, Al, Pb) and magnetic deflection.
\end{enumerate}

\textbf{Common Mistakes}\par
\begin{itemize}
\item Confusing activity (Bq) with count rate (cps); count rate = activity $\times$ detection efficiency.
\item Forgetting to use atomic masses (include electrons) in mass defect calculations.
\item Misapplying half-life: $N = N_0 (1/2)^{t/T_{1/2}}$ only for integer half-lives; use exponential for exact.
\item Mixing up $\beta^-$ and $\beta^+$ decay equations (change in Z).
\item Neglecting neutrinos in $\beta$ decay energy balance.
\item Using $c = 3.00 \times 10^8\ \mathrm{m/s}$ but forgetting to convert MeV to joules when needed.
\item Assuming all radiation is equally ionising; $\alpha$ is most ionising, $\gamma$ least.
\end{itemize}

\textbf{GCE Tips}\par
\begin{itemize}
\item \textbf{Definitions}: Learn precise wording (e.g., "half-life is the time taken for the activity of a sample to fall to half its initial value").
\item \textbf{Graphs}: Be able to sketch $N$ vs $t$, $\ln N$ vs $t$, and determine half-life from graph.
\item \textbf{Calculations}: Show all steps; keep units consistent (kg, MeV, J). Use $1\ \mathrm{u} = 1.66 \times 10^{-27}\ \mathrm{kg} = 931.5\ \mathrm{MeV}$.
\item \textbf{Nuclear equations}: Balance A and Z; include neutrinos/antineutrinos for $\beta$ decay.
\item \textbf{Applications}: Know uses: $\alpha$ (smoke detectors), $\beta$ (thickness gauges), $\gamma$ (sterilisation, tracers), fission (power), fusion (future energy).
\item \textbf{Safety}: Time, distance, shielding; $\alpha$ stopped by paper, $\beta$ by few mm Al, $\gamma$ by thick Pb.
\item \textbf{Cameroon context}: Mention potential for nuclear energy (e.g., uranium deposits in Cameroon) and medical isotopes.
\end{itemize}
\end{bilan}

\findecours

\end{document}
